\section{农产品加工：规模最大、利润最薄}
\label{sec:processing}

农产品加工板块 \num{23} 家公司（其他农产品加工 \num{11}、粮油加工 \num{7}、果蔬加工 \num{5}），
\hl{是公司数量最多的板块}，流通市值合计约 \num{1464} 亿元。
按 2026 年半年报口径，\num{8} 家亏损，半年 ROE 中位数 \num{1.40}\%（未年化），
2026 年 6 月末资产负债率超过 \num{65}\% 的仅 \num{2} 家。

\begingroup\scriptsize\setlength{\tabcolsep}{2pt}
\begin{longtable}{@{}llrrrrrrr@{}}
\caption{农业上市公司分行业明细（粮油加工、果蔬加工、其他农产品加工；财务为 2026 年半年报，ROE 未年化，公告计数窗口 2023-09---2026-09）}\label{tab:comp-processing}\\
\toprule
代码 & 公司简称 & 市值（亿元） & 营收（亿元） & 同比（\%） & 净利（亿元） & ROE（\%） & 再融资公告 & 重组公告 \\
\midrule
\endfirsthead
\multicolumn{9}{l}{\small（续）}\\
\toprule
代码 & 公司简称 & 市值（亿元） & 营收（亿元） & 同比（\%） & 净利（亿元） & ROE（\%） & 再融资公告 & 重组公告 \\
\midrule
\endhead
\bottomrule
\endlastfoot
600737 & 中粮糖业 & 315.5 & 92.8 & -21.2 & 6.08 & 5.2 & 1 & 0 \\
605198 & 安德利 & 172.9 & 7.2 & -24.4 & 1.23 & 4.3 & 17 & 2 \\
300999 & 金龙鱼 & 133.2 & 1,228.5 & 6.2 & 22.94 & 2.4 & 0 & 0 \\
600127 & 金健米业 & 105.8 & 19.7 & 25.2 & -0.10 & -1.4 & 0 & 4 \\
000930 & 中粮科技 & 97.6 & 81.9 & -7.0 & -1.69 & -1.6 & 0 & 0 \\
600251 & 冠农股份 & 77.9 & 13.9 & 0.3 & 5.37 & 13.3 & 0 & 1 \\
600962 & 国投中鲁 & 68.6 & 8.8 & -12.4 & 0.02 & 0.2 & 85 & 32 \\
002481 & 双塔食品 & 46.5 & 11.2 & 7.3 & 0.33 & 1.4 & 1 & 2 \\
001338 & 永顺泰 & 45.7 & 19.1 & -13.9 & 0.51 & 1.4 & 1 & 0 \\
300138 & 晨光生物 & 44.6 & 31.6 & -13.6 & 1.79 & 5.2 & 0 & 0 \\
600191 & 华资实业 & 43.9 & 3.5 & 58.0 & 0.36 & 2.1 & 7 & 2 \\
000505 & 京粮控股 & 42.1 & 31.1 & -26.0 & 0.08 & 0.3 & 0 & 2 \\
002086 & 东方海洋 & 32.3 & 1.6 & 0.9 & -0.48 & -3.8 & 0 & 0 \\
002286 & 保龄宝 & 31.3 & 15.0 & 6.9 & 0.82 & 3.6 & 9 & 0 \\
000972 & 中基健康 & 30.5 & 2.0 & -19.6 & -0.66 & — & 36 & 16 \\
000019 & 深粮控股 & 28.2 & 23.9 & 0.1 & 1.12 & 2.2 & 0 & 0 \\
300175 & 朗源股份 & 25.2 & 1.7 & 67.7 & -0.13 & -2.7 & 0 & 2 \\
002852 & 道道全 & 23.1 & 28.6 & 2.6 & 0.17 & 0.8 & 3 & 0 \\
603182 & 嘉华股份 & 22.2 & 8.6 & 25.7 & 0.57 & 5.0 & 9 & 6 \\
003030 & 祖名股份 & 21.9 & 10.1 & 10.0 & -0.24 & -2.5 & 0 & 3 \\
000911 & *ST广糖 & 20.5 & 15.8 & 17.4 & -0.16 & 0.0 & 78 & 1 \\
603231 & 索宝蛋白 & 18.2 & 9.7 & 24.9 & 0.68 & 3.4 & 0 & 1 \\
300268 & 佳沃食品 & 16.0 & 4.2 & -66.1 & -0.04 & -1.4 & 0 & 20 \\
\end{longtable}
\endgroup


\subsection{规模与利润的严重错配}
\label{sec:processing-scale}

加工板块最直观的对照来自两家公司：\textbf{金龙鱼} 2026 年上半年营收 \num{1228.5} 亿元
（占 \num{104} 家农业上市公司上半年营收总额 \num{6085} 亿元的 \num{20.2}\%、
占本板块营收的 \num{73.5}\%）实现 \num{22.94} 亿元净利润，
\hl{净利率仅 \num{1.94}\%}；而 \textbf{安德利}以 \num{7.16} 亿元营收
实现 \num{1.23} 亿元净利润、\hl{净利率 \num{17.2}\%}，
其流通市值（\num{172.9} 亿元）甚至超过金龙鱼（\num{133.2} 亿元）。
\emph{这组对照揭示了加工行业的周期本质}：
以大宗农产品为原料、产品同质化的“量大利薄”模式（粮油加工）
没有定价权，利润被原料价格波动与渠道费用吞噬；
而以浓缩苹果汁、番茄制品等\hl{细分品类 + 出口市场}为核心的加工企业
（安德利半年 ROE \num{4.3}\%、冠农股份 \num{13.3}\%、晨光生物 \num{5.2}\%）
可以获得接近品牌消费品的利润率。

\subsection{食糖：当周期回到加工环节}
\label{sec:processing-sugar}

\textbf{中粮糖业}（流通市值 \num{315.5} 亿元，板块第一）2026 年上半年实现营收
\num{92.8} 亿元、净利润 \num{6.08} 亿元（半年 ROE \num{5.19}\%），
\hl{营收同比 \num{-21.2}\%，利润却同比 \num{+36.7}\%}（上年同期 \num{4.45} 亿元）。
把这家公司放回第 \ref{sec:soft} 章的框架中看，其业绩完全由白糖周期解释：
2024 年以来国内食糖产量恢复（广西、云南连续两年丰产）、
配额外进口成本下降，2026 年 8 月末广西白糖现货价格下滑至 \num{5420} 元/吨，
\hl{郑商所白糖主力合约 2021 年 2 月至 2026 年 9 月仅累计上涨 \num{0.8}\%}
——基本回到基期水平，明显弱于其他软商品（棉花 \num{3.5}\%、橡胶 \num{24.8}\%）。
对加工企业而言，原料价格下跌在库存周转周期内先体现为\hl{存货减值}，
因此糖价下跌的第一年（2025 年）加工企业的利润往往承压，
而第二年才因成本下降而受益。\emph{这是加工环节与种植环节的典型时滞差异}，
中粮糖业 2026 年上半年“收入降、利润增”的组合正是这一逻辑的兑现。

\subsection{困境公司的集中区}
\label{sec:processing-distress}

加工板块是农业上市公司中\hl{亏损集中度最高}的区域之一（2026 年上半年 \num{8} 家亏损）：
\textbf{中粮科技}（\num{-1.69} 亿元）、\textbf{中基健康}（\num{-0.66} 亿元）、
\textbf{东方海洋}（\num{-0.48} 亿元）、\textbf{祖名股份}（\num{-0.24} 亿元）
四家合计亏损 \num{3.1} 亿元，占板块全部亏损额（\num{3.5} 亿元）的 \num{88.6}\%。
从公告行为看，这批公司与高公告密度群体高度重合：
\textbf{*ST广糖}近三年再融资类公告 \num{78} 条（在全部 \num{104} 家公司中排第三），
2026 年 6 月末资产负债率 \num{100.2}\%；
\textbf{国投中鲁}（果蔬加工，2026 年上半年微利 \num{0.02} 亿元）
以 \num{85} 条再融资类公告居全样本第二，并伴随 \num{32} 条重组并购类公告（全样本第一）；
中基健康的资产负债率更高达 \num{104.1}\%——\hl{净资产已经为负}，
其半年 ROE 因此无法计算。
\hl{高公告密度 + 持续亏损 + ST 标识，构成一组清晰的“重组与保壳”信号}，
这将在第 \ref{sec:financing} 节中作为融资需求分类的依据。

\begin{keybox}[农产品加工的画像结论]
\normalsize
\textbf{① 规模不产生利润}：板块 \num{23} 家公司的上半年营业收入合计 \num{1671} 亿元、
占全样本的 \num{27.5}\%，而净利润合计只有 \num{38.6} 亿元，
其中金龙鱼一家占 \num{59.5}\%。\par\vspace{3pt}
\textbf{② 细分品类 + 出口是唯一的高毛利路径}：冠农股份（半年 ROE \num{13.3}\%）、
晨光生物（\num{5.2}\%）、安德利（\num{4.3}\%）均明显高于板块中位数 \num{1.40}\%，
而粮油加工环节的 ROE 中位数只有 \num{2.22}\%、其他农产品加工仅 \num{1.36}\%。\par\vspace{3pt}
\textbf{③ 食糖加工是“周期放大器”}：\hl{糖价下跌第一年存货减值、
第二年成本改善}，中粮糖业 2026 年上半年营收同比 \num{-21.2}\%、
利润同比 \num{+36.7}\% 的背离正源于此。\par\vspace{3pt}
\textbf{④ 利润逆势回升。}板块 \num{23} 家公司 2026 年上半年合计归母净利润
\num{38.58} 亿元（上年同期 \num{33.46} 亿元），\num{8} 家亏损；
\hl{低毛利、强刚需的加工环节在猪价下行期的相对优势再次显现}。
\end{keybox}

\subsection{上市公司画像（\num{23} 家）}
\label{sec:co-processing}

本节给出农产品加工行业全部 \num{23} 家上市公司的画像，按流通市值降序排列。
数据口径与第 \ref{sec:comp-overview} 节一致：\hl{财务数据统一采用 2026 年半年报}
（合并报表口径、未年化，资产负债率为 2026 年 6 月末），
同时列出 2025 年半年报作为同口径对照、2023---2025 年年报作为趋势对照；
股价为月度收盘价，最后一个交易日为 2026 年 9 月 24 日；周期分段使用与第一部分相同的
ZigZag 算法，阈值取 25\%（个股波动大于商品价格，故高于生猪现货的 20\%）；
公告计数窗口为 2023 年 9 月 1 日至 2026 年 9 月 27 日。

与前一版不同，本版画像正文由\hl{逐家检索的公开资料底稿}驱动：底稿只收录有来源链接的事实
（公司沿革、股权与控制权变动、融资历史、风险事件、战略动作、融资需求线索），
正文按事实生成并以脚注给出来源链接，因此各家公司的段落结构与重点并不相同。
每家公司在统一顺序下展开：基本情况与主营结构、历史股价、过去周期分析、盈利情况、财务分析、
重大战略转型、历史融投资情况、未来潜在融投资需求。每家公司配图两张：
股价与行业指数（申万农林牧渔）的对比图，以及近年主要财务指标的组合图
（营业收入与归母净利润、ROE 与资产负债率、货币资金与有息负债），
并给出关键财务指标表与历史融资/资本运作表。
未来融资需求按两个口径粗略度量：\hl{短期债务缺口 $=$ 期末短期债务 $-$ 货币资金}
与\hl{年化经营缺口 $=$ 半年归母亏损 $\times 2$}；
前者衡量滚动偿付压力，后者衡量维持当期盈利水平的现金消耗，
两者都是压力情形下的静态估算，不构成对融资规模或时点的预测。

% 农业上市公司画像：粮油加工、果蔬加工、其他农产品加工（共 23 家，按流通市值降序）
% 由 scripts/make_company_profiles.py 自动生成，请勿手工修改

\subsubsection{中粮糖业（600737）}
\label{co:600737}

\pfl{公司基本情况} 中粮糖业成立于 1993 年，中粮集团旗下食糖与番茄双主业平台，国内食糖行业龙头，拥有从国内外制糖、进口炼糖到贸易、仓储物流的全产业链；番茄制品规模居亚洲第一、世界第二。经历 2004 年德隆系危机后由中粮重组脱困，近年受糖价与番茄酱价格下行影响业绩承压。 公司于 1996-07-31 在上交所首发上市，注册资本 21.39 亿元，总股本 21.39 亿股。 截至 2026-09-24收盘价 14.75 元，流通市值 315.5 亿元，近一年-11.5\%，流通市值在三级行业“其他农产品加工”的 11 家公司中排第 \zhnumber{1} 位，近一年涨跌幅低于全样本中位数（-10.1\%）1.4 个百分点。 按 2025 年年报披露的行业构成，收入占比前三位为工业 66.3\%；贸易 32.9\%；其他(补充) 0.7\%。 与 2021 年年报相比，第一大行业由“贸易”变为“工业”，主业重心已经迁移。 发展过程中的关键节点包括：1996 年 7 月 31 日在上海证券交易所挂牌上市，证券简称“新疆屯河”，证券代码 ……（1996）；德隆集团入主新疆屯河，推动公司由水泥建材向番茄酱加工（“红色产业”）转型，先后投入 2 ……（1999）；受德隆系危机冲击，披露 2004 年净利润 -6.66 亿元；截至 2005 年 4 月……（2004）。\footnote{\texttt{http://www.sse.com.cn/disclosure/listedinfo/announcement/c/new/2021-04-20/600737\_2021042}}\footnote{\texttt{https://pdf.dfcfw.com/pdf/H3\_AP201906281336801577\_1.pdf}}\footnote{\texttt{http://static.cninfo.com.cn/finalpage/2005-04-15/15489757.PDF}}

\pfl{历史股价} 本报告样本区间（2021 年 1 月 至 2026 年 9 月）内，股价由 9.19 元变为 14.75 元，累计 +60.5\%，区间最高 17.74 元（2026 年 1 月）、最低 6.62 元（2022 年 10 月）；当前价相当于区间最高价的 83.1\%，为区间最低价的 2.23 倍。 月度收益的年化波动率 33.5\%，低于全样本中位数 37.4\%，在三级行业“其他农产品加工”的 11 家公司中波动率排第 \zhnumber{8} 位。

\begin{center}
\begin{tikzpicture}
\begin{axis}[
  width=12.6cm, height=3.9cm, scale only axis,
  xmin=2020.922, xmax=2026.888, ymin=6.16, ymax=18.98,
  xtick={2021,2022,2023,2024,2025,2026},
  yearticks,
  yticklabel style={font=\scriptsize},
  ylabel={元/股}, ylabel style={font=\scriptsize},
  grid=major, grid style={LightGray, line width=0.3pt},
  axis line style={InkGray, line width=0.5pt},
  every axis plot/.append style={line width=0.9pt},
]
\addplot[AgriGreen] table[x=x,y=y]{data/clean/profiles/px_600737.dat};
\addplot[SoftRed, dashed, line width=0.7pt] table[x=x,y=y]{data/clean/profiles/pxzz_600737.dat};
\addplot[only marks, mark=*, mark size=1.1pt, SoftRed] table[x=x,y=y]{data/clean/profiles/pxzz_600737.dat};
\end{axis}
\end{tikzpicture}

\captionof{figure}[中粮糖业（600737）股价与周期骨架]{中粮糖业（600737）月度收盘价（元/股）与 ZigZag 周期骨架（阈值 25\%，圆点为周期转折点）}
\end{center}

\pfl{过去周期分析} 以 25\% 的 ZigZag 阈值划分，样本区间内股价形成 2 个主升段与 2 个主跌段。 最大主跌段为 2021 年 6 月 至 2022 年 10 月，16 个月-33.9\%；最大主升段出现在 2022 年 10 月 至 2026 年 1 月，39 个月+168.0\%。 最近一次转折出现在 2026 年 6 月（下跌段自 2026 年 1 月起，5 个月-30.5\%），当前处于该段的延续之中。 公司股价较申万农林牧渔指数提前约 43 个月见底，是板块中较早定价周期反转的公司之一。 该子行业横跨糖、番茄制品、添加剂等多个品种，各自的供给周期与政策（进口配额、关税）是主要变量，公司 2026 年上半年实现盈利、半年 ROE 5.2\%（未年化），好于全样本中位数（0.75\%），下行周期对其报表的冲击小于同业。

\pfl{盈利情况} 2026 年上半年营业收入 92.77 亿元（同比 -21.2\%），归母净利润 6.08 亿元，同比 +36.7\%。 以年报为趋势对照，2023---2025 年营业收入由 331.14 亿元变为 265.11 亿元（两年年均复合增速 -10.5\%），归母净利润由 20.73 亿元变为 9.51 亿元。 按半年报口径，2026 年上半年的 ROE 为 5.2\%（未年化）、销售净利率 6.7\%、毛利率 14.6\%，ROE 在“其他农产品加工”11 家公司中排第 \zhnumber{1} 位。 同期全样本半年 ROE 中位数 0.75\%，公司与之相差 4.44 个百分点（略好于中位数公司）。 结合公司披露，年报披露商誉账面原值合计 358,586,291.61 元，累计计提减值准备 358,568,664.9……；2025 年上半年受国内外糖价震荡走弱、大包装番茄酱价格下行影响，营业收入 117.67 亿元，同比减少……。\footnote{\texttt{https://static.cninfo.com.cn/finalpage/2026-04-28/1225220574.PDF}}\footnote{\texttt{http://dataclouds.cninfo.com.cn/shgonggao/2025/2025-08-28/2359255d832e11f0a692fa163e957f}}

\begin{center}\small\setlength{\tabcolsep}{4pt}
\captionof{table}[中粮糖业（600737）关键财务指标]{中粮糖业（600737）关键财务指标（合并报表口径；两个半年报期均未年化；现金短债比 $=$ 货币资金 $\div$ 短期债务）}
\begin{tabular}{@{}lrrrrr@{}}
\toprule
指标 & 2023 年 & 2024 年 & 2025 年 & 2025 年半年报 & 2026 年半年报 \\
\midrule
营业收入（亿元） & 331.14 & 324.97 & 265.11 & 117.67 & 92.77 \\
归母净利润（亿元） & 20.73 & 17.13 & 9.51 & 4.45 & 6.08 \\
毛利率（\%） & 14.1 & 12.5 & 11.4 & 8.5 & 14.6 \\
ROE（\%） & 19.4 & 14.4 & 8.2 & 3.9 & 5.2 \\
资产负债率（\%） & 44.9 & 41.5 & 44.1 & 48.0 & 39.1 \\
经营活动现金流净额（亿元） & 8.87 & 26.34 & 12.88 & 10.28 & 11.98 \\
货币资金（亿元） & 10.06 & 8.37 & 14.65 & 21.72 & 8.58 \\
有息负债合计（亿元） & 37.58 & 35.66 & 46.37 & 42.49 & 30.71 \\
现金短债比（倍） & 0.27 & 0.24 & 0.32 & 0.52 & 0.28 \\
\bottomrule
\end{tabular}
\end{center}

\begin{center}
\begin{tikzpicture}
\begin{groupplot}[
  group style={group size=2 by 1, horizontal sep=0.95cm},
  width=7.2cm, height=4.6cm,
  tick label style={font=\tiny},
  title style={font=\scriptsize\heiti},
  yticklabel style={font=\tiny},
  grid=major, grid style={LightGray, line width=0.3pt},
  axis line style={InkGray, line width=0.5pt},
  enlarge x limits={abs=0.8},
  legend style={font=\scriptsize, at={(0.5,-0.22)}, anchor=north,
                legend columns=2, draw=none, fill=none, column sep=6pt},
]
\nextgroupplot[
  ybar, bar width=4pt,
  ymin=-33.11, ymax=370.88,
  xtick={0,1,2,3,4,5.7},
  xticklabels={2021,2022,2023,2024,2025,2026H1},
  ylabel={亿元},
]
\addplot[fill=AgriGreen!75, draw=AgriGreen, bar shift=-2.6pt] table[x=i, y=rev]{data/clean/profiles/fin_600737.dat};
\addplot[fill=SoftRed!60, draw=SoftRed, bar shift=2.6pt] table[x=i, y=ni]{data/clean/profiles/fin_600737.dat};
\addplot[fill=AgriGreen!30, draw=AgriGreen, pattern=north east lines, bar shift=-2.6pt] table[x=x, y=rev]{data/clean/profiles/finh_600737.dat};
\addplot[fill=SoftRed!20, draw=SoftRed, pattern=north east lines, bar shift=2.6pt] table[x=x, y=ni]{data/clean/profiles/finh_600737.dat};
\legend{营业收入（年/半年）, 归母净利润（年/半年）}
\nextgroupplot[
  ymin=-3.7, ymax=50.4,
  xtick={0,1,2,3,4,5.7},
  xticklabels={2021,2022,2023,2024,2025,2026H1},
  ylabel={\%},
]
\addplot[AgriGold, mark=*, mark size=1.3pt] table[x=i, y=roe]{data/clean/profiles/fin_600737.dat};
\addplot[OilBlue, mark=square*, mark size=1.3pt] table[x=i, y=debt]{data/clean/profiles/fin_600737.dat};
\addplot[only marks, AgriGold, mark=*, mark size=2.0pt] table[x=x, y=roe]{data/clean/profiles/finh_600737.dat};
\addplot[only marks, OilBlue, mark=square*, mark size=2.0pt] table[x=x, y=debt]{data/clean/profiles/finh_600737.dat};
\legend{ROE, 资产负债率}
\end{groupplot}
\end{tikzpicture}
\begin{tikzpicture}
\begin{axis}[
  width=15.2cm, height=3.1cm, scale only axis,
  ymin=-5.37, ymax=52.01,
  xtick={0,1,2,3,4,5.7},
  xticklabels={2021,2022,2023,2024,2025,2026H1},
  tick label style={font=\tiny},
  yticklabel style={font=\tiny},
  ylabel={亿元}, ylabel style={font=\scriptsize},
  ybar, bar width=4pt,
  grid=major, grid style={LightGray, line width=0.3pt},
  axis line style={InkGray, line width=0.5pt},
  enlarge x limits={abs=0.8},
  legend style={font=\scriptsize, at={(0.5,-0.30)}, anchor=north,
                legend columns=3, draw=none, fill=none, column sep=10pt},
]
\addplot[fill=AgriGreen!55, draw=AgriGreen, bar shift=-3.0pt] table[x=i, y=cash]{data/clean/profiles/fin_600737.dat};
\addplot[fill=SoftRed!45, draw=SoftRed, bar shift=3.0pt] table[x=i, y=ibd]{data/clean/profiles/fin_600737.dat};
\addplot[fill=AgriGreen!25, draw=AgriGreen, pattern=north east lines, bar shift=-3.0pt] table[x=x, y=cash]{data/clean/profiles/finh_600737.dat};
\addplot[fill=SoftRed!20, draw=SoftRed, pattern=north east lines, bar shift=3.0pt] table[x=x, y=ibd]{data/clean/profiles/finh_600737.dat};
\addplot[OilBlue, mark=*, mark size=1.1pt, line width=0.8pt] table[x=i, y=ocf]{data/clean/profiles/fin_600737.dat};
\addplot[only marks, OilBlue, mark=*, mark size=1.8pt] table[x=x, y=ocf]{data/clean/profiles/finh_600737.dat};
\legend{货币资金, 有息负债, 经营活动现金流净额}
\end{axis}
\end{tikzpicture}

\captionof{figure}[中粮糖业（600737）近年主要财务指标]{中粮糖业（600737）近年主要财务指标（左上：营业收入与归母净利润；右上：ROE 与资产负债率；下：货币资金、有息负债与经营活动现金流净额。
2021---2025 年为年报口径，2026H1 为 2026 年半年报/半年末，与其前一年报之间留出间隔、以斜纹或空心点区分；半年报数据未年化）}
\end{center}

\pfl{财务分析} 2026 年 6 月末资产负债率 39.1\%，较 2025 年末 下降 5.1 个百分点（样本中位数 52.2\%，所处三级行业内排第 \zhnumber{6} 位）。 2026 年 6 月末货币资金 8.58 亿元、短期债务（短期借款与一年内到期非流动负债）30.28 亿元、长期债务 0.43 亿元，有息负债合计 30.71 亿元，现金短债比 0.28 倍——缺口明显。 净有息负债 22.13 亿元，相当于其 2026 年上半年营业收入的 24\%。 流动比率 1.82、速动比率 0.72，短期指标尚可。 2026 年上半年经营活动现金流净额 11.98 亿元，经营现金流与利润同向为正，内生资金可以支撑运营。 2026 年 6 月末在建工程 1.38 亿元，较上年末增加 0.31 亿元，仍有在建投入。 公司披露的财务与合规风险包括：德隆系危机爆发，2004 年净利润 -666,007,835.43 元；截至 2005 ……；2025 年度报告披露公司无重大诉讼、仲裁事项；董事、高管、控股股东、实际控制人不存在涉……。\footnote{\texttt{http://static.cninfo.com.cn/finalpage/2005-04-15/15489757.PDF}}\footnote{\texttt{https://static.cninfo.com.cn/finalpage/2026-04-28/1225220574.PDF}}

\pfl{重大战略转型} 公司在报告期内的战略动作可归为 3 项：1999 德隆入主后由水泥建材（“灰色产业”）向番茄酱加工（“红色产业”）战略转型，建成全国最大番茄酱生产企业；2005—2007 中粮集团收购重组后积极布局糖产业，逐步形成以食糖业务为主、番茄酱加工为辅的业务体系；2025 立足 2B 大包装番茄酱出口优势，拓展国内小包装番茄酱、番茄饮品、调味品、保健品等高附加值业务；食糖端稳定自产糖……。 公告口径上，近三年（2023 年 9 月---2026 年 9 月）公司披露重大重组与并购类公告 0 条、对外投资与转型类公告 0 条、回购与股权激励类公告 0 条，合计公告 282 条。 主营结构的迁移已经落到报表上：2021 年年报的第一大行业为“贸易”（65.9\%），2025 年年报变为“工业”（66.3\%）。 综合判断，公司在报告期内有明确的外延扩张或业务结构调整动作，转型方向与融资安排之间存在直接关联。\footnote{\texttt{https://pdf.dfcfw.com/pdf/H3\_AP201906281336801577\_1.pdf}}\footnote{\texttt{http://dataclouds.cninfo.com.cn/shgonggao/2025/2025-08-28/2359255d832e11f0a692fa163e957f}}

\begin{center}\small\setlength{\tabcolsep}{4pt}
\captionof{table}[中粮糖业（600737）历史融资与资本运作]{中粮糖业（600737）历史融资与资本运作（据公司公告与公开报道整理；金额为披露值，含首次公开发行、再融资、债券、银行借款与重整投资等）}
\begin{tabular}{@{}p{1.5cm}p{3.2cm}p{2.4cm}p{6.3cm}@{}}
\toprule
时间 & 方式 & 规模 & 用途与说明 \\
\midrule
1996-07-31 & IPO（上海证券交易所，前身新疆屯河） & 实际募集资金合计 10,115 万元 & 招股公告日 1996-06-29，发行价 5.78 元/股，主承销商君安证券有限公司 \\
2007—2008 & 非公开发行（募集资金全部用于偿还对中粮集团的借款） & --- & 偿还对中粮集团的借款；经公司 2007 年第二次临时股东大会审议…… \\
2019-04 & 非公开发行 A 股（上交所） & 发行 86,972,073 股，发行价 7.52 元/股，募集资金总额 654,029,988.96 元（净额约 6.41 亿元） & 甘蔗糖技术升级改造及配套优质高产高糖糖料蔗基地建设项目、甜菜糖项目…… \\
2022-12 & 子公司收购关联方资产 & 6,843.87 万元 & 全资子公司中粮屯河番茄有限公司收购中粮集团持有的内蒙古中粮番茄制品有限公司 100\% 股…… \\
\bottomrule
\end{tabular}
\end{center}

\pfl{历史融投资情况} 从现金流量表看，2025 年公司取得借款收到的现金 104.59 亿元、偿还债务支付的现金 95.16 亿元、吸收权益性投资 — 亿元，筹资活动现金流净额 -0.55 亿元（2023 年为取得借款 123.93 亿元、偿还债务 125.08 亿元）。 2026 年上半年筹资活动现金流净额为 -15.84 亿元（取得借款 24.64 亿元、偿还债务 39.79 亿元，未年化），已转为净流出，处于净还债状态。 2025 年分配股利、利润或偿付利息支付的现金合计 9.59 亿元。 公司披露的股本与债务融资脉络包括：IPO（上海证券交易所，前身新疆屯河）、子公司收购关联方资产、非公开发行 A 股（上交所）、非公开发行（募集资金全部用于偿还对中粮集团的借款）。 其中规模较大的两笔为：1996-07-31IPO（上海证券交易所，前身新疆屯河）（实际募集资金合计 10,115 万元）；2022-12子公司收购关联方资产（6,843.87 万元）。 股本方面，近三年披露股本变动 3 次；总股本由 2021 年末的 21.39 亿股变为 21.39 亿股（+0.00\%）。 分红方面，近三年现金分红 6 次、累计每股派现 2.08 元（最近一次 2025-12-31）——在利润承压的阶段仍维持了分红。 近三年“再融资”类公告 1 条，涉及定向增发、募集资金使用。\footnote{\texttt{https://vip.stock.finance.sina.com.cn/corp/go.php/vISSUE\_NewStock/stockid/600737.phtml}}\footnote{\texttt{http://static.cninfo.com.cn/finalpage/2008-04-22/39017468.PDF}}\footnote{\texttt{https://vip.stock.finance.sina.com.cn/corp/view/vCB\_AllBulletinDetail.php}}

\pfl{未来潜在融资需求分析} 2026 年 6 月末货币资金 8.58 亿元，而短期债务 30.28 亿元（现金短债比 0.28 倍），2026 年上半年经营现金流净额 11.98 亿元。按“短期债务减去账面货币资金”的滚动偿付口径，静态缺口约 21.70 亿元；若再计入上半年亏损的年化额（0.00 亿元），维持一年正常经营所需的外部资金约 21.70---21.70 亿元。 按第 \ref{sec:financing-need} 节的三档划分，公司属于第三档“扩张型”：2026 年上半年 ROE 5.2\%（未年化，折合约 10.4\%）、6 月末资产负债率 39.1\%，资产负债表仍具备加杠杆空间；融资用途以产能扩张、品类并购与海外布局为主，单笔规模通常在 5---20 亿元量级，会被市场定价为成长而非补血。 公开披露的下一步融资线索包括：2019—2023 2019 年非公开发行募集资金截至 2023-12-31 剩余 1,078.73 万元，募集资金使用率 80.92\%。 资金需求还要叠加在建工程：期末在建工程 1.38 亿元、2026 年上半年购建固定资产等支出 2.04 亿元（未年化），这部分支出在周期底部通常会被主动放缓。\footnote{\texttt{https://basic.10jqka.com.cn/600737/capital.html}}


\subsubsection{安德利（605198）}
\label{co:605198}

\pfl{公司基本情况} 安德利成立于 1996 年，烟台北方安德利果汁股份有限公司，山东安德利集团旗下企业，全球主要浓缩果汁加工企业之一，主营浓缩苹果汁、浓缩梨汁等多品种浓缩果汁；2003 年 H 股在港上市、2020 年回归 A 股，为国内首家饮料类“A+H”双上市企业，2026 年拟跨界收购电子材料资产。 公司于 2020-09-18 在上交所首发上市，注册资本 3.28 亿元，总股本 3.49 亿股。 截至 2026-09-24收盘价 63.91 元，流通市值 172.9 亿元，近一年+43.8\%，流通市值在三级行业“果蔬加工”的 5 家公司中排第 \zhnumber{1} 位，近一年涨跌幅高于全样本中位数（-10.1\%）53.9 个百分点。 按 2025 年年报披露的行业构成，收入占比前三位为饮料制造业 97.2\%；果渣制造业 2.5\%；其他(补充) 0.3\%。 发展过程中的关键节点包括：2003 年 4 月 22 日公司股票在香港联交所创业板上市（股份代号 8259），为国……（2003）；2011 年 1 月 19 日 H 股由香港联交所创业板转往主板买卖（股份代号 0221……（2011）；2020 年 9 月 18 日 A 股在上海证券交易所上市（证券简称后由“德利股份”变更……（2020）。\footnote{\texttt{https://www.hkexnews.hk/listedco/listconews/sehk/2026/0327/2026032703935\_c.pdf}}\footnote{\texttt{https://www.sse.com.cn/disclosure/listedinfo/announcement/c/new/2024-06-15/605198\_202406}}

\pfl{历史股价} 本报告样本区间（2021 年 1 月 至 2026 年 9 月）内，股价由 21.10 元变为 63.91 元，累计 +202.9\%，区间最高 67.75 元（2026 年 8 月）、最低 15.56 元（2022 年 4 月）；当前价相当于区间最高价的 94.3\%，为区间最低价的 4.11 倍。 月度收益的年化波动率 53.6\%，高于全样本中位数 37.4\%，在三级行业“果蔬加工”的 5 家公司中波动率排第 \zhnumber{1} 位。 从事件看，2024 年 8 月 至 2025 年 5 月的上涨（+172.1\%）与公司终止不超 3 亿元简易程序定增，原因为“目前本次发行上市条件尚不成熟”等市场……（2025-05）在时间上重合——股价的拐点多数能在公司自身的资本运作与风险事件上找到对应。\footnote{\texttt{http://finance.ce.cn/stock/gsgdbd/202505/t20250521\_2249324.shtml}}

\begin{center}
\begin{tikzpicture}
\begin{axis}[
  width=12.6cm, height=3.9cm, scale only axis,
  xmin=2020.922, xmax=2026.888, ymin=14.47, ymax=72.49,
  xtick={2021,2022,2023,2024,2025,2026},
  yearticks,
  yticklabel style={font=\scriptsize},
  ylabel={元/股}, ylabel style={font=\scriptsize},
  grid=major, grid style={LightGray, line width=0.3pt},
  axis line style={InkGray, line width=0.5pt},
  every axis plot/.append style={line width=0.9pt},
]
\addplot[AgriGreen] table[x=x,y=y]{data/clean/profiles/px_605198.dat};
\addplot[SoftRed, dashed, line width=0.7pt] table[x=x,y=y]{data/clean/profiles/pxzz_605198.dat};
\addplot[only marks, mark=*, mark size=1.1pt, SoftRed] table[x=x,y=y]{data/clean/profiles/pxzz_605198.dat};
\end{axis}
\end{tikzpicture}

\captionof{figure}[安德利（605198）股价与周期骨架]{安德利（605198）月度收盘价（元/股）与 ZigZag 周期骨架（阈值 25\%，圆点为周期转折点）}
\end{center}

\pfl{过去周期分析} 以 25\% 的 ZigZag 阈值划分，样本区间内股价形成 5 个主升段与 4 个主跌段。 最大主升段出现在 2024 年 8 月 至 2025 年 5 月，9 个月+172.1\%；最大主跌段为 2024 年 4 月 至 2024 年 8 月，4 个月-44.8\%。 最近一次转折出现在 2026 年 9 月（上涨段自 2025 年 12 月起，9 个月+69.9\%），当前处于该段之后的震荡之中。 公司股价较申万农林牧渔指数提前约 49 个月见底，是板块中较早定价周期反转的公司之一。 果蔬加工的外向型特征明显，出口订单与原料产地价格决定利润率，公司 2026 年上半年实现盈利、半年 ROE 4.3\%（未年化），好于全样本中位数（0.75\%），下行周期对其报表的冲击小于同业。

\pfl{盈利情况} 2026 年上半年营业收入 7.16 亿元（同比 -24.4\%），归母净利润 1.23 亿元，同比 -38.7\%。 以年报为趋势对照，2023---2025 年营业收入由 8.76 亿元变为 16.77 亿元（两年年均复合增速 +38.4\%），归母净利润由 2.56 亿元变为 3.30 亿元。 按半年报口径，2026 年上半年的 ROE 为 4.3\%（未年化）、销售净利率 17.2\%、毛利率 24.3\%，ROE 在“果蔬加工”5 家公司中排第 \zhnumber{2} 位。 同期全样本半年 ROE 中位数 0.75\%，公司与之相差 3.58 个百分点（略好于中位数公司）。 结合公司披露，收购公告提示风险：标的公司 2024 年、2025 年连续亏损，2026 年一季度已扭亏为盈但全年盈利仍……。\footnote{\texttt{https://paper.cnstock.com/html/2026-09/19/content\_2270544.htm}}

\begin{center}\small\setlength{\tabcolsep}{4pt}
\captionof{table}[安德利（605198）关键财务指标]{安德利（605198）关键财务指标（合并报表口径；两个半年报期均未年化；现金短债比 $=$ 货币资金 $\div$ 短期债务）}
\begin{tabular}{@{}lrrrrr@{}}
\toprule
指标 & 2023 年 & 2024 年 & 2025 年 & 2025 年半年报 & 2026 年半年报 \\
\midrule
营业收入（亿元） & 8.76 & 14.18 & 16.77 & 9.48 & 7.16 \\
归母净利润（亿元） & 2.56 & 2.61 & 3.30 & 2.01 & 1.23 \\
毛利率（\%） & 33.0 & 24.8 & 24.3 & 23.9 & 24.3 \\
ROE（\%） & 10.3 & 9.9 & 11.7 & 7.4 & 4.3 \\
资产负债率（\%） & 3.1 & 4.8 & 5.3 & 7.1 & 7.8 \\
经营活动现金流净额（亿元） & -0.30 & -1.09 & 4.92 & 7.92 & 5.38 \\
货币资金（亿元） & 5.44 & 2.36 & 4.09 & 9.59 & 7.39 \\
有息负债合计（亿元） & --- & --- & --- & --- & --- \\
现金短债比（倍） & --- & --- & --- & --- & --- \\
\bottomrule
\end{tabular}
\end{center}

\begin{center}
\begin{tikzpicture}
\begin{groupplot}[
  group style={group size=2 by 1, horizontal sep=0.95cm},
  width=7.2cm, height=4.6cm,
  tick label style={font=\tiny},
  title style={font=\scriptsize\heiti},
  yticklabel style={font=\tiny},
  grid=major, grid style={LightGray, line width=0.3pt},
  axis line style={InkGray, line width=0.5pt},
  enlarge x limits={abs=0.8},
  legend style={font=\scriptsize, at={(0.5,-0.22)}, anchor=north,
                legend columns=2, draw=none, fill=none, column sep=6pt},
]
\nextgroupplot[
  ybar, bar width=4pt,
  ymin=-1.68, ymax=18.78,
  xtick={0,1,2,3,4,5.7},
  xticklabels={2021,2022,2023,2024,2025,2026H1},
  ylabel={亿元},
]
\addplot[fill=AgriGreen!75, draw=AgriGreen, bar shift=-2.6pt] table[x=i, y=rev]{data/clean/profiles/fin_605198.dat};
\addplot[fill=SoftRed!60, draw=SoftRed, bar shift=2.6pt] table[x=i, y=ni]{data/clean/profiles/fin_605198.dat};
\addplot[fill=AgriGreen!30, draw=AgriGreen, pattern=north east lines, bar shift=-2.6pt] table[x=x, y=rev]{data/clean/profiles/finh_605198.dat};
\addplot[fill=SoftRed!20, draw=SoftRed, pattern=north east lines, bar shift=2.6pt] table[x=x, y=ni]{data/clean/profiles/finh_605198.dat};
\legend{营业收入（年/半年）, 归母净利润（年/半年）}
\nextgroupplot[
  ymin=-0.9, ymax=12.9,
  xtick={0,1,2,3,4,5.7},
  xticklabels={2021,2022,2023,2024,2025,2026H1},
  ylabel={\%},
]
\addplot[AgriGold, mark=*, mark size=1.3pt] table[x=i, y=roe]{data/clean/profiles/fin_605198.dat};
\addplot[OilBlue, mark=square*, mark size=1.3pt] table[x=i, y=debt]{data/clean/profiles/fin_605198.dat};
\addplot[only marks, AgriGold, mark=*, mark size=2.0pt] table[x=x, y=roe]{data/clean/profiles/finh_605198.dat};
\addplot[only marks, OilBlue, mark=square*, mark size=2.0pt] table[x=x, y=debt]{data/clean/profiles/finh_605198.dat};
\legend{ROE, 资产负债率}
\end{groupplot}
\end{tikzpicture}
\begin{tikzpicture}
\begin{axis}[
  width=15.2cm, height=3.1cm, scale only axis,
  ymin=-1.94, ymax=8.41,
  xtick={0,1,2,3,4,5.7},
  xticklabels={2021,2022,2023,2024,2025,2026H1},
  tick label style={font=\tiny},
  yticklabel style={font=\tiny},
  ylabel={亿元}, ylabel style={font=\scriptsize},
  ybar, bar width=4pt,
  grid=major, grid style={LightGray, line width=0.3pt},
  axis line style={InkGray, line width=0.5pt},
  enlarge x limits={abs=0.8},
  legend style={font=\scriptsize, at={(0.5,-0.30)}, anchor=north,
                legend columns=3, draw=none, fill=none, column sep=10pt},
]
\addplot[fill=AgriGreen!55, draw=AgriGreen, bar shift=-3.0pt] table[x=i, y=cash]{data/clean/profiles/fin_605198.dat};
\addplot[fill=SoftRed!45, draw=SoftRed, bar shift=3.0pt] table[x=i, y=ibd]{data/clean/profiles/fin_605198.dat};
\addplot[fill=AgriGreen!25, draw=AgriGreen, pattern=north east lines, bar shift=-3.0pt] table[x=x, y=cash]{data/clean/profiles/finh_605198.dat};
\addplot[fill=SoftRed!20, draw=SoftRed, pattern=north east lines, bar shift=3.0pt] table[x=x, y=ibd]{data/clean/profiles/finh_605198.dat};
\addplot[OilBlue, mark=*, mark size=1.1pt, line width=0.8pt] table[x=i, y=ocf]{data/clean/profiles/fin_605198.dat};
\addplot[only marks, OilBlue, mark=*, mark size=1.8pt] table[x=x, y=ocf]{data/clean/profiles/finh_605198.dat};
\legend{货币资金, 有息负债, 经营活动现金流净额}
\end{axis}
\end{tikzpicture}

\captionof{figure}[安德利（605198）近年主要财务指标]{安德利（605198）近年主要财务指标（左上：营业收入与归母净利润；右上：ROE 与资产负债率；下：货币资金、有息负债与经营活动现金流净额。
2021---2025 年为年报口径，2026H1 为 2026 年半年报/半年末，与其前一年报之间留出间隔、以斜纹或空心点区分；半年报数据未年化）}
\end{center}

\pfl{财务分析} 2026 年 6 月末资产负债率 7.8\%，较 2025 年末 上升 2.5 个百分点（样本中位数 52.2\%，所处三级行业内排第 \zhnumber{1} 位）。 流动比率 8.54、速动比率 5.56，短期指标尚可。 2026 年上半年经营活动现金流净额 5.38 亿元，经营现金流与利润同向为正，内生资金可以支撑运营。 公司披露的财务与合规风险包括：收购公告披露：截至评估基准日，标的公司股东丽水毓甬创业投资合伙企业（有限合伙）持有标的公……。\footnote{\texttt{https://paper.cnstock.com/html/2026-09/19/content\_2270544.htm}}

\pfl{重大战略转型} 从公开披露看，公司的战略推进包括：2023 扩产与并购并行：新建浓缩桃汁、浓缩山楂汁生产线，竞拍取得新疆厂区，产业布局由 6 省 8 厂扩至 7 省 9 厂；2024-05 拟定增投向脱色脱酸浓缩果汁与 NFC 果汁产线，延伸产品结构（该次定增已于 2025 年 5 月终止）；2026-09 以约 7.93 亿元现金收购宁波甬强科技约 62.06\% 股权，跨界电子材料领域，取得控制权并将其纳入合并报表。 公告口径上，近三年（2023 年 9 月---2026 年 9 月）公司披露重大重组与并购类公告 2 条、对外投资与转型类公告 5 条、回购与股权激励类公告 35 条，合计公告 481 条。 第一大行业仍是“饮料制造业”，收入占比 97.2\%，与 2021 年年报相比没有方向性变化。 综合判断，公司在报告期内有明确的外延扩张或业务结构调整动作，转型方向与融资安排之间存在直接关联。\footnote{\texttt{https://www.stcn.com/article/detail/1200611.html}}\footnote{\texttt{https://finance.sina.com.cn/cj/2026-09-21/doc-inispuec0911536.shtml}}

\begin{center}\small\setlength{\tabcolsep}{4pt}
\captionof{table}[安德利（605198）历史融资与资本运作]{安德利（605198）历史融资与资本运作（据公司公告与公开报道整理；金额为披露值，含首次公开发行、再融资、债券、银行借款与重整投资等）}
\begin{tabular}{@{}p{1.5cm}p{3.2cm}p{2.4cm}p{6.3cm}@{}}
\toprule
时间 & 方式 & 规模 & 用途与说明 \\
\midrule
2003-04-22 & H 股 IPO（香港联交所创业板，股份代号 8259） & --- & 国内果汁企业首家在香港上市的公司；2011-01-19 转主板，股份代号 02218 \\
2020-09-18 & A 股 IPO（上海证券交易所） & 发行 20,000,000 股，发行价 7.60 元/股，募集资金总额 15,200.00 万元，扣除发行费用 3,050.00 万元后募集资金净额 12,150.00 万元 & 多品种浓缩果汁生产线建设项目（多品种果汁生产项目）…… \\
2024-05-09 & 拟以简易程序向特定对象发行 A 股股票 & 拟募集资金总额不超过 30,000.00 万元（含本数） & 年产 7,200 吨脱色脱酸浓缩果汁生产线项目、年产 1.2 万吨 NFC 果汁项目…… \\
2024-06—2024-12 & H 股回购 & 累计回购 H 股 780 万股，支付资金总额 7,374.05 万港元（不含佣金等费用） & 注销/回购 H 股股份；2024 年 6 月 6 日至 12 月 31 日期间实施 H …… \\
2026-09 & 现金收购（跨界取得控制权） & 合计约 7.93 亿元收购宁波甬强科技有限公司约 62.06\% 股权 & 取得甬强科技控制权，标的纳入合并报表范围；差异化定价：1.575 亿元收购创始团队 12…… \\
\bottomrule
\end{tabular}
\end{center}

\pfl{历史融投资情况} 从现金流量表看，2025 年公司取得借款收到的现金 — 亿元、偿还债务支付的现金 — 亿元、吸收权益性投资 — 亿元，筹资活动现金流净额 -1.94 亿元（2023 年为取得借款 0.30 亿元、偿还债务 1.60 亿元）。 2026 年上半年筹资活动现金流净额为 -0.97 亿元（取得借款 — 亿元、偿还债务 — 亿元，未年化），已转为净流出，处于净还债状态。 2025 年分配股利、利润或偿付利息支付的现金合计 0.83 亿元。 公开信息中与公司直接相关的融资与资本运作包括：A 股 IPO（上海证券交易所）、H 股 IPO（香港联交所创业板，股份代号 8259）、H 股回购、拟以简易程序向特定对象发行 A 股股票、现金收购（跨界取得控制权）。 其中规模较大的两笔为：2024-05-09拟以简易程序向特定对象发行 A 股股票（拟募集资金总额不超过 30,000.00 万元（含本数））；2020-09-18A 股 IPO（上海证券交易所）（发行 20,000,000 股，发行价 7.60 元/股，募……）。 股本方面，近三年披露股本变动 5 次（限售股份上市 1 次、承诺限售 1 次、股份回购 1 次）；总股本由 2021 年末的 3.67 亿股变为 3.49 亿股（-4.98\%）。 分红方面，近三年现金分红 5 次、累计每股派现 0.92 元（最近一次 2026-07-08）——在利润承压的阶段仍维持了分红。 近三年“再融资”类公告 17 条，涉及定向增发、募集资金使用。\footnote{\texttt{https://www.hkexnews.hk/listedco/listconews/sehk/2026/0327/2026032703935\_c.pdf}}\footnote{\texttt{http://file.finance.sina.com.cn/211.154.219.97:9494/MRGG/CNSESH\_STOCK/2020/2020-9/2020-0}}\footnote{\texttt{http://finance.ce.cn/stock/gsgdbd/202505/t20250521\_2249324.shtml}}

\pfl{未来潜在融资需求分析} 公司 2026 年 6 月末资产负债率 7.8\%、2026 年上半年 ROE 4.3\%（未年化），资产负债表尚有余量，融资需求不是“补血型”的刚性需求，而是围绕并购整合、项目投入与营运资金的主动性需求。 公开披露的下一步融资线索包括：2025-05 终止不超 3 亿元简易程序定增，原因为“目前本次发行上市条件尚不成熟”等市场环境因素（公告口径）；2026-06 收购甬强科技的交易对价初步预估 6—8 亿元，资金来源为自有、自筹等方式；截至 2026-03-31 公司货币资金及交易……。\footnote{\texttt{http://finance.ce.cn/stock/gsgdbd/202505/t20250521\_2249324.shtml}}\footnote{\texttt{http://file.finance.sina.com.cn/211.154.219.97:9494/MRGG/CNSESH\_STOCK/2026/2026-6/2026-0}}


\subsubsection{金龙鱼（300999）}
\label{co:300999}

\pfl{公司基本情况} 益海嘉里旗下粮油食品巨头，主营厨房食品（食用油、大米、面粉）与饲料原料、油脂科技，旗下拥有金龙鱼、胡姬花、香满园、欧丽薇兰等品牌；2020 年登陆创业板并创下当时创业板 IPO 募资最大规模，2023—2024 年利润连续下滑后于 2025 年触底回暖。 公司于 2020-10-15 在深交所创业板首发上市，注册资本 54.22 亿元，总股本 54.22 亿股。 截至 2026-09-24收盘价 24.55 元，流通市值 133.2 亿元，近一年-22.5\%，流通市值在三级行业“粮油加工”的 7 家公司中排第 \zhnumber{1} 位，近一年涨跌幅低于全样本中位数（-10.1\%）12.4 个百分点。 按 2025 年年报披露的行业构成，收入占比前三位为农副食品加工业 100.0\%。 发展过程中的关键节点包括：上市以来首次出现营业收入与净利润“双降”：营业收入 2,515.24 亿元，同比下降 2……（2023）；公司名称由“益海嘉里金龙鱼粮油食品股份有限公司”变更为“益海嘉里金龙鱼食品集团股份有限公……（2024）；营业收入 2,388.66 亿元，同比下降 5.03\%；利润总额 36.56 亿元，同比……（2024）。\footnote{\texttt{http://static.cninfo.com.cn/finalpage/2024-02-21/1219146870.PDF}}\footnote{\texttt{http://money.finance.sina.com.cn/corp/view/vCB\_AllBulletinDetail.php}}\footnote{\texttt{http://money.finance.sina.com.cn/corp/view/vCB\_AllBulletinDetail.php}}

\pfl{历史股价} 本报告样本区间（2021 年 1 月 至 2026 年 9 月）内，股价由 121.92 元变为 24.55 元，累计 -79.9\%，区间最高 121.92 元（2021 年 1 月）、最低 23.26 元（2026 年 6 月）；当前价相当于区间最高价的 20.1\%，为区间最低价的 1.06 倍。 月度收益的年化波动率 32.2\%，低于全样本中位数 37.4\%，在三级行业“粮油加工”的 7 家公司中波动率排第 \zhnumber{4} 位。 从事件看，2024 年 8 月 至 2024 年 9 月的上涨（+37.2\%）与公司2024 年 9 月 3 日董事会审议通过；交易完成后公司持有鲁花集团 10……（2024-09）在时间上重合——股价的拐点多数能在公司自身的资本运作与风险事件上找到对应。\footnote{\texttt{https://finance.sina.com.cn/stock/relnews/cn/2024-12-09/doc-incywqmm5688220.shtml}}

\begin{center}
\begin{tikzpicture}
\begin{axis}[
  width=12.6cm, height=3.9cm, scale only axis,
  xmin=2020.922, xmax=2026.888, ymin=21.63, ymax=130.45,
  xtick={2021,2022,2023,2024,2025,2026},
  yearticks,
  yticklabel style={font=\scriptsize},
  ylabel={元/股}, ylabel style={font=\scriptsize},
  grid=major, grid style={LightGray, line width=0.3pt},
  axis line style={InkGray, line width=0.5pt},
  every axis plot/.append style={line width=0.9pt},
]
\addplot[AgriGreen] table[x=x,y=y]{data/clean/profiles/px_300999.dat};
\addplot[SoftRed, dashed, line width=0.7pt] table[x=x,y=y]{data/clean/profiles/pxzz_300999.dat};
\addplot[only marks, mark=*, mark size=1.1pt, SoftRed] table[x=x,y=y]{data/clean/profiles/pxzz_300999.dat};
\end{axis}
\end{tikzpicture}

\captionof{figure}[金龙鱼（300999）股价与周期骨架]{金龙鱼（300999）月度收盘价（元/股）与 ZigZag 周期骨架（阈值 25\%，圆点为周期转折点）}
\end{center}

\pfl{过去周期分析} 以 25\% 的 ZigZag 阈值划分，样本区间内股价形成 1 个主升段与 2 个主跌段。 最大主升段出现在 2024 年 8 月 至 2024 年 9 月，1 个月+37.2\%；最大主跌段为 2021 年 1 月 至 2024 年 8 月，43 个月-78.8\%。 最近一次转折出现在 2026 年 6 月（下跌段自 2024 年 9 月起，21 个月-34.5\%），当前处于该段的延续之中。 公司股价与申万农林牧渔指数几乎同步见底，个股并未走出独立行情。 粮油加工是“成本加成”型生意（第 \ref{sec:grain} 章），利润来自原料与成品的价差及规模效率，收入规模大而净利率薄，公司 2026 年上半年实现盈利、半年 ROE 2.4\%（未年化），好于全样本中位数（0.75\%），下行周期对其报表的冲击小于同业。

\pfl{盈利情况} 2026 年上半年营业收入 1,228.52 亿元（同比 +6.2\%），归母净利润 22.94 亿元，同比 +30.7\%。 以年报为趋势对照，2023---2025 年营业收入由 2,515.24 亿元变为 2,451.26 亿元（两年年均复合增速 -1.3\%），归母净利润由 28.48 亿元变为 31.53 亿元。 按半年报口径，2026 年上半年的 ROE 为 2.4\%（未年化）、销售净利率 1.9\%、毛利率 6.3\%，ROE 在“粮油加工”7 家公司中排第 \zhnumber{3} 位。 同期全样本半年 ROE 中位数 0.75\%，公司与之相差 1.61 个百分点（略好于中位数公司）。 结合公司披露，两起子公司诉讼合计计提预计负债损失约 7.33 亿元，均为一审未生效判决、相关子公司已提起上诉，最终结果……；下属子公司益海（广州）粮油工业有限公司收到安徽省淮北市中级人民法院一审《刑事判决书》，被判决构成合同诈骗……。\footnote{\texttt{https://finance.sina.com.cn/stock/zqgd/2026-03-04/doc-inhpvpqk3863433.shtml}}

\begin{center}\small\setlength{\tabcolsep}{4pt}
\captionof{table}[金龙鱼（300999）关键财务指标]{金龙鱼（300999）关键财务指标（合并报表口径；两个半年报期均未年化；现金短债比 $=$ 货币资金 $\div$ 短期债务）}
\begin{tabular}{@{}lrrrrr@{}}
\toprule
指标 & 2023 年 & 2024 年 & 2025 年 & 2025 年半年报 & 2026 年半年报 \\
\midrule
营业收入（亿元） & 2,515.24 & 2,382.79 & 2,451.26 & 1,156.82 & 1,228.52 \\
归母净利润（亿元） & 28.48 & 25.02 & 31.53 & 17.56 & 22.94 \\
毛利率（\%） & 4.8 & 5.4 & 6.4 & 6.5 & 6.3 \\
ROE（\%） & 3.2 & 2.7 & 3.3 & 1.9 & 2.4 \\
资产负债率（\%） & 59.8 & 56.4 & 55.9 & 55.9 & 56.2 \\
经营活动现金流净额（亿元） & 149.21 & 60.38 & 154.45 & 146.25 & 10.01 \\
货币资金（亿元） & 592.96 & 529.06 & 442.04 & 443.72 & 485.84 \\
有息负债合计（亿元） & 1,157.03 & 943.10 & 864.73 & 887.24 & 927.22 \\
现金短债比（倍） & 0.56 & 0.60 & 0.54 & 0.53 & 0.55 \\
\bottomrule
\end{tabular}
\end{center}

\begin{center}
\begin{tikzpicture}
\begin{groupplot}[
  group style={group size=2 by 1, horizontal sep=0.95cm},
  width=7.2cm, height=4.6cm,
  tick label style={font=\tiny},
  title style={font=\scriptsize\heiti},
  yticklabel style={font=\tiny},
  grid=major, grid style={LightGray, line width=0.3pt},
  axis line style={InkGray, line width=0.5pt},
  enlarge x limits={abs=0.8},
  legend style={font=\scriptsize, at={(0.5,-0.22)}, anchor=north,
                legend columns=2, draw=none, fill=none, column sep=6pt},
]
\nextgroupplot[
  ybar, bar width=4pt,
  ymin=-257.49, ymax=2883.84,
  xtick={0,1,2,3,4,5.7},
  xticklabels={2021,2022,2023,2024,2025,2026H1},
  ylabel={亿元},
]
\addplot[fill=AgriGreen!75, draw=AgriGreen, bar shift=-2.6pt] table[x=i, y=rev]{data/clean/profiles/fin_300999.dat};
\addplot[fill=SoftRed!60, draw=SoftRed, bar shift=2.6pt] table[x=i, y=ni]{data/clean/profiles/fin_300999.dat};
\addplot[fill=AgriGreen!30, draw=AgriGreen, pattern=north east lines, bar shift=-2.6pt] table[x=x, y=rev]{data/clean/profiles/finh_300999.dat};
\addplot[fill=SoftRed!20, draw=SoftRed, pattern=north east lines, bar shift=2.6pt] table[x=x, y=ni]{data/clean/profiles/finh_300999.dat};
\legend{营业收入（年/半年）, 归母净利润（年/半年）}
\nextgroupplot[
  ymin=-4.8, ymax=65.7,
  xtick={0,1,2,3,4,5.7},
  xticklabels={2021,2022,2023,2024,2025,2026H1},
  ylabel={\%},
]
\addplot[AgriGold, mark=*, mark size=1.3pt] table[x=i, y=roe]{data/clean/profiles/fin_300999.dat};
\addplot[OilBlue, mark=square*, mark size=1.3pt] table[x=i, y=debt]{data/clean/profiles/fin_300999.dat};
\addplot[only marks, AgriGold, mark=*, mark size=2.0pt] table[x=x, y=roe]{data/clean/profiles/finh_300999.dat};
\addplot[only marks, OilBlue, mark=square*, mark size=2.0pt] table[x=x, y=debt]{data/clean/profiles/finh_300999.dat};
\legend{ROE, 资产负债率}
\end{groupplot}
\end{tikzpicture}
\begin{tikzpicture}
\begin{axis}[
  width=15.2cm, height=3.1cm, scale only axis,
  ymin=-115.70, ymax=1295.87,
  xtick={0,1,2,3,4,5.7},
  xticklabels={2021,2022,2023,2024,2025,2026H1},
  tick label style={font=\tiny},
  yticklabel style={font=\tiny},
  ylabel={亿元}, ylabel style={font=\scriptsize},
  ybar, bar width=4pt,
  grid=major, grid style={LightGray, line width=0.3pt},
  axis line style={InkGray, line width=0.5pt},
  enlarge x limits={abs=0.8},
  legend style={font=\scriptsize, at={(0.5,-0.30)}, anchor=north,
                legend columns=3, draw=none, fill=none, column sep=10pt},
]
\addplot[fill=AgriGreen!55, draw=AgriGreen, bar shift=-3.0pt] table[x=i, y=cash]{data/clean/profiles/fin_300999.dat};
\addplot[fill=SoftRed!45, draw=SoftRed, bar shift=3.0pt] table[x=i, y=ibd]{data/clean/profiles/fin_300999.dat};
\addplot[fill=AgriGreen!25, draw=AgriGreen, pattern=north east lines, bar shift=-3.0pt] table[x=x, y=cash]{data/clean/profiles/finh_300999.dat};
\addplot[fill=SoftRed!20, draw=SoftRed, pattern=north east lines, bar shift=3.0pt] table[x=x, y=ibd]{data/clean/profiles/finh_300999.dat};
\addplot[OilBlue, mark=*, mark size=1.1pt, line width=0.8pt] table[x=i, y=ocf]{data/clean/profiles/fin_300999.dat};
\addplot[only marks, OilBlue, mark=*, mark size=1.8pt] table[x=x, y=ocf]{data/clean/profiles/finh_300999.dat};
\legend{货币资金, 有息负债, 经营活动现金流净额}
\end{axis}
\end{tikzpicture}

\captionof{figure}[金龙鱼（300999）近年主要财务指标]{金龙鱼（300999）近年主要财务指标（左上：营业收入与归母净利润；右上：ROE 与资产负债率；下：货币资金、有息负债与经营活动现金流净额。
2021---2025 年为年报口径，2026H1 为 2026 年半年报/半年末，与其前一年报之间留出间隔、以斜纹或空心点区分；半年报数据未年化）}
\end{center}

\pfl{财务分析} 2026 年 6 月末资产负债率 56.2\%，较 2025 年末 上升 0.4 个百分点（样本中位数 52.2\%，所处三级行业内排第 \zhnumber{6} 位）。 2026 年 6 月末货币资金 485.84 亿元、短期债务（短期借款与一年内到期非流动负债）887.12 亿元、长期债务 40.10 亿元，有息负债合计 927.22 亿元，现金短债比 0.55 倍——覆盖偏紧。 净有息负债 441.38 亿元，相当于其 2026 年上半年营业收入的 36\%。 流动比率 1.20、速动比率 0.78，短期指标尚可。 2026 年上半年经营活动现金流净额 10.01 亿元，经营现金流与利润同向为正，内生资金可以支撑运营。 2026 年 6 月末在建工程 42.03 亿元，较上年末增加 0.83 亿元，仍有在建投入。 2026 年 6 月末商誉 61.11 亿元，占归母净资产的 6.0\%，是净资产中最脆弱的一块。 公司披露的财务与合规风险包括：下属子公司益海（广州）粮油工业有限公司收到安徽省淮北市中级人民法院一审《刑事判决书》，被……；两起子公司诉讼合计计提预计负债损失约 7.33 亿元，均为一审未生效判决、相关子公司已提……。\footnote{\texttt{https://finance.sina.com.cn/stock/zqgd/2026-03-04/doc-inhpvpqk3863433.shtml}}

\pfl{重大战略转型} 公司在报告期内的战略动作可归为 2 项：2024-09 以参股公司股权作价增资鲁花集团，公司与关联方合计取得鲁花集团 26.64\% 股权，推动食用油行业整合；2026-08 2026 年上半年油米面产品在消费品、餐饮、食品工业及电商渠道销量均有所增长；控股股东连续四年自愿延长股份锁定期……。 公告口径上，近三年（2023 年 9 月---2026 年 9 月）公司披露重大重组与并购类公告 0 条、对外投资与转型类公告 3 条、回购与股权激励类公告 14 条，合计公告 306 条。 第一大行业仍是“农副食品加工业”，收入占比 100.0\%，与 2021 年年报相比没有方向性变化。 综合判断，公司在报告期内有明确的外延扩张或业务结构调整动作，转型方向与融资安排之间存在直接关联。\footnote{\texttt{https://www.21jingji.com/article/20240905/herald/1d1ba39e0c0f48de29504d91eaab03e3.html}}\footnote{\texttt{https://finance.sina.com.cn/roll/2026-08-12/doc-ininamih5176022.shtml}}

\pfl{历史融投资情况} 从现金流量表看，2025 年公司取得借款收到的现金 1,569.21 亿元、偿还债务支付的现金 1,649.00 亿元、吸收权益性投资 0.70 亿元，筹资活动现金流净额 -46.92 亿元（2023 年为取得借款 1,395.70 亿元、偿还债务 1,225.17 亿元）。 2026 年上半年筹资活动现金流净额为 52.24 亿元（取得借款 870.26 亿元、偿还债务 806.36 亿元，未年化），仍处于净融入状态。 2025 年分配股利、利润或偿付利息支付的现金合计 25.46 亿元。 公司披露的股本与债务融资脉络包括：IPO（深圳证券交易所创业板）、以参股公司股权作价对外增资（关联交易）。 其中规模较大的两笔为：2020-10IPO（深圳证券交易所创业板）（首次公开发行 542,159,154 股，发行价 25.70……）；2024-09以参股公司股权作价对外增资（关联交易）（公司出资作价 22.90 亿元、关联方香港嘉银作价 32.7……）。 股本方面，近三年披露股本变动 3 次；总股本由 2021 年末的 54.22 亿股变为 54.22 亿股（+0.00\%）。 分红方面，近三年现金分红 4 次、累计每股派现 0.53 元（最近一次 2025-12-31）——在利润承压的阶段仍维持了分红。\footnote{\texttt{http://money.finance.sina.com.cn/corp/view/vCB\_AllBulletinDetail.php}}\footnote{\texttt{https://finance.sina.com.cn/stock/relnews/cn/2024-12-09/doc-incywqmm5688220.shtml}}

\pfl{未来潜在融资需求分析} 2026 年 6 月末货币资金 485.84 亿元，而短期债务 887.12 亿元（现金短债比 0.55 倍），2026 年上半年经营现金流净额 10.01 亿元。按“短期债务减去账面货币资金”的滚动偿付口径，静态缺口约 401.28 亿元；若再计入上半年亏损的年化额（0.00 亿元），维持一年正常经营所需的外部资金约 401.28---401.28 亿元。 公司 2026 年 6 月末资产负债率 56.2\%、2026 年上半年 ROE 2.4\%（未年化），资产负债表尚有余量，融资需求不是“补血型”的刚性需求，而是围绕并购整合、项目投入与营运资金的主动性需求。 公开披露的下一步融资线索包括：2020—2024 IPO 募集资金净额 136.93 亿元，按公告披露投向兰州新区粮油食品加工基地（一期）等四个项目；2024 2024 年报披露报告期投资额 821.63 亿元，上年同期 675.65 亿元，同比增长 21.61\%（投资口径，非融……。 资金需求还要叠加在建工程：期末在建工程 42.03 亿元、2026 年上半年购建固定资产等支出 20.91 亿元（未年化），这部分支出在周期底部通常会被主动放缓。\footnote{\texttt{http://money.finance.sina.com.cn/corp/view/vCB\_AllBulletinDetail.php}}


\subsubsection{金健米业（600127）}
\label{co:600127}

\pfl{公司基本情况} 公司前身可追溯至 1998 年，湖南常德地方国资背景的粮油食品加工企业，主营粮油食品加工、农产品贸易与乳品，控股股东为湖南粮食集团，实际控制人为湖南省国资委；收入规模数十亿元但盈利极薄，累计可供分配利润为负，2024 年与控股股东完成资产置换后主业结构调整。 公司于 1998-05-06 在上交所首发上市，注册资本 6.42 亿元，总股本 6.42 亿股。 截至 2026-09-24收盘价 16.49 元，流通市值 105.8 亿元，近一年+146.1\%，流通市值在三级行业“粮油加工”的 7 家公司中排第 \zhnumber{2} 位，近一年涨跌幅高于全样本中位数（-10.1\%）156.2 个百分点。 按 2026 年半年报披露的行业构成，收入占比前三位为粮油业务 94.5\%；食品行业 3.4\%；乳品业务 2.3\%。 与 2021 年年报相比，第一大行业由“粮油食品加工业”变为“粮油业务”，主业重心已经迁移。 发展过程中的关键节点包括：1998 年 4 月 8 日至 15 日以“全额预缴款、比例配售、余款即退”方式向社会发……（1998）；经中国证监会证监公司字[2000]229 号文批准，2000 年 12 月 29 日至 ……（2001）；2021 年 7 月公司实际控制人由长沙市国资委变更为湖南省国资委；划转前湖南粮食集团通……（2021）。\footnote{\texttt{https://vip.stock.finance.sina.com.cn/corp/go.php/vISSUE\_NewStock/stockid/600127.phtml}}\footnote{\texttt{http://www.sse.com.cn/disclosure/listedinfo/announcement/c/2003-04-29/600127\_2002\_n.pdf}}\footnote{\texttt{http://static.cninfo.com.cn/finalpage/2021-07-16/1210496470.PDF}}

\pfl{历史股价} 自 2021 年 1 月至 2026 年 9 月，股价由 9.10 元变为 16.49 元，累计 +81.2\%，区间最高 16.49 元（2026 年 9 月）、最低 5.20 元（2026 年 6 月）；当前价相当于区间最高价的 100.0\%，为区间最低价的 3.17 倍。 月度收益的年化波动率 45.3\%，高于全样本中位数 37.4\%，在三级行业“粮油加工”的 7 家公司中波动率排第 \zhnumber{1} 位。

\begin{center}
\begin{tikzpicture}
\begin{axis}[
  width=12.6cm, height=3.9cm, scale only axis,
  xmin=2020.922, xmax=2026.888, ymin=4.84, ymax=17.64,
  xtick={2021,2022,2023,2024,2025,2026},
  yearticks,
  yticklabel style={font=\scriptsize},
  ylabel={元/股}, ylabel style={font=\scriptsize},
  grid=major, grid style={LightGray, line width=0.3pt},
  axis line style={InkGray, line width=0.5pt},
  every axis plot/.append style={line width=0.9pt},
]
\addplot[AgriGreen] table[x=x,y=y]{data/clean/profiles/px_600127.dat};
\addplot[SoftRed, dashed, line width=0.7pt] table[x=x,y=y]{data/clean/profiles/pxzz_600127.dat};
\addplot[only marks, mark=*, mark size=1.1pt, SoftRed] table[x=x,y=y]{data/clean/profiles/pxzz_600127.dat};
\end{axis}
\end{tikzpicture}

\captionof{figure}[金健米业（600127）股价与周期骨架]{金健米业（600127）月度收盘价（元/股）与 ZigZag 周期骨架（阈值 25\%，圆点为周期转折点）}
\end{center}

\pfl{过去周期分析} 以 25\% 的 ZigZag 阈值划分，样本区间内股价形成 3 个主升段与 2 个主跌段。 最大主升段出现在 2026 年 6 月 至 2026 年 9 月，3 个月+217.1\%；最大主跌段为 2022 年 5 月 至 2024 年 7 月，26 个月-42.1\%。 最近一次转折出现在 2026 年 9 月（上涨段自 2026 年 6 月起，3 个月+217.1\%），当前处于该段之后的震荡之中。 公司股价与申万农林牧渔指数几乎同步见底，个股并未走出独立行情。 粮油加工是“成本加成”型生意（第 \ref{sec:grain} 章），利润来自原料与成品的价差及规模效率，收入规模大而净利率薄，公司 2026 年 6 月末资产负债率 62.2\%、半年 ROE -1.4\%（未年化），在本轮下行中属于随行业同步调整的一类。

\pfl{盈利情况} 2026 年上半年营业收入 19.70 亿元（同比 +25.2\%），归母净利润 -0.10 亿元，由上年同期的盈利转为亏损。 以年报为趋势对照，2023---2025 年营业收入由 52.69 亿元变为 33.58 亿元（两年年均复合增速 -20.2\%），归母净利润由 0.14 亿元变为 0.04 亿元。 按半年报口径，2026 年上半年的 ROE 为 -1.4\%（未年化）、销售净利率 -0.5\%、毛利率 6.0\%，ROE 在“粮油加工”7 家公司中排第 \zhnumber{7} 位。 同期全样本半年 ROE 中位数 0.75\%，公司与之相差 2.19 个百分点（弱于中位数公司）。 结合公司披露，营业收入 462,716 万元，同比下降 12.19\%；归母净利润 223 万元，同比下降 83.94\%……；2024 年累计实现营业收入 462,716 万元，较上年度 526,943 万元减少 64,227 万……。\footnote{\texttt{https://www.sse.com.cn/disclosure/listedinfo/announcement/c/new/2025-04-15/600127\_202504}}

\begin{center}\small\setlength{\tabcolsep}{4pt}
\captionof{table}[金健米业（600127）关键财务指标]{金健米业（600127）关键财务指标（合并报表口径；两个半年报期均未年化；现金短债比 $=$ 货币资金 $\div$ 短期债务）}
\begin{tabular}{@{}lrrrrr@{}}
\toprule
指标 & 2023 年 & 2024 年 & 2025 年 & 2025 年半年报 & 2026 年半年报 \\
\midrule
营业收入（亿元） & 52.69 & 46.27 & 33.58 & 15.74 & 19.70 \\
归母净利润（亿元） & 0.14 & 0.02 & 0.04 & 0.12 & -0.10 \\
毛利率（\%） & 4.6 & 5.8 & 7.7 & 8.3 & 6.0 \\
ROE（\%） & 1.6 & 0.3 & 0.6 & 1.7 & -1.4 \\
资产负债率（\%） & 54.8 & 59.4 & 61.3 & 57.4 & 62.2 \\
经营活动现金流净额（亿元） & 1.90 & -0.44 & 0.72 & 1.07 & 0.38 \\
货币资金（亿元） & 1.49 & 1.45 & 1.86 & 1.57 & 1.50 \\
有息负债合计（亿元） & 7.12 & 7.37 & 7.78 & 6.75 & 7.31 \\
现金短债比（倍） & 0.23 & 0.20 & 0.24 & 0.23 & 0.21 \\
\bottomrule
\end{tabular}
\end{center}

\begin{center}
\begin{tikzpicture}
\begin{groupplot}[
  group style={group size=2 by 1, horizontal sep=0.95cm},
  width=7.2cm, height=4.6cm,
  tick label style={font=\tiny},
  title style={font=\scriptsize\heiti},
  yticklabel style={font=\tiny},
  grid=major, grid style={LightGray, line width=0.3pt},
  axis line style={InkGray, line width=0.5pt},
  enlarge x limits={abs=0.8},
  legend style={font=\scriptsize, at={(0.5,-0.22)}, anchor=north,
                legend columns=2, draw=none, fill=none, column sep=6pt},
]
\nextgroupplot[
  ybar, bar width=4pt,
  ymin=-7.00, ymax=72.20,
  xtick={0,1,2,3,4,5.7},
  xticklabels={2021,2022,2023,2024,2025,2026H1},
  ylabel={亿元},
]
\addplot[fill=AgriGreen!75, draw=AgriGreen, bar shift=-2.6pt] table[x=i, y=rev]{data/clean/profiles/fin_600127.dat};
\addplot[fill=SoftRed!60, draw=SoftRed, bar shift=2.6pt] table[x=i, y=ni]{data/clean/profiles/fin_600127.dat};
\addplot[fill=AgriGreen!30, draw=AgriGreen, pattern=north east lines, bar shift=-2.6pt] table[x=x, y=rev]{data/clean/profiles/finh_600127.dat};
\addplot[fill=SoftRed!20, draw=SoftRed, pattern=north east lines, bar shift=2.6pt] table[x=x, y=ni]{data/clean/profiles/finh_600127.dat};
\legend{营业收入（年/半年）, 归母净利润（年/半年）}
\nextgroupplot[
  ymin=-13.0, ymax=70.7,
  xtick={0,1,2,3,4,5.7},
  xticklabels={2021,2022,2023,2024,2025,2026H1},
  ylabel={\%},
]
\addplot[AgriGold, mark=*, mark size=1.3pt] table[x=i, y=roe]{data/clean/profiles/fin_600127.dat};
\addplot[OilBlue, mark=square*, mark size=1.3pt] table[x=i, y=debt]{data/clean/profiles/fin_600127.dat};
\addplot[only marks, AgriGold, mark=*, mark size=2.0pt] table[x=x, y=roe]{data/clean/profiles/finh_600127.dat};
\addplot[only marks, OilBlue, mark=square*, mark size=2.0pt] table[x=x, y=debt]{data/clean/profiles/finh_600127.dat};
\legend{ROE, 资产负债率}
\end{groupplot}
\end{tikzpicture}
\begin{tikzpicture}
\begin{axis}[
  width=15.2cm, height=3.1cm, scale only axis,
  ymin=-1.44, ymax=10.69,
  xtick={0,1,2,3,4,5.7},
  xticklabels={2021,2022,2023,2024,2025,2026H1},
  tick label style={font=\tiny},
  yticklabel style={font=\tiny},
  ylabel={亿元}, ylabel style={font=\scriptsize},
  ybar, bar width=4pt,
  grid=major, grid style={LightGray, line width=0.3pt},
  axis line style={InkGray, line width=0.5pt},
  enlarge x limits={abs=0.8},
  legend style={font=\scriptsize, at={(0.5,-0.30)}, anchor=north,
                legend columns=3, draw=none, fill=none, column sep=10pt},
]
\addplot[fill=AgriGreen!55, draw=AgriGreen, bar shift=-3.0pt] table[x=i, y=cash]{data/clean/profiles/fin_600127.dat};
\addplot[fill=SoftRed!45, draw=SoftRed, bar shift=3.0pt] table[x=i, y=ibd]{data/clean/profiles/fin_600127.dat};
\addplot[fill=AgriGreen!25, draw=AgriGreen, pattern=north east lines, bar shift=-3.0pt] table[x=x, y=cash]{data/clean/profiles/finh_600127.dat};
\addplot[fill=SoftRed!20, draw=SoftRed, pattern=north east lines, bar shift=3.0pt] table[x=x, y=ibd]{data/clean/profiles/finh_600127.dat};
\addplot[OilBlue, mark=*, mark size=1.1pt, line width=0.8pt] table[x=i, y=ocf]{data/clean/profiles/fin_600127.dat};
\addplot[only marks, OilBlue, mark=*, mark size=1.8pt] table[x=x, y=ocf]{data/clean/profiles/finh_600127.dat};
\legend{货币资金, 有息负债, 经营活动现金流净额}
\end{axis}
\end{tikzpicture}

\captionof{figure}[金健米业（600127）近年主要财务指标]{金健米业（600127）近年主要财务指标（左上：营业收入与归母净利润；右上：ROE 与资产负债率；下：货币资金、有息负债与经营活动现金流净额。
2021---2025 年为年报口径，2026H1 为 2026 年半年报/半年末，与其前一年报之间留出间隔、以斜纹或空心点区分；半年报数据未年化）}
\end{center}

\pfl{财务分析} 2026 年 6 月末资产负债率 62.2\%，较 2025 年末 上升 0.9 个百分点（样本中位数 52.2\%，所处三级行业内排第 \zhnumber{7} 位）。 2026 年 6 月末货币资金 1.50 亿元、短期债务（短期借款与一年内到期非流动负债）7.28 亿元、长期债务 0.03 亿元，有息负债合计 7.31 亿元，现金短债比 0.21 倍——缺口明显。 净有息负债 5.81 亿元，相当于其 2026 年上半年营业收入的 30\%。 流动比率 0.97、速动比率 0.47，短期偿付压力较大。 2026 年上半年经营活动现金流净额 0.38 亿元，净利润为负而经营现金流为正，差额主要来自折旧摊销与营运资本变动，说明亏损尚未完全转化为现金流出。 公司披露的财务与合规风险包括：收到上海证券交易所《关于金健米业股份有限公司 2024 年年度报告的信息披露监管问询函》……。\footnote{\texttt{http://static.cninfo.com.cn/finalpage/2025-06-30/1224029813.PDF}}

\pfl{重大战略转型} 近三年公司的战略动作集中在以下几个方面：2021-07-15 湖南省国资委在《金健米业股份有限公司详式权益变动报告书》中承诺，在符合公司发展实际、不损害中小股东权益的前提下，……；2024-11 以资产置换方式优化产业布局：置出从事贸易业务的进出口与农产品贸易子公司，置入湖南裕湘食品与中南粮油食品科学研究院……；2026-06 2026 年半年度报告披露公司存在与业绩承诺相关的补偿安排：相关标的业绩未达承诺则现金补偿，承诺期至 2027 ……。 公告口径上，近三年（2023 年 9 月---2026 年 9 月）公司披露重大重组与并购类公告 4 条、对外投资与转型类公告 2 条、回购与股权激励类公告 0 条，合计公告 359 条。 主营结构的迁移已经落到报表上：2021 年年报的第一大行业为“粮油食品加工业”（48.1\%），2026 年半年报变为“粮油业务”（94.5\%）。 综合判断，公司在报告期内有明确的外延扩张或业务结构调整动作，转型方向与融资安排之间存在直接关联。\footnote{\texttt{https://static.cninfo.com.cn/finalpage/2025-08-26/1224566240.PDF}}\footnote{\texttt{https://www.sse.com.cn/disclosure/listedinfo/announcement/c/new/2024-12-28/600127\_202412}}\footnote{\texttt{https://stockmc.xueqiu.com/202608/600127\_20260829\_VIPV.pdf}}

\begin{center}\small\setlength{\tabcolsep}{4pt}
\captionof{table}[金健米业（600127）历史融资与资本运作]{金健米业（600127）历史融资与资本运作（据公司公告与公开报道整理；金额为披露值，含首次公开发行、再融资、债券、银行借款与重整投资等）}
\begin{tabular}{@{}p{1.5cm}p{3.2cm}p{2.4cm}p{6.3cm}@{}}
\toprule
时间 & 方式 & 规模 & 用途与说明 \\
\midrule
1998-04 & IPO（上海证券交易所） & 发行公众股 5,000 万股，发行价 5.30 元/股，实际募集资金合计 26,501.00 万元、募集资金净额 25,500.00 万元 & 含公司职工股 500 万股；发行方式为全额预缴款、比例配售、余款即退…… \\
2000-12—2001-01 & 公募增发 A 股（上海证券交易所） & 增发 5,000 万股，发行价格 16.10 元/股 & 经中国证监会证监公司字[2000]229 号文批准…… \\
2024-11—2024-12 & 与控股股东资产置换（关联交易） & 公司向控股股东支付资产置换差价 8,361.76 万元 & 置出湖南金健进出口有限责任公司 100\% 股权…… \\
\bottomrule
\end{tabular}
\end{center}

\pfl{历史融投资情况} 从现金流量表看，2025 年公司取得借款收到的现金 5.86 亿元、偿还债务支付的现金 5.44 亿元、吸收权益性投资 — 亿元，筹资活动现金流净额 0.28 亿元（2023 年为取得借款 6.21 亿元、偿还债务 8.86 亿元）。 2026 年上半年筹资活动现金流净额为 -0.53 亿元（取得借款 2.62 亿元、偿还债务 3.10 亿元，未年化），已转为净流出，处于净还债状态。 2025 年分配股利、利润或偿付利息支付的现金合计 0.16 亿元。 公开信息中与公司直接相关的融资与资本运作包括：IPO（上海证券交易所）、与控股股东资产置换（关联交易）、公募增发 A 股（上海证券交易所）。 其中规模较大的两笔为：1998-04IPO（上海证券交易所）（发行公众股 5,000 万股，发行价 5.30 元/股，实际……）；2024-11—2024-12与控股股东资产置换（关联交易）（公司向控股股东支付资产置换差价 8,361.76 万元）。 股本方面，近三年披露股本变动 3 次；总股本由 2021 年末的 6.42 亿股变为 6.42 亿股（+0.00\%）。 分红方面，近三年未实施现金分红，在全部样本中属于分红能力偏弱的一端。\footnote{\texttt{https://vip.stock.finance.sina.com.cn/corp/go.php/vISSUE\_NewStock/stockid/600127.phtml}}\footnote{\texttt{http://www.sse.com.cn/disclosure/listedinfo/announcement/c/2003-04-29/600127\_2002\_n.pdf}}\footnote{\texttt{https://www.sse.com.cn/disclosure/listedinfo/announcement/c/new/2025-06-30/600127\_202506}}

\pfl{未来潜在融资需求分析} 2026 年 6 月末货币资金 1.50 亿元，而短期债务 7.28 亿元（现金短债比 0.21 倍），2026 年上半年经营现金流净额 0.38 亿元。按“短期债务减去账面货币资金”的滚动偿付口径，静态缺口约 5.78 亿元；若再计入上半年亏损的年化额（0.19 亿元），维持一年正常经营所需的外部资金约 5.78---5.97 亿元。 公司 2026 年上半年归母净利润为负（-0.10 亿元），但 6 月末资产负债率 62.2\% 尚未进入第一档（亏损且负债率 $>$65\%）的区间，当前融资需求以补充流动资金、支撑低谷期经营性支出为主。 公开披露的下一步融资线索包括：2021—2026 公司实际控制人变更为湖南省国资委后形成同业竞争，控股股东及国资方就解决同业竞争作出承诺（截至 2025 年半年报仍在履行……；2024-12 资产置换中公司向控股股东支付置换差价 8,361.76 万元（公告中作为模拟测算的现金支付安排披露）。\footnote{\texttt{https://static.cninfo.com.cn/finalpage/2025-08-26/1224566240.PDF}}\footnote{\texttt{https://www.sse.com.cn/disclosure/listedinfo/announcement/c/new/2025-06-30/600127\_202506}}


\subsubsection{中粮科技（000930）}
\label{co:000930}

\pfl{公司基本情况} 1998 年设立至今，中粮科技原名安徽丰原生物化学，现为中粮集团旗下玉米深加工与生物能源平台，主营燃料乙醇、淀粉及淀粉糖、味精等产品，国内燃料乙醇市场份额居前；2025 年盈利微薄，2026 年上半年因产品价格下行转为亏损。 公司于 1999-07-12 在深交所主板首发上市，注册资本 18.56 亿元，总股本 18.60 亿股。 截至 2026-09-24收盘价 5.26 元，流通市值 97.6 亿元，近一年-11.9\%，流通市值在三级行业“其他农产品加工”的 11 家公司中排第 \zhnumber{2} 位，近一年涨跌幅低于全样本中位数（-10.1\%）1.8 个百分点。 按 2026 年半年报披露的行业构成，收入占比前三位为农产品加工及销售 98.3\%；其他(补充) 1.7\%。 发展过程中的关键节点包括：经中国证监会证监发行字〔1999〕67 号文件批准，1999 年 7 月 12 日股票在……（1999）；经中国证监会证监公司字〔2003〕36 号文件批准，2003 年 4 月 24 日向社会……（2003）；经 2005 年第二次临时股东大会决议，以股权登记日（2005-12-06）总股本 48……（2005）。\footnote{\texttt{https://vip.stock.finance.sina.com.cn/corp/go.php/vISSUE\_NewStock/stockid/000930.phtml}}\footnote{\texttt{http://money.finance.sina.com.cn/corp/view/vISSUE\_MarketBulletinDetail.php}}\footnote{\texttt{https://vip.stock.finance.sina.com.cn/corp/view/vCB\_AllBulletinDetail.php}}

\pfl{历史股价} 本报告样本区间（2021 年 1 月 至 2026 年 9 月）内，股价由 10.89 元变为 5.26 元，累计 -51.7\%，区间最高 11.68 元（2021 年 11 月）、最低 4.25 元（2026 年 6 月）；当前价相当于区间最高价的 45.0\%，为区间最低价的 1.24 倍。 月度收益的年化波动率 28.2\%，低于全样本中位数 37.4\%，在三级行业“其他农产品加工”的 11 家公司中波动率排第 \zhnumber{11} 位。

\begin{center}
\begin{tikzpicture}
\begin{axis}[
  width=12.6cm, height=3.9cm, scale only axis,
  xmin=2020.922, xmax=2026.888, ymin=3.95, ymax=12.50,
  xtick={2021,2022,2023,2024,2025,2026},
  yearticks,
  yticklabel style={font=\scriptsize},
  ylabel={元/股}, ylabel style={font=\scriptsize},
  grid=major, grid style={LightGray, line width=0.3pt},
  axis line style={InkGray, line width=0.5pt},
  every axis plot/.append style={line width=0.9pt},
]
\addplot[AgriGreen] table[x=x,y=y]{data/clean/profiles/px_000930.dat};
\addplot[SoftRed, dashed, line width=0.7pt] table[x=x,y=y]{data/clean/profiles/pxzz_000930.dat};
\addplot[only marks, mark=*, mark size=1.1pt, SoftRed] table[x=x,y=y]{data/clean/profiles/pxzz_000930.dat};
\end{axis}
\end{tikzpicture}

\captionof{figure}[中粮科技（000930）股价与周期骨架]{中粮科技（000930）月度收盘价（元/股）与 ZigZag 周期骨架（阈值 25\%，圆点为周期转折点）}
\end{center}

\pfl{过去周期分析} 以 25\% 的 ZigZag 阈值划分，样本区间内股价形成 2 个主升段与 3 个主跌段。 最大主跌段为 2021 年 11 月 至 2024 年 8 月，33 个月-58.3\%；最大主升段出现在 2024 年 8 月 至 2026 年 2 月，18 个月+39.0\%。 最近一次转折出现在 2026 年 6 月（下跌段自 2026 年 2 月起，4 个月-37.2\%），当前处于该段的延续之中。 公司股价与申万农林牧渔指数几乎同步见底，个股并未走出独立行情。 该子行业横跨糖、番茄制品、添加剂等多个品种，各自的供给周期与政策（进口配额、关税）是主要变量，公司 2026 年 6 月末资产负债率 40.7\%、半年 ROE -1.6\%（未年化），在本轮下行中属于随行业同步调整的一类。

\pfl{盈利情况} 2026 年上半年营业收入 81.94 亿元（同比 -7.0\%），归母净利润 -1.69 亿元，由上年同期的盈利转为亏损。 以年报为趋势对照，2023---2025 年营业收入由 203.79 亿元变为 175.12 亿元（两年年均复合增速 -7.3\%），归母净利润由 -6.02 亿元变为 0.38 亿元。 按半年报口径，2026 年上半年的 ROE 为 -1.6\%（未年化）、销售净利率 -2.1\%、毛利率 5.0\%，ROE 在“其他农产品加工”11 家公司中排第 \zhnumber{9} 位。 同期全样本半年 ROE 中位数 0.75\%，公司与之相差 2.37 个百分点（弱于中位数公司）。 结合公司披露，2026 年上半年归母净利润 -1.69 亿元，同比由盈转亏（下降 257.85\%）；扣非净利润 -1.……；2025 年营业收入 175.12 亿元，同比下降 12.67\%；归母净利润仅 3,836.76 万元（……。\footnote{\texttt{https://vip.stock.finance.sina.com.cn/corp/view/vCB\_AllBulletinDetail.php}}\footnote{\texttt{http://static.cninfo.com.cn/finalpage/2026-04-24/1225174161.PDF}}

\begin{center}\small\setlength{\tabcolsep}{4pt}
\captionof{table}[中粮科技（000930）关键财务指标]{中粮科技（000930）关键财务指标（合并报表口径；两个半年报期均未年化；现金短债比 $=$ 货币资金 $\div$ 短期债务）}
\begin{tabular}{@{}lrrrrr@{}}
\toprule
指标 & 2023 年 & 2024 年 & 2025 年 & 2025 年半年报 & 2026 年半年报 \\
\midrule
营业收入（亿元） & 203.79 & 200.53 & 175.12 & 88.12 & 81.94 \\
归母净利润（亿元） & -6.02 & 0.25 & 0.38 & 1.07 & -1.69 \\
毛利率（\%） & 5.5 & 7.8 & 8.4 & 9.0 & 5.0 \\
ROE（\%） & -5.3 & 0.2 & 0.4 & 1.0 & -1.6 \\
资产负债率（\%） & 34.5 & 37.0 & 34.3 & 37.6 & 40.7 \\
经营活动现金流净额（亿元） & 19.58 & 1.56 & 18.47 & 7.92 & -8.75 \\
货币资金（亿元） & 14.81 & 19.09 & 24.88 & 23.88 & 21.29 \\
有息负债合计（亿元） & 36.43 & 44.21 & 24.33 & 44.90 & 49.14 \\
现金短债比（倍） & 0.45 & 0.47 & 1.30 & 0.60 & 0.62 \\
\bottomrule
\end{tabular}
\end{center}

\begin{center}
\begin{tikzpicture}
\begin{groupplot}[
  group style={group size=2 by 1, horizontal sep=0.95cm},
  width=7.2cm, height=4.6cm,
  tick label style={font=\tiny},
  title style={font=\scriptsize\heiti},
  yticklabel style={font=\tiny},
  grid=major, grid style={LightGray, line width=0.3pt},
  axis line style={InkGray, line width=0.5pt},
  enlarge x limits={abs=0.8},
  legend style={font=\scriptsize, at={(0.5,-0.22)}, anchor=north,
                legend columns=2, draw=none, fill=none, column sep=6pt},
]
\nextgroupplot[
  ybar, bar width=4pt,
  ymin=-30.18, ymax=264.59,
  xtick={0,1,2,3,4,5.7},
  xticklabels={2021,2022,2023,2024,2025,2026H1},
  ylabel={亿元},
]
\addplot[fill=AgriGreen!75, draw=AgriGreen, bar shift=-2.6pt] table[x=i, y=rev]{data/clean/profiles/fin_000930.dat};
\addplot[fill=SoftRed!60, draw=SoftRed, bar shift=2.6pt] table[x=i, y=ni]{data/clean/profiles/fin_000930.dat};
\addplot[fill=AgriGreen!30, draw=AgriGreen, pattern=north east lines, bar shift=-2.6pt] table[x=x, y=rev]{data/clean/profiles/finh_000930.dat};
\addplot[fill=SoftRed!20, draw=SoftRed, pattern=north east lines, bar shift=2.6pt] table[x=x, y=ni]{data/clean/profiles/finh_000930.dat};
\legend{营业收入（年/半年）, 归母净利润（年/半年）}
\nextgroupplot[
  ymin=-9.0, ymax=45.3,
  xtick={0,1,2,3,4,5.7},
  xticklabels={2021,2022,2023,2024,2025,2026H1},
  ylabel={\%},
]
\addplot[AgriGold, mark=*, mark size=1.3pt] table[x=i, y=roe]{data/clean/profiles/fin_000930.dat};
\addplot[OilBlue, mark=square*, mark size=1.3pt] table[x=i, y=debt]{data/clean/profiles/fin_000930.dat};
\addplot[only marks, AgriGold, mark=*, mark size=2.0pt] table[x=x, y=roe]{data/clean/profiles/finh_000930.dat};
\addplot[only marks, OilBlue, mark=square*, mark size=2.0pt] table[x=x, y=debt]{data/clean/profiles/finh_000930.dat};
\legend{ROE, 资产负债率}
\end{groupplot}
\end{tikzpicture}
\begin{tikzpicture}
\begin{axis}[
  width=15.2cm, height=3.1cm, scale only axis,
  ymin=-14.54, ymax=56.09,
  xtick={0,1,2,3,4,5.7},
  xticklabels={2021,2022,2023,2024,2025,2026H1},
  tick label style={font=\tiny},
  yticklabel style={font=\tiny},
  ylabel={亿元}, ylabel style={font=\scriptsize},
  ybar, bar width=4pt,
  grid=major, grid style={LightGray, line width=0.3pt},
  axis line style={InkGray, line width=0.5pt},
  enlarge x limits={abs=0.8},
  legend style={font=\scriptsize, at={(0.5,-0.30)}, anchor=north,
                legend columns=3, draw=none, fill=none, column sep=10pt},
]
\addplot[fill=AgriGreen!55, draw=AgriGreen, bar shift=-3.0pt] table[x=i, y=cash]{data/clean/profiles/fin_000930.dat};
\addplot[fill=SoftRed!45, draw=SoftRed, bar shift=3.0pt] table[x=i, y=ibd]{data/clean/profiles/fin_000930.dat};
\addplot[fill=AgriGreen!25, draw=AgriGreen, pattern=north east lines, bar shift=-3.0pt] table[x=x, y=cash]{data/clean/profiles/finh_000930.dat};
\addplot[fill=SoftRed!20, draw=SoftRed, pattern=north east lines, bar shift=3.0pt] table[x=x, y=ibd]{data/clean/profiles/finh_000930.dat};
\addplot[OilBlue, mark=*, mark size=1.1pt, line width=0.8pt] table[x=i, y=ocf]{data/clean/profiles/fin_000930.dat};
\addplot[only marks, OilBlue, mark=*, mark size=1.8pt] table[x=x, y=ocf]{data/clean/profiles/finh_000930.dat};
\legend{货币资金, 有息负债, 经营活动现金流净额}
\end{axis}
\end{tikzpicture}

\captionof{figure}[中粮科技（000930）近年主要财务指标]{中粮科技（000930）近年主要财务指标（左上：营业收入与归母净利润；右上：ROE 与资产负债率；下：货币资金、有息负债与经营活动现金流净额。
2021---2025 年为年报口径，2026H1 为 2026 年半年报/半年末，与其前一年报之间留出间隔、以斜纹或空心点区分；半年报数据未年化）}
\end{center}

\pfl{财务分析} 2026 年 6 月末资产负债率 40.7\%，较 2025 年末 上升 6.4 个百分点（样本中位数 52.2\%，所处三级行业内排第 \zhnumber{7} 位）。 2026 年 6 月末货币资金 21.29 亿元、短期债务（短期借款与一年内到期非流动负债）34.61 亿元、长期债务 14.53 亿元，有息负债合计 49.14 亿元，现金短债比 0.62 倍——覆盖偏紧。 净有息负债 27.85 亿元，相当于其 2026 年上半年营业收入的 34\%。 流动比率 1.27、速动比率 0.86，短期指标尚可。 2026 年上半年经营活动现金流净额 -8.75 亿元，经营现金流与利润同时为负，现金消耗较快。 2026 年 6 月末在建工程 7.92 亿元，较上年末增加 0.40 亿元，仍有在建投入。 2026 年 6 月末商誉 3.66 亿元，占归母净资产的 3.4\%，是净资产中最脆弱的一块。

\pfl{重大战略转型} 近三年公司的战略动作集中在以下几个方面：2018 中粮集团将境内玉米深加工与燃料乙醇资产整体注入上市公司，形成生物能源、淀粉、淀粉糖、饲料原料等多元产品结构；2025 贯彻“以客户为中心”理念，柔性生产无水乙醇、食用酒精、药用乙醇等产品；酒精业务可使用玉米、不宜食用的水稻和小麦、……；2025 广西公司热电系统改造升级项目推进并网（2025 年年报披露）。 公告口径上，近三年（2023 年 9 月---2026 年 9 月）公司披露重大重组与并购类公告 0 条、对外投资与转型类公告 1 条、回购与股权激励类公告 8 条，合计公告 336 条。 第一大行业仍是“农产品加工及销售”，收入占比 98.3\%，与 2021 年年报相比没有方向性变化。 综合判断，公司在报告期内有明确的外延扩张或业务结构调整动作，转型方向与融资安排之间存在直接关联。\footnote{\texttt{https://pdf.dfcfw.com/pdf/H3\_AP202111261531119434\_1.pdf}}\footnote{\texttt{https://stock.stockstar.com/notice/SN2026042400019909.shtml}}

\begin{center}\small\setlength{\tabcolsep}{4pt}
\captionof{table}[中粮科技（000930）历史融资与资本运作]{中粮科技（000930）历史融资与资本运作（据公司公告与公开报道整理；金额为披露值，含首次公开发行、再融资、债券、银行借款与重整投资等）}
\begin{tabular}{@{}p{1.5cm}p{3.2cm}p{2.4cm}p{6.3cm}@{}}
\toprule
时间 & 方式 & 规模 & 用途与说明 \\
\midrule
1999-07-12 & IPO（深圳证券交易所，原安徽丰原生物化学股份有限公司） & 实际发行量 6,000 万股，发行价 6.50 元/股，实际募集资金合计 39,000.00 万元，发行费用总额 1,440.00 万元，募集资金净额 37,560.00 万元 & 上市地深圳证券交易所，主承销商安徽省信托投资公司，发行方式网上定价发行…… \\
2003-04-24 & 可转换公司债券（“丰原转债”） & 5 亿元（500,000,000 元） & 经中国证监会证监公司字〔2003〕36 号核准…… \\
2018 & 发行股份购买资产（重大资产重组暨关联交易） & 标的资产评估值合计 828,472.80 万元；新增股份 883,223,262 股 & 收购生化能源、生物化学、桦力投资 100\% 股权，注入中粮集团玉米深加工与燃料乙醇资产…… \\
2022 & 现金分红 & 以 1,864,720,561 股为基数每 10 股派发现金红利 5.70 元（含税） & 股东回报；该年度公司盈利水平较高，分红力度显著大于后续年度 \\
\bottomrule
\end{tabular}
\end{center}

\pfl{历史融投资情况} 从现金流量表看，2025 年公司取得借款收到的现金 69.69 亿元、偿还债务支付的现金 76.63 亿元、吸收权益性投资 — 亿元，筹资活动现金流净额 -7.88 亿元（2023 年为取得借款 75.65 亿元、偿还债务 72.42 亿元）。 2026 年上半年筹资活动现金流净额为 14.03 亿元（取得借款 39.63 亿元、偿还债务 25.42 亿元，未年化），仍处于净融入状态。 2025 年分配股利、利润或偿付利息支付的现金合计 0.86 亿元。 公开信息中与公司直接相关的融资与资本运作包括：IPO（深圳证券交易所，原安徽丰原生物化学股份有限公司）、发行股份购买资产（重大资产重组暨关联交易）、可转换公司债券（“丰原转债”）、现金分红。 其中规模较大的两笔为：2003-04-24可转换公司债券（“丰原转债”）（5 亿元（500,000,000 元））；1999-07-12IPO（深圳证券交易所，原安徽丰原生物化学股份有限公司）（实际发行量 6,000 万股，发行价 6.50 元/股，实际……）。 股本方面，近三年披露股本变动 6 次（股份回购 2 次、激励股份解禁 1 次）；总股本由 2021 年末的 18.66 亿股变为 18.60 亿股（-0.32\%）。 分红方面，近三年现金分红 2 次、累计每股派现 0.60 元（最近一次 2025-12-31）——在利润承压的阶段仍维持了分红。\footnote{\texttt{https://vip.stock.finance.sina.com.cn/corp/go.php/vISSUE\_NewStock/stockid/000930.phtml}}\footnote{\texttt{http://money.finance.sina.com.cn/corp/view/vCB\_AllBulletinDetail.php}}\footnote{\texttt{http://vip.stock.finance.sina.com.cn/corp/view/vCB\_AllBulletinDetail.php}}

\pfl{未来潜在融资需求分析} 2026 年 6 月末货币资金 21.29 亿元，而短期债务 34.61 亿元（现金短债比 0.62 倍），2026 年上半年经营现金流净额 -8.75 亿元。按“短期债务减去账面货币资金”的滚动偿付口径，静态缺口约 13.32 亿元；若再计入上半年亏损的年化额（3.39 亿元），维持一年正常经营所需的外部资金约 13.32---16.70 亿元。 公司 2026 年上半年归母净利润为负（-1.69 亿元），但 6 月末资产负债率 40.7\% 尚未进入第一档（亏损且负债率 $>$65\%）的区间，当前融资需求以补充流动资金、支撑低谷期经营性支出为主。 公开披露的下一步融资线索包括：2019—2023 公司历史上曾实施限制性股票激励计划：2022 年 2 月披露首次解锁安排，符合解锁条件的激励对象 412 人，实际可解锁……；2026-08 2026 年半年报披露主业毛利压缩、净利转亏、现金流恶化，未披露新的股权融资或发债预案。 资金需求还要叠加在建工程：期末在建工程 7.92 亿元、2026 年上半年购建固定资产等支出 5.45 亿元（未年化），这部分支出在周期底部通常会被主动放缓。\footnote{\texttt{http://money.finance.sina.com.cn/corp/view/vCB\_AllBulletinDetail.php}}\footnote{\texttt{https://finance.sina.com.cn/stock/relnews/cn/2026-08-28/doc-inipwhrt5706332.shtml}}


\subsubsection{冠农股份（600251）}
\label{co:600251}

\pfl{公司基本情况} 冠农股份成立于 2003 年，新疆生产建设兵团第二师国资背景的农产品精深加工企业，产业涉及棉花加工、番茄制品、甜菜制糖、棉籽蛋白与生物发酵饲料；核心利润来自对国投罗钾（20.3\%）和开都河水电（25.28\%）的权益法投资收益，主业加工板块盈利能力较弱。 公司于 2003-06-09 在上交所首发上市，注册资本 7.77 亿元，总股本 7.77 亿股。 截至 2026-09-24收盘价 10.03 元，流通市值 77.9 亿元，近一年+21.3\%，流通市值在三级行业“果蔬加工”的 5 家公司中排第 \zhnumber{2} 位，近一年涨跌幅高于全样本中位数（-10.1\%）31.4 个百分点。 按 2026 年半年报披露的产品构成，收入占比前三位为油脂产品 38.7\%；棉花产品 34.4\%；果蔬产品 19.7\%。 与 2021 年年报相比，第一大产品由“皮棉”变为“油脂产品”，主业重心已经迁移。 发展过程中的关键节点包括：2017 年 12 月新疆生产建设兵团第二师国资委将其持有新疆冠源投资有限责任公司（现已……（2017）；2019 年 10 月 25 日新疆证监局下发《关于对新疆冠农果茸股份有限公司的监管关注……（2019）；2020 年 1 月 14 日上海证券交易所对公司及时任财务总监邱照亮予以监管关注（上证……（2020）。\footnote{\texttt{http://www.sse.com.cn/disclosure/listedinfo/announcement/c/2017-12-20/600251\_20171220\_3.}}\footnote{\texttt{http://www.sse.com.cn/disclosure/listedinfo/announcement/c/new/2021-06-22/600251\_2021062}}

\pfl{历史股价} 以 2021 年 1 月为起点、2026 年 9 月为终点，股价由 5.80 元变为 10.03 元，累计 +72.9\%，区间最高 12.95 元（2026 年 4 月）、最低 5.80 元（2021 年 1 月）；当前价相当于区间最高价的 77.5\%，为区间最低价的 1.73 倍。 月度收益的年化波动率 35.0\%，低于全样本中位数 37.4\%，在三级行业“果蔬加工”的 5 家公司中波动率排第 \zhnumber{5} 位。

\begin{center}
\begin{tikzpicture}
\begin{axis}[
  width=12.6cm, height=3.9cm, scale only axis,
  xmin=2020.922, xmax=2026.888, ymin=5.39, ymax=13.86,
  xtick={2021,2022,2023,2024,2025,2026},
  yearticks,
  yticklabel style={font=\scriptsize},
  ylabel={元/股}, ylabel style={font=\scriptsize},
  grid=major, grid style={LightGray, line width=0.3pt},
  axis line style={InkGray, line width=0.5pt},
  every axis plot/.append style={line width=0.9pt},
]
\addplot[AgriGreen] table[x=x,y=y]{data/clean/profiles/px_600251.dat};
\addplot[SoftRed, dashed, line width=0.7pt] table[x=x,y=y]{data/clean/profiles/pxzz_600251.dat};
\addplot[only marks, mark=*, mark size=1.1pt, SoftRed] table[x=x,y=y]{data/clean/profiles/pxzz_600251.dat};
\end{axis}
\end{tikzpicture}

\captionof{figure}[冠农股份（600251）股价与周期骨架]{冠农股份（600251）月度收盘价（元/股）与 ZigZag 周期骨架（阈值 25\%，圆点为周期转折点）}
\end{center}

\pfl{过去周期分析} 以 25\% 的 ZigZag 阈值划分，样本区间内股价形成 3 个主升段与 3 个主跌段。 最大主升段出现在 2024 年 8 月 至 2026 年 4 月，20 个月+93.0\%；最大主跌段为 2022 年 7 月 至 2024 年 8 月，25 个月-41.1\%。 最近一次转折出现在 2026 年 6 月（下跌段自 2026 年 4 月起，2 个月-33.3\%），当前处于该段的延续之中。 公司股价较申万农林牧渔指数提前约 64 个月见底，是板块中较早定价周期反转的公司之一。 果蔬加工的外向型特征明显，出口订单与原料产地价格决定利润率，公司 2026 年上半年实现盈利、半年 ROE 13.3\%（未年化），好于全样本中位数（0.75\%），下行周期对其报表的冲击小于同业。

\pfl{盈利情况} 2026 年上半年营业收入 13.89 亿元（同比 +0.3\%），归母净利润 5.37 亿元，同比 +79.1\%。 以年报为趋势对照，2023---2025 年营业收入由 43.48 亿元变为 25.15 亿元（两年年均复合增速 -23.9\%），归母净利润由 7.20 亿元变为 3.48 亿元。 按半年报口径，2026 年上半年的 ROE 为 13.3\%（未年化）、销售净利率 39.5\%、毛利率 8.6\%，ROE 在“果蔬加工”5 家公司中排第 \zhnumber{1} 位。 行业同口径半年 ROE 的中位数为 0.75\%，该公司高于中位数 12.59 个百分点。 结合公司披露，2025 年年报披露商誉减值情况：子公司 1、子公司 4 的商誉已全额计提减值准备，子公司 2 商誉金额……；2025 年营业收入 25.15 亿元，同比下降 28.40\%（2024 年调整后 35.13 亿元，调……。\footnote{\texttt{https://money.finance.sina.com.cn/corp/view/vCB\_AllBulletinDetail.php}}

\begin{center}\small\setlength{\tabcolsep}{4pt}
\captionof{table}[冠农股份（600251）关键财务指标]{冠农股份（600251）关键财务指标（合并报表口径；两个半年报期均未年化；现金短债比 $=$ 货币资金 $\div$ 短期债务）}
\begin{tabular}{@{}lrrrrr@{}}
\toprule
指标 & 2023 年 & 2024 年 & 2025 年 & 2025 年半年报 & 2026 年半年报 \\
\midrule
营业收入（亿元） & 43.48 & 35.13 & 25.15 & 13.86 & 13.89 \\
归母净利润（亿元） & 7.20 & 2.05 & 3.48 & 3.00 & 5.37 \\
毛利率（\%） & 16.8 & 8.1 & 8.8 & 7.0 & 8.6 \\
ROE（\%） & 20.4 & 5.5 & 9.4 & 8.0 & 13.3 \\
资产负债率（\%） & 48.4 & 46.3 & 46.9 & 34.2 & 30.8 \\
经营活动现金流净额（亿元） & 2.05 & 6.13 & 1.26 & 8.36 & 5.12 \\
货币资金（亿元） & 18.77 & 16.92 & 10.80 & 18.78 & 13.50 \\
有息负债合计（亿元） & 26.81 & 23.42 & 22.77 & 14.86 & 11.87 \\
现金短债比（倍） & 0.76 & 0.78 & 0.51 & 1.43 & 1.19 \\
\bottomrule
\end{tabular}
\end{center}

\begin{center}
\begin{tikzpicture}
\begin{groupplot}[
  group style={group size=2 by 1, horizontal sep=0.95cm},
  width=7.2cm, height=4.6cm,
  tick label style={font=\tiny},
  title style={font=\scriptsize\heiti},
  yticklabel style={font=\tiny},
  grid=major, grid style={LightGray, line width=0.3pt},
  axis line style={InkGray, line width=0.5pt},
  enlarge x limits={abs=0.8},
  legend style={font=\scriptsize, at={(0.5,-0.22)}, anchor=north,
                legend columns=2, draw=none, fill=none, column sep=6pt},
]
\nextgroupplot[
  ybar, bar width=4pt,
  ymin=-4.35, ymax=48.69,
  xtick={0,1,2,3,4,5.7},
  xticklabels={2021,2022,2023,2024,2025,2026H1},
  ylabel={亿元},
]
\addplot[fill=AgriGreen!75, draw=AgriGreen, bar shift=-2.6pt] table[x=i, y=rev]{data/clean/profiles/fin_600251.dat};
\addplot[fill=SoftRed!60, draw=SoftRed, bar shift=2.6pt] table[x=i, y=ni]{data/clean/profiles/fin_600251.dat};
\addplot[fill=AgriGreen!30, draw=AgriGreen, pattern=north east lines, bar shift=-2.6pt] table[x=x, y=rev]{data/clean/profiles/finh_600251.dat};
\addplot[fill=SoftRed!20, draw=SoftRed, pattern=north east lines, bar shift=2.6pt] table[x=x, y=ni]{data/clean/profiles/finh_600251.dat};
\legend{营业收入（年/半年）, 归母净利润（年/半年）}
\nextgroupplot[
  ymin=-4.5, ymax=61.4,
  xtick={0,1,2,3,4,5.7},
  xticklabels={2021,2022,2023,2024,2025,2026H1},
  ylabel={\%},
]
\addplot[AgriGold, mark=*, mark size=1.3pt] table[x=i, y=roe]{data/clean/profiles/fin_600251.dat};
\addplot[OilBlue, mark=square*, mark size=1.3pt] table[x=i, y=debt]{data/clean/profiles/fin_600251.dat};
\addplot[only marks, AgriGold, mark=*, mark size=2.0pt] table[x=x, y=roe]{data/clean/profiles/finh_600251.dat};
\addplot[only marks, OilBlue, mark=square*, mark size=2.0pt] table[x=x, y=debt]{data/clean/profiles/finh_600251.dat};
\legend{ROE, 资产负债率}
\end{groupplot}
\end{tikzpicture}
\begin{tikzpicture}
\begin{axis}[
  width=15.2cm, height=3.1cm, scale only axis,
  ymin=-3.26, ymax=36.47,
  xtick={0,1,2,3,4,5.7},
  xticklabels={2021,2022,2023,2024,2025,2026H1},
  tick label style={font=\tiny},
  yticklabel style={font=\tiny},
  ylabel={亿元}, ylabel style={font=\scriptsize},
  ybar, bar width=4pt,
  grid=major, grid style={LightGray, line width=0.3pt},
  axis line style={InkGray, line width=0.5pt},
  enlarge x limits={abs=0.8},
  legend style={font=\scriptsize, at={(0.5,-0.30)}, anchor=north,
                legend columns=3, draw=none, fill=none, column sep=10pt},
]
\addplot[fill=AgriGreen!55, draw=AgriGreen, bar shift=-3.0pt] table[x=i, y=cash]{data/clean/profiles/fin_600251.dat};
\addplot[fill=SoftRed!45, draw=SoftRed, bar shift=3.0pt] table[x=i, y=ibd]{data/clean/profiles/fin_600251.dat};
\addplot[fill=AgriGreen!25, draw=AgriGreen, pattern=north east lines, bar shift=-3.0pt] table[x=x, y=cash]{data/clean/profiles/finh_600251.dat};
\addplot[fill=SoftRed!20, draw=SoftRed, pattern=north east lines, bar shift=3.0pt] table[x=x, y=ibd]{data/clean/profiles/finh_600251.dat};
\addplot[OilBlue, mark=*, mark size=1.1pt, line width=0.8pt] table[x=i, y=ocf]{data/clean/profiles/fin_600251.dat};
\addplot[only marks, OilBlue, mark=*, mark size=1.8pt] table[x=x, y=ocf]{data/clean/profiles/finh_600251.dat};
\legend{货币资金, 有息负债, 经营活动现金流净额}
\end{axis}
\end{tikzpicture}

\captionof{figure}[冠农股份（600251）近年主要财务指标]{冠农股份（600251）近年主要财务指标（左上：营业收入与归母净利润；右上：ROE 与资产负债率；下：货币资金、有息负债与经营活动现金流净额。
2021---2025 年为年报口径，2026H1 为 2026 年半年报/半年末，与其前一年报之间留出间隔、以斜纹或空心点区分；半年报数据未年化）}
\end{center}

\pfl{财务分析} 2026 年 6 月末资产负债率 30.8\%，较 2025 年末 下降 16.1 个百分点（样本中位数 52.2\%，所处三级行业内排第 \zhnumber{3} 位）。 2026 年 6 月末货币资金 13.50 亿元、短期债务（短期借款与一年内到期非流动负债）11.36 亿元、长期债务 0.50 亿元，有息负债合计 11.87 亿元，现金短债比 1.19 倍——覆盖充分。 净有息负债 -1.63 亿元，相当于其 2026 年上半年营业收入的 -12\%。 流动比率 1.92、速动比率 1.54，短期指标尚可。 2026 年上半年经营活动现金流净额 5.12 亿元，经营现金流与利润同向为正，内生资金可以支撑运营。

\pfl{重大战略转型} 公司在报告期内的战略动作可归为 4 项：2021 筹划公开发行可转债以支持主业扩产，后于 2023 年 3 月终止，转向依靠参股公司分红与自有资金支持项目投资；2023 参股国投罗钾（控股股东为国投矿业投资有限公司）增资扩产，公司按持股比例出资 7,202.97 万元；2024 公司承诺在 36 个月内继续推进绿洲农业将新疆新建番茄制品有限公司 100\% 股权或控股权对外出让，以解决同业竞……；2026-06 公司披露对外投资严格遵循长期持有、战略协同、稳健增值原则，核心资产为对国投罗钾、开都河水电的战略股权投资。 公告口径上，近三年（2023 年 9 月---2026 年 9 月）公司披露重大重组与并购类公告 1 条、对外投资与转型类公告 0 条、回购与股权激励类公告 0 条，合计公告 346 条。 主营结构的迁移已经落到报表上：2021 年年报的第一大产品为“皮棉”（64.8\%），2026 年半年报变为“油脂产品”（38.7\%）。 综合判断，公司在报告期内有明确的外延扩张或业务结构调整动作，转型方向与融资安排之间存在直接关联。\footnote{\texttt{http://www.sse.com.cn/disclosure/listedinfo/announcement/c/new/2023-03-22/600251\_2023032}}\footnote{\texttt{https://pdf.dfcfw.com/pdf/H2\_AN202312011613169716\_1.pdf}}\footnote{\texttt{https://www.sse.com.cn/disclosure/listedinfo/announcement/c/new/2025-04-03/600251\_202504}}

\begin{center}\small\setlength{\tabcolsep}{4pt}
\captionof{table}[冠农股份（600251）历史融资与资本运作]{冠农股份（600251）历史融资与资本运作（据公司公告与公开报道整理；金额为披露值，含首次公开发行、再融资、债券、银行借款与重整投资等）}
\begin{tabular}{@{}p{1.5cm}p{3.2cm}p{2.4cm}p{6.3cm}@{}}
\toprule
时间 & 方式 & 规模 & 用途与说明 \\
\midrule
2003-05 & IPO（上海证券交易所） & 发行 4,000 万股，发行价 5.75 元/股 & 经中国证监会证监发行字[2003]46 号核准，采用向二级市场投资者定价配售方式…… \\
2017-12 & 国有股权无偿划转（间接控股股东变更） & 冠源投资 100\% 国有股权无偿划转 & 兵团国资体系内部整合；第二师国资委将所持冠源投资 100\% 股权无偿划转至绿原国资…… \\
2021—2023 & 拟公开发行可转换公司债券（最终终止） & 拟募集资金总额 110,000.00 万元 & 公司公开发行可转换公司债券募投项目（公告未在本底稿核实的页面中列明具体项目）…… \\
2023-12 & 参股公司增资（对外投资） & 投资金额 7,202.97 万元 & 对参股公司国投新疆罗布泊钾盐有限责任公司增资；不构成重大资产重组，不构成关联交易…… \\
2025 & 参股公司现金分红（投资收益回收） & 收到国投罗钾 2024 年度现金分红款 41,222.05 万元 & 为公司生产经营和项目投资提供资金保障；公司持有国投罗钾 20.3\% 股权…… \\
\bottomrule
\end{tabular}
\end{center}

\pfl{历史融投资情况} 从现金流量表看，2025 年公司取得借款收到的现金 29.56 亿元、偿还债务支付的现金 30.76 亿元、吸收权益性投资 — 亿元，筹资活动现金流净额 -3.87 亿元（2023 年为取得借款 31.99 亿元、偿还债务 29.70 亿元）。 2026 年上半年筹资活动现金流净额为 -11.60 亿元（取得借款 6.47 亿元、偿还债务 17.81 亿元，未年化），已转为净流出，处于净还债状态。 2025 年分配股利、利润或偿付利息支付的现金合计 2.42 亿元。 公司披露的股本与债务融资脉络包括：IPO（上海证券交易所）、参股公司增资（对外投资）、参股公司现金分红（投资收益回收）、国有股权无偿划转（间接控股股东变更）、拟公开发行可转换公司债券（最终终止）。 其中规模较大的两笔为：2021—2023拟公开发行可转换公司债券（最终终止）（拟募集资金总额 110,000.00 万元）；2025参股公司现金分红（投资收益回收）（收到国投罗钾 2024 年度现金分红款 41,222.05 ……）。 股本方面，近三年披露股本变动 5 次（股份回购 2 次）；总股本由 2021 年末的 7.81 亿股变为 7.77 亿股（-0.57\%）。 分红方面，近三年现金分红 5 次、累计每股派现 0.88 元（最近一次 2025-12-31）——在利润承压的阶段仍维持了分红。\footnote{\texttt{http://money.finance.sina.com.cn/corp/view/vISSUE\_MarketBulletinDetail.php}}\footnote{\texttt{http://www.sse.com.cn/disclosure/listedinfo/announcement/c/2017-12-20/600251\_20171220\_3.}}\footnote{\texttt{http://www.sse.com.cn/disclosure/listedinfo/announcement/c/new/2023-03-22/600251\_2023032}}

\pfl{未来潜在融资需求分析} 2026 年 6 月末货币资金 13.50 亿元，而短期债务 11.36 亿元（现金短债比 1.19 倍），2026 年上半年经营现金流净额 5.12 亿元。账面货币资金可以覆盖短期债务，缺口不体现在静态偿付层面。 按第 \ref{sec:financing-need} 节的三档划分，公司属于第三档“扩张型”：2026 年上半年 ROE 13.3\%（未年化，折合约 26.7\%）、6 月末资产负债率 30.8\%，资产负债表仍具备加杠杆空间；融资用途以产能扩张、品类并购与海外布局为主，单笔规模通常在 5---20 亿元量级，会被市场定价为成长而非补血。 公开披露的下一步融资线索包括：2023-03 终止向不特定对象发行可转换公司债券（原拟募集资金 110,000.00 万元）；2025-06 公司收到国投罗钾 2024 年度现金分红款 41,222.05 万元，公告称“为公司生产经营和项目投资提供了更加充足的资……；2026 2026 年半年度报告期后事项包括 2026-06-18 披露 2025 年年度权益分派实施公告、2026-06-27 ……。\footnote{\texttt{http://www.sse.com.cn/disclosure/listedinfo/announcement/c/new/2023-03-22/600251\_2023032}}\footnote{\texttt{http://vip.stock.finance.sina.com.cn/corp/view/vCB\_AllBulletinDetail.php}}\footnote{\texttt{https://money.finance.sina.com.cn/corp/go.php/vCB\_AllBulletin/stockid/600251.phtml}}


\subsubsection{国投中鲁（600962）}
\label{co:600962}

\pfl{公司基本情况} 国投中鲁成立于 2001 年，国家开发投资集团有限公司控股的浓缩苹果汁加工企业，是国内主板第一家上市的浓缩果汁企业，主业规模与盈利能力有限、母公司长期存在未弥补亏损；2025 年底起筹划以发行股份方式收购中国电子工程设计院 100\% 股份，跨界转型电子信息产业。 公司于 2004-06-22 在上交所首发上市，注册资本 2.62 亿元，总股本 2.62 亿股。 截至 2026-09-24收盘价 26.15 元，流通市值 68.6 亿元，近一年+23.0\%，流通市值在三级行业“果蔬加工”的 5 家公司中排第 \zhnumber{3} 位，近一年涨跌幅高于全样本中位数（-10.1\%）33.1 个百分点。 按 2025 年年报披露的行业构成，收入占比前三位为饮料制造业 99.5\%；其他(补充) 0.5\%。 发展过程中的关键节点包括：经中国证监会证监发行字〔2004〕72 号文核准，2004 年 6 月 7 日采取上网定……（2004）；成为浓缩苹果汁行业第一家在国内主板上市的国有控股企业（2004）；2025 年 7 月起市场关注公司筹划以发行股份方式收购中国电子工程设计院股份有限公司（……（2025）。\footnote{\texttt{http://www.ce.cn/ztpd/tszt/gushi/2005/ssnb/PDF/200503/20/P020050320462494982913.pdf}}\footnote{\texttt{https://zhonglu.sdic.com.cn/gtzl/about/A2118index\_1.htm}}\footnote{\texttt{https://www.nbd.com.cn/articles/2025-12-30/4202501.html}}

\pfl{历史股价} 以 2021 年 1 月为起点、2026 年 9 月为终点，股价由 7.05 元变为 26.15 元，累计 +270.9\%，区间最高 30.17 元（2026 年 6 月）、最低 7.04 元（2021 年 2 月）；当前价相当于区间最高价的 86.7\%，为区间最低价的 3.71 倍。 月度收益的年化波动率 45.2\%，高于全样本中位数 37.4\%，在三级行业“果蔬加工”的 5 家公司中波动率排第 \zhnumber{3} 位。 从事件看，2025 年 3 月 至 2026 年 6 月的上涨（+137.9\%）与公司公司收到上海证券交易所《关于国投中鲁果汁股份有限公司发行股份购买资产并募集配……（2026-03-10）在时间上重合——股价的拐点多数能在公司自身的资本运作与风险事件上找到对应。\footnote{\texttt{http://static.sse.com.cn/stock/disclosure/announcement/c/202607/600962\_20260724\_ETJK.pdf}}

\begin{center}
\begin{tikzpicture}
\begin{axis}[
  width=12.6cm, height=3.9cm, scale only axis,
  xmin=2020.922, xmax=2026.888, ymin=6.55, ymax=32.28,
  xtick={2021,2022,2023,2024,2025,2026},
  yearticks,
  yticklabel style={font=\scriptsize},
  ylabel={元/股}, ylabel style={font=\scriptsize},
  grid=major, grid style={LightGray, line width=0.3pt},
  axis line style={InkGray, line width=0.5pt},
  every axis plot/.append style={line width=0.9pt},
]
\addplot[AgriGreen] table[x=x,y=y]{data/clean/profiles/px_600962.dat};
\addplot[SoftRed, dashed, line width=0.7pt] table[x=x,y=y]{data/clean/profiles/pxzz_600962.dat};
\addplot[only marks, mark=*, mark size=1.1pt, SoftRed] table[x=x,y=y]{data/clean/profiles/pxzz_600962.dat};
\end{axis}
\end{tikzpicture}

\captionof{figure}[国投中鲁（600962）股价与周期骨架]{国投中鲁（600962）月度收盘价（元/股）与 ZigZag 周期骨架（阈值 25\%，圆点为周期转折点）}
\end{center}

\pfl{过去周期分析} 以 25\% 的 ZigZag 阈值划分，样本区间内股价形成 3 个主升段与 2 个主跌段。 最大主升段出现在 2025 年 3 月 至 2026 年 6 月，15 个月+137.9\%；最大主跌段为 2023 年 3 月 至 2024 年 2 月，11 个月-37.5\%。 最近一次转折出现在 2026 年 6 月（上涨段自 2025 年 3 月起，15 个月+137.9\%），当前处于该段之后的震荡之中。 公司股价较申万农林牧渔指数提前约 63 个月见底，是板块中较早定价周期反转的公司之一。 果蔬加工的外向型特征明显，出口订单与原料产地价格决定利润率，公司 2026 年 6 月末资产负债率 52.0\%、半年 ROE 0.2\%（未年化），在本轮下行中属于随行业同步调整的一类。

\pfl{盈利情况} 2026 年上半年营业收入 8.83 亿元（同比 -12.4\%），归母净利润 0.02 亿元，同比 -92.7\%。 以年报为趋势对照，2023---2025 年营业收入由 14.87 亿元变为 19.84 亿元（两年年均复合增速 +15.5\%），归母净利润由 0.58 亿元变为 0.41 亿元。 按半年报口径，2026 年上半年的 ROE 为 0.2\%（未年化）、销售净利率 0.7\%、毛利率 14.2\%，ROE 在“果蔬加工”5 家公司中排第 \zhnumber{3} 位。 同期全样本半年 ROE 中位数 0.75\%，公司与之相差 0.56 个百分点（弱于中位数公司）。 结合公司披露，2026 年上半年归母净利润仅 182.12 万元，同比下降 92.74\%；扣非归母净利润 -114.2……；2025 年度归母净利润 4,092.15 万元，但截至报告期末母公司仍存在未弥补亏损，2025 年度拟……。\footnote{\texttt{https://stockmc.xueqiu.com/202608/600962\_20260829\_U1KD.pdf}}\footnote{\texttt{http://money.finance.sina.com.cn/corp/view/vCB\_AllBulletinDetail.php}}

\begin{center}\small\setlength{\tabcolsep}{4pt}
\captionof{table}[国投中鲁（600962）关键财务指标]{国投中鲁（600962）关键财务指标（合并报表口径；两个半年报期均未年化；现金短债比 $=$ 货币资金 $\div$ 短期债务）}
\begin{tabular}{@{}lrrrrr@{}}
\toprule
指标 & 2023 年 & 2024 年 & 2025 年 & 2025 年半年报 & 2026 年半年报 \\
\midrule
营业收入（亿元） & 14.87 & 19.87 & 19.84 & 10.07 & 8.83 \\
归母净利润（亿元） & 0.58 & 0.29 & 0.41 & 0.25 & 0.02 \\
毛利率（\%） & 22.7 & 15.5 & 16.7 & 14.6 & 14.2 \\
ROE（\%） & 7.2 & 3.3 & 4.4 & 2.7 & 0.2 \\
资产负债率（\%） & 59.5 & 65.4 & 61.4 & 50.0 & 52.0 \\
经营活动现金流净额（亿元） & -1.66 & -2.83 & 2.42 & 8.75 & 5.35 \\
货币资金（亿元） & 1.73 & 1.62 & 1.68 & 2.56 & 1.57 \\
有息负债合计（亿元） & 13.11 & 16.70 & 15.18 & 9.29 & 10.30 \\
现金短债比（倍） & 0.14 & 0.10 & 0.11 & 0.29 & 0.17 \\
\bottomrule
\end{tabular}
\end{center}

\begin{center}
\begin{tikzpicture}
\begin{groupplot}[
  group style={group size=2 by 1, horizontal sep=0.95cm},
  width=7.2cm, height=4.6cm,
  tick label style={font=\tiny},
  title style={font=\scriptsize\heiti},
  yticklabel style={font=\tiny},
  grid=major, grid style={LightGray, line width=0.3pt},
  axis line style={InkGray, line width=0.5pt},
  enlarge x limits={abs=0.8},
  legend style={font=\scriptsize, at={(0.5,-0.22)}, anchor=north,
                legend columns=2, draw=none, fill=none, column sep=6pt},
]
\nextgroupplot[
  ybar, bar width=4pt,
  ymin=-2.20, ymax=22.28,
  xtick={0,1,2,3,4,5.7},
  xticklabels={2021,2022,2023,2024,2025,2026H1},
  ylabel={亿元},
]
\addplot[fill=AgriGreen!75, draw=AgriGreen, bar shift=-2.6pt] table[x=i, y=rev]{data/clean/profiles/fin_600962.dat};
\addplot[fill=SoftRed!60, draw=SoftRed, bar shift=2.6pt] table[x=i, y=ni]{data/clean/profiles/fin_600962.dat};
\addplot[fill=AgriGreen!30, draw=AgriGreen, pattern=north east lines, bar shift=-2.6pt] table[x=x, y=rev]{data/clean/profiles/finh_600962.dat};
\addplot[fill=SoftRed!20, draw=SoftRed, pattern=north east lines, bar shift=2.6pt] table[x=x, y=ni]{data/clean/profiles/finh_600962.dat};
\legend{营业收入（年/半年）, 归母净利润（年/半年）}
\nextgroupplot[
  ymin=-8.3, ymax=72.2,
  xtick={0,1,2,3,4,5.7},
  xticklabels={2021,2022,2023,2024,2025,2026H1},
  ylabel={\%},
]
\addplot[AgriGold, mark=*, mark size=1.3pt] table[x=i, y=roe]{data/clean/profiles/fin_600962.dat};
\addplot[OilBlue, mark=square*, mark size=1.3pt] table[x=i, y=debt]{data/clean/profiles/fin_600962.dat};
\addplot[only marks, AgriGold, mark=*, mark size=2.0pt] table[x=x, y=roe]{data/clean/profiles/finh_600962.dat};
\addplot[only marks, OilBlue, mark=square*, mark size=2.0pt] table[x=x, y=debt]{data/clean/profiles/finh_600962.dat};
\legend{ROE, 资产负债率}
\end{groupplot}
\end{tikzpicture}
\begin{tikzpicture}
\begin{axis}[
  width=15.2cm, height=3.1cm, scale only axis,
  ymin=-4.78, ymax=19.05,
  xtick={0,1,2,3,4,5.7},
  xticklabels={2021,2022,2023,2024,2025,2026H1},
  tick label style={font=\tiny},
  yticklabel style={font=\tiny},
  ylabel={亿元}, ylabel style={font=\scriptsize},
  ybar, bar width=4pt,
  grid=major, grid style={LightGray, line width=0.3pt},
  axis line style={InkGray, line width=0.5pt},
  enlarge x limits={abs=0.8},
  legend style={font=\scriptsize, at={(0.5,-0.30)}, anchor=north,
                legend columns=3, draw=none, fill=none, column sep=10pt},
]
\addplot[fill=AgriGreen!55, draw=AgriGreen, bar shift=-3.0pt] table[x=i, y=cash]{data/clean/profiles/fin_600962.dat};
\addplot[fill=SoftRed!45, draw=SoftRed, bar shift=3.0pt] table[x=i, y=ibd]{data/clean/profiles/fin_600962.dat};
\addplot[fill=AgriGreen!25, draw=AgriGreen, pattern=north east lines, bar shift=-3.0pt] table[x=x, y=cash]{data/clean/profiles/finh_600962.dat};
\addplot[fill=SoftRed!20, draw=SoftRed, pattern=north east lines, bar shift=3.0pt] table[x=x, y=ibd]{data/clean/profiles/finh_600962.dat};
\addplot[OilBlue, mark=*, mark size=1.1pt, line width=0.8pt] table[x=i, y=ocf]{data/clean/profiles/fin_600962.dat};
\addplot[only marks, OilBlue, mark=*, mark size=1.8pt] table[x=x, y=ocf]{data/clean/profiles/finh_600962.dat};
\legend{货币资金, 有息负债, 经营活动现金流净额}
\end{axis}
\end{tikzpicture}

\captionof{figure}[国投中鲁（600962）近年主要财务指标]{国投中鲁（600962）近年主要财务指标（左上：营业收入与归母净利润；右上：ROE 与资产负债率；下：货币资金、有息负债与经营活动现金流净额。
2021---2025 年为年报口径，2026H1 为 2026 年半年报/半年末，与其前一年报之间留出间隔、以斜纹或空心点区分；半年报数据未年化）}
\end{center}

\pfl{财务分析} 2026 年 6 月末资产负债率 52.0\%，较 2025 年末 下降 9.4 个百分点（样本中位数 52.2\%，所处三级行业内排第 \zhnumber{4} 位）。 2026 年 6 月末货币资金 1.57 亿元、短期债务（短期借款与一年内到期非流动负债）9.37 亿元、长期债务 0.93 亿元，有息负债合计 10.30 亿元，现金短债比 0.17 倍——缺口明显。 净有息负债 8.73 亿元，相当于其 2026 年上半年营业收入的 99\%。 流动比率 1.37、速动比率 0.59，短期指标尚可。 2026 年上半年经营活动现金流净额 5.35 亿元，经营现金流与利润同向为正，内生资金可以支撑运营。 公司披露的财务与合规风险包括：公司收到上海证券交易所《关于国投中鲁果汁股份有限公司发行股份购买资产并募集配套资金暨关联……。\footnote{\texttt{http://static.sse.com.cn/stock/disclosure/announcement/c/202607/600962\_20260724\_ETJK.pdf}}

\pfl{重大战略转型} 从公开披露看，公司的战略推进包括：2025-12 标的电子院成立于 1992 年 8 月 27 日，主营业务聚焦电子信息产业，为半导体、新型显示及其他相关领域客户……；2025-12 重组包含业绩承诺与补偿安排：标的资产 2026 年度至 2028 年度扣非净利润分别不低于约 3.12 亿元、3……；2025-12-30 第九届董事会第 12 次会议审议通过发行股份购买资产并募集配套资金方案，旨在注入国投集团旗下电子信息产业优质资产……；2026-08 重组获上交所并购重组审核委员会审核通过并取得中国证监会同意注册批复；公司披露重组定位由 2025 年报的“加速战……。 公告口径上，近三年（2023 年 9 月---2026 年 9 月）公司披露重大重组与并购类公告 32 条、对外投资与转型类公告 0 条、回购与股权激励类公告 0 条，合计公告 398 条。 第一大行业仍是“饮料制造业”，收入占比 99.5\%，与 2021 年年报相比没有方向性变化。 综合判断，公司在报告期内有明确的外延扩张或业务结构调整动作，转型方向与融资安排之间存在直接关联。\footnote{\texttt{https://www.nbd.com.cn/articles/2025-12-30/4202501.html}}

\pfl{历史融投资情况} 从现金流量表看，2025 年公司取得借款收到的现金 14.41 亿元、偿还债务支付的现金 16.17 亿元、吸收权益性投资 — 亿元，筹资活动现金流净额 -2.04 亿元（2023 年为取得借款 12.36 亿元、偿还债务 10.40 亿元）。 2026 年上半年筹资活动现金流净额为 -4.99 亿元（取得借款 0.58 亿元、偿还债务 5.47 亿元，未年化），已转为净流出，处于净还债状态。 2025 年分配股利、利润或偿付利息支付的现金合计 0.27 亿元。 公司披露的股本与债务融资脉络包括：IPO（上海证券交易所）、发行股份购买资产并募集配套资金（重大资产重组暨关联交易）。 其中规模较大的两笔为：2025-12-30发行股份购买资产并募集配套资金（重大资产重组暨关联交易）（发行股份购买电子院 100\% 股份的交易价格 60.26 亿……）；2004-06-07IPO（上海证券交易所）（发行 6,500 万股，发行价 4.80 元/股，扣除发行费……）。 股本方面，近三年披露股本变动 3 次；总股本由 2021 年末的 2.62 亿股变为 2.62 亿股（+0.00\%）。 分红方面，近三年未实施现金分红，在全部样本中属于分红能力偏弱的一端。 近三年“再融资”类公告 85 条，涉及发行股份购买资产。\footnote{\texttt{http://www.ce.cn/ztpd/tszt/gushi/2005/ssnb/PDF/200503/20/P020050320462494982913.pdf}}\footnote{\texttt{https://www.nbd.com.cn/articles/2025-12-30/4202501.html}}

\pfl{未来潜在融资需求分析} 2026 年 6 月末货币资金 1.57 亿元，而短期债务 9.37 亿元（现金短债比 0.17 倍），2026 年上半年经营现金流净额 5.35 亿元。按“短期债务减去账面货币资金”的滚动偿付口径，静态缺口约 7.80 亿元；若再计入上半年亏损的年化额（0.00 亿元），维持一年正常经营所需的外部资金约 7.80---7.80 亿元。 按第 \ref{sec:financing-need} 节的三档划分，公司属于第二档“保壳型”：近三年重大重组与并购类公告 32 条，2026 年上半年归母净利润 0.02 亿元，融资需求更多体现为“资产置换 + 维持上市地位”，而不是扩产。 公开披露的下一步融资线索包括：2025-12-30 拟向不超过 35 名特定投资者发行股份募集配套资金不超过 17.26 亿元，其中补充流动资金 7 亿元；2026-08-26 收到中国证监会同意发行股份购买资产并募集配套资金注册的批复（证监许可〔2026〕2187 号），重组进入实施阶段。\footnote{\texttt{https://www.nbd.com.cn/articles/2025-12-30/4202501.html}}\footnote{\texttt{https://stockmc.xueqiu.com/202608/600962\_20260829\_U1KD.pdf}}


\subsubsection{双塔食品（002481）}
\label{co:002481}

\pfl{公司基本情况} 双塔食品成立于 2010 年，山东招远龙口粉丝行业唯一上市公司，主营豌豆蛋白、粉丝、膳食纤维，是全球领先的豌豆蛋白与豌豆淀粉生产基地；近年受豌豆蛋白海外出口及淀粉售价影响盈利大幅回落，控股股东高比例股权质押构成治理与债务风险关注点。 公司于 2010-09-21 在深交所主板首发上市，注册资本 12.34 亿元，总股本 12.34 亿股。 截至 2026-09-24收盘价 4.16 元，流通市值 46.5 亿元，近一年-25.2\%，流通市值在三级行业“其他农产品加工”的 11 家公司中排第 \zhnumber{3} 位，近一年涨跌幅低于全样本中位数（-10.1\%）15.1 个百分点。 按 2026 年半年报披露的行业构成，收入占比前三位为农副产品加工业 100.0\%。 与 2021 年年报相比，第一大行业由“农副食品加工业”变为“农副产品加工业”，主业重心已经迁移。 发展过程中的关键节点包括：公司申报的“龙口粉丝传统制作技艺”（项目编号 Ⅷ-234）被列入第四批国家级非物质文化遗……（2014）；非公开发行股票 7,335.60 万股、发行价 17.50 元/股，募集资金总额 128……（2015）；2021 年 3 月 12 日完成 2021 年限制性股票激励计划授予登记，对核心管理人……（2021）。\footnote{\texttt{https://www.cs.com.cn/cj2020/202404/t20240412\_6401503.html}}\footnote{\texttt{http://money.finance.sina.com.cn/corp/view/vISSUE\_MarketBulletinDetail.php}}\footnote{\texttt{http://file.finance.sina.com.cn/211.154.219.97:9494/MRGG/CNSESZ\_STOCK/2021/2021-4/2021-0}}

\pfl{历史股价} 以 2021 年 1 月为起点、2026 年 9 月为终点，股价由 12.00 元变为 4.16 元，累计 -65.3\%，区间最高 13.36 元（2021 年 5 月）、最低 3.79 元（2026 年 6 月）；当前价相当于区间最高价的 31.1\%，为区间最低价的 1.10 倍。 月度收益的年化波动率 34.9\%，低于全样本中位数 37.4\%，在三级行业“其他农产品加工”的 11 家公司中波动率排第 \zhnumber{6} 位。

\begin{center}
\begin{tikzpicture}
\begin{axis}[
  width=12.6cm, height=3.9cm, scale only axis,
  xmin=2020.922, xmax=2026.888, ymin=3.52, ymax=14.30,
  xtick={2021,2022,2023,2024,2025,2026},
  yearticks,
  yticklabel style={font=\scriptsize},
  ylabel={元/股}, ylabel style={font=\scriptsize},
  grid=major, grid style={LightGray, line width=0.3pt},
  axis line style={InkGray, line width=0.5pt},
  every axis plot/.append style={line width=0.9pt},
]
\addplot[AgriGreen] table[x=x,y=y]{data/clean/profiles/px_002481.dat};
\addplot[SoftRed, dashed, line width=0.7pt] table[x=x,y=y]{data/clean/profiles/pxzz_002481.dat};
\addplot[only marks, mark=*, mark size=1.1pt, SoftRed] table[x=x,y=y]{data/clean/profiles/pxzz_002481.dat};
\end{axis}
\end{tikzpicture}

\captionof{figure}[双塔食品（002481）股价与周期骨架]{双塔食品（002481）月度收盘价（元/股）与 ZigZag 周期骨架（阈值 25\%，圆点为周期转折点）}
\end{center}

\pfl{过去周期分析} 以 25\% 的 ZigZag 阈值划分，样本区间内股价形成 2 个主升段与 3 个主跌段。 最大主升段出现在 2024 年 6 月 至 2025 年 8 月，14 个月+46.5\%；最大主跌段为 2022 年 5 月 至 2024 年 6 月，25 个月-58.5\%。 最近一次转折出现在 2026 年 6 月（下跌段自 2025 年 8 月起，10 个月-33.2\%），当前处于该段的延续之中。 公司股价与申万农林牧渔指数几乎同步见底，个股并未走出独立行情。 该子行业横跨糖、番茄制品、添加剂等多个品种，各自的供给周期与政策（进口配额、关税）是主要变量，公司 2026 年上半年实现盈利、半年 ROE 1.4\%（未年化），好于全样本中位数（0.75\%），下行周期对其报表的冲击小于同业。

\pfl{盈利情况} 2026 年上半年营业收入 11.23 亿元（同比 +7.3\%），归母净利润 0.33 亿元，同比 -39.0\%。 以年报为趋势对照，2023---2025 年营业收入由 24.56 亿元变为 21.21 亿元（两年年均复合增速 -7.1\%），归母净利润由 0.93 亿元变为 0.33 亿元。 按半年报口径，2026 年上半年的 ROE 为 1.4\%（未年化）、销售净利率 2.9\%、毛利率 19.4\%，ROE 在“其他农产品加工”11 家公司中排第 \zhnumber{6} 位。 行业同口径半年 ROE 的中位数为 0.75\%，该公司高于中位数 0.61 个百分点。 结合公司披露，2025 年度业绩预告预计归母净利润 2,800 万元至 4,000 万元，同比下降 70.41\% 至 ……；2022 年营业收入 2,381,327,771.40 元（同比增长 9.97\%），归母净利润 -310……。\footnote{\texttt{http://jjdb.sdenews.com/jjdb/jjdb/PDF/20260206/03.pdf}}\footnote{\texttt{https://disc.static.szse.cn/disc/disk03/finalpage/2023-04-25/01a4f3c3-9a0f-4731-9de8-3fc}}

\begin{center}\small\setlength{\tabcolsep}{4pt}
\captionof{table}[双塔食品（002481）关键财务指标]{双塔食品（002481）关键财务指标（合并报表口径；两个半年报期均未年化；现金短债比 $=$ 货币资金 $\div$ 短期债务）}
\begin{tabular}{@{}lrrrrr@{}}
\toprule
指标 & 2023 年 & 2024 年 & 2025 年 & 2025 年半年报 & 2026 年半年报 \\
\midrule
营业收入（亿元） & 24.56 & 24.48 & 21.21 & 10.47 & 11.23 \\
归母净利润（亿元） & 0.93 & 0.95 & 0.33 & 0.54 & 0.33 \\
毛利率（\%） & 14.8 & 15.0 & 20.2 & 17.9 & 19.4 \\
ROE（\%） & 3.9 & 3.6 & 1.3 & 2.0 & 1.4 \\
资产负债率（\%） & 32.5 & 43.7 & 47.6 & 41.5 & 44.8 \\
经营活动现金流净额（亿元） & 2.31 & 3.25 & 0.91 & -1.79 & 1.27 \\
货币资金（亿元） & 4.08 & 11.28 & 7.32 & 6.10 & 3.84 \\
有息负债合计（亿元） & 7.20 & 14.30 & 16.02 & 12.94 & 12.29 \\
现金短债比（倍） & 0.69 & 0.86 & 0.56 & 0.54 & 0.43 \\
\bottomrule
\end{tabular}
\end{center}

\begin{center}
\begin{tikzpicture}
\begin{groupplot}[
  group style={group size=2 by 1, horizontal sep=0.95cm},
  width=7.2cm, height=4.6cm,
  tick label style={font=\tiny},
  title style={font=\scriptsize\heiti},
  yticklabel style={font=\tiny},
  grid=major, grid style={LightGray, line width=0.3pt},
  axis line style={InkGray, line width=0.5pt},
  enlarge x limits={abs=0.8},
  legend style={font=\scriptsize, at={(0.5,-0.22)}, anchor=north,
                legend columns=2, draw=none, fill=none, column sep=6pt},
]
\nextgroupplot[
  ybar, bar width=4pt,
  ymin=-5.88, ymax=27.88,
  xtick={0,1,2,3,4,5.7},
  xticklabels={2021,2022,2023,2024,2025,2026H1},
  ylabel={亿元},
]
\addplot[fill=AgriGreen!75, draw=AgriGreen, bar shift=-2.6pt] table[x=i, y=rev]{data/clean/profiles/fin_002481.dat};
\addplot[fill=SoftRed!60, draw=SoftRed, bar shift=2.6pt] table[x=i, y=ni]{data/clean/profiles/fin_002481.dat};
\addplot[fill=AgriGreen!30, draw=AgriGreen, pattern=north east lines, bar shift=-2.6pt] table[x=x, y=rev]{data/clean/profiles/finh_002481.dat};
\addplot[fill=SoftRed!20, draw=SoftRed, pattern=north east lines, bar shift=2.6pt] table[x=x, y=ni]{data/clean/profiles/finh_002481.dat};
\legend{营业收入（年/半年）, 归母净利润（年/半年）}
\nextgroupplot[
  ymin=-16.9, ymax=53.5,
  xtick={0,1,2,3,4,5.7},
  xticklabels={2021,2022,2023,2024,2025,2026H1},
  ylabel={\%},
]
\addplot[AgriGold, mark=*, mark size=1.3pt] table[x=i, y=roe]{data/clean/profiles/fin_002481.dat};
\addplot[OilBlue, mark=square*, mark size=1.3pt] table[x=i, y=debt]{data/clean/profiles/fin_002481.dat};
\addplot[only marks, AgriGold, mark=*, mark size=2.0pt] table[x=x, y=roe]{data/clean/profiles/finh_002481.dat};
\addplot[only marks, OilBlue, mark=square*, mark size=2.0pt] table[x=x, y=debt]{data/clean/profiles/finh_002481.dat};
\legend{ROE, 资产负债率}
\end{groupplot}
\end{tikzpicture}
\begin{tikzpicture}
\begin{axis}[
  width=15.2cm, height=3.1cm, scale only axis,
  ymin=-1.60, ymax=17.95,
  xtick={0,1,2,3,4,5.7},
  xticklabels={2021,2022,2023,2024,2025,2026H1},
  tick label style={font=\tiny},
  yticklabel style={font=\tiny},
  ylabel={亿元}, ylabel style={font=\scriptsize},
  ybar, bar width=4pt,
  grid=major, grid style={LightGray, line width=0.3pt},
  axis line style={InkGray, line width=0.5pt},
  enlarge x limits={abs=0.8},
  legend style={font=\scriptsize, at={(0.5,-0.30)}, anchor=north,
                legend columns=3, draw=none, fill=none, column sep=10pt},
]
\addplot[fill=AgriGreen!55, draw=AgriGreen, bar shift=-3.0pt] table[x=i, y=cash]{data/clean/profiles/fin_002481.dat};
\addplot[fill=SoftRed!45, draw=SoftRed, bar shift=3.0pt] table[x=i, y=ibd]{data/clean/profiles/fin_002481.dat};
\addplot[fill=AgriGreen!25, draw=AgriGreen, pattern=north east lines, bar shift=-3.0pt] table[x=x, y=cash]{data/clean/profiles/finh_002481.dat};
\addplot[fill=SoftRed!20, draw=SoftRed, pattern=north east lines, bar shift=3.0pt] table[x=x, y=ibd]{data/clean/profiles/finh_002481.dat};
\addplot[OilBlue, mark=*, mark size=1.1pt, line width=0.8pt] table[x=i, y=ocf]{data/clean/profiles/fin_002481.dat};
\addplot[only marks, OilBlue, mark=*, mark size=1.8pt] table[x=x, y=ocf]{data/clean/profiles/finh_002481.dat};
\legend{货币资金, 有息负债, 经营活动现金流净额}
\end{axis}
\end{tikzpicture}

\captionof{figure}[双塔食品（002481）近年主要财务指标]{双塔食品（002481）近年主要财务指标（左上：营业收入与归母净利润；右上：ROE 与资产负债率；下：货币资金、有息负债与经营活动现金流净额。
2021---2025 年为年报口径，2026H1 为 2026 年半年报/半年末，与其前一年报之间留出间隔、以斜纹或空心点区分；半年报数据未年化）}
\end{center}

\pfl{财务分析} 2026 年 6 月末资产负债率 44.8\%，较 2025 年末 下降 2.8 个百分点（样本中位数 52.2\%，所处三级行业内排第 \zhnumber{8} 位）。 2026 年 6 月末货币资金 3.84 亿元、短期债务（短期借款与一年内到期非流动负债）8.95 亿元、长期债务 3.34 亿元，有息负债合计 12.29 亿元，现金短债比 0.43 倍——覆盖偏紧。 净有息负债 8.44 亿元，相当于其 2026 年上半年营业收入的 75\%。 流动比率 1.21、速动比率 0.59，短期指标尚可。 2026 年上半年经营活动现金流净额 1.27 亿元，经营现金流与利润同向为正，内生资金可以支撑运营。 公司披露的财务与合规风险包括：股权质押风险：截至 2026 年 9 月 5 日公司整体质押股份 3.21 亿股，占总股……；控股股东部分股份解除质押及再质押公告披露：控股股东及其一致行动人未来半年内合计到期的质押……。\footnote{\texttt{https://finance.sina.com.cn/stock/aiassist/zy/2026-09-07/doc-iniqxyne9658815.shtml}}\footnote{\texttt{https://finance.sina.com.cn/roll/2026-09-22/doc-inisrnpr7417131.shtml}}

\pfl{重大战略转型} 从公开披露看，公司的战略推进包括：2020—2026 持续聚焦豌豆蛋白、粉丝、膳食纤维主业，优化产品结构、发展循环产业链条、加大研发投入以提升产品附加值；产品延伸至人……；2020-03 发布非公开发行预案，拟投向万吨白蛋白提取综合加工、万吨豌豆功能低聚糖提取加工、豌豆精深加工综合利用等项目，扩张高……。 公告口径上，近三年（2023 年 9 月---2026 年 9 月）公司披露重大重组与并购类公告 2 条、对外投资与转型类公告 1 条、回购与股权激励类公告 33 条，合计公告 322 条。 主营结构的迁移已经落到报表上：2021 年年报的第一大行业为“农副食品加工业”（100.0\%），2026 年半年报变为“农副产品加工业”（100.0\%）。 综合判断，公司在报告期内有明确的外延扩张或业务结构调整动作，转型方向与融资安排之间存在直接关联。\footnote{\texttt{https://money.finance.sina.com.cn/corp/view/vCB\_AllBulletinDetail.php}}\footnote{\texttt{https://pdf.dfcfw.com/pdf/H3\_AP202003041375833368\_1.pdf}}

\begin{center}\small\setlength{\tabcolsep}{4pt}
\captionof{table}[双塔食品（002481）历史融资与资本运作]{双塔食品（002481）历史融资与资本运作（据公司公告与公开报道整理；金额为披露值，含首次公开发行、再融资、债券、银行借款与重整投资等）}
\begin{tabular}{@{}p{1.5cm}p{3.2cm}p{2.4cm}p{6.3cm}@{}}
\toprule
时间 & 方式 & 规模 & 用途与说明 \\
\midrule
2010-09-08 & IPO（深圳证券交易所） & 发行 1,500 万股，发行价 39.80 元/股，占发行后总股本的 25.00\%，发行后总股本 6,000 万股 & 年产 20,000 吨粉丝项目；2010 年 7 月 23 日首发申请获通过…… \\
2015-01-14 & 非公开发行 A 股（深交所） & 发行 7,335.60 万股，发行价 17.50 元/股，募集资金总额 128,373.00 万元、募集资金净额 126,524.32 万元 & 新增股份上市首日股价不除权；认购股票限售期十二个月…… \\
2020-03-03 & 非公开发行股票预案（未在本底稿核实到实施完成的公开信息） & 拟募集资金总额不超过 21.5 亿元 & 万吨白蛋白提取综合加工项目 5.1 亿元、万吨豌豆功能低聚糖提取加工项目 4.45 亿元…… \\
2020-11—2021 & 股份回购（用于股权激励） & --- & 回购股份用于股权激励；2020 年 11 月 28 日第五届董事会第六次会议审议通过以集…… \\
\bottomrule
\end{tabular}
\end{center}

\pfl{历史融投资情况} 从现金流量表看，2025 年公司取得借款收到的现金 26.46 亿元、偿还债务支付的现金 24.71 亿元、吸收权益性投资 0.01 亿元，筹资活动现金流净额 0.23 亿元（2023 年为取得借款 8.41 亿元、偿还债务 11.03 亿元）。 2026 年上半年筹资活动现金流净额为 -0.97 亿元（取得借款 10.72 亿元、偿还债务 14.41 亿元，未年化），已转为净流出，处于净还债状态。 2025 年分配股利、利润或偿付利息支付的现金合计 0.55 亿元。 公司披露的股本与债务融资脉络包括：IPO（深圳证券交易所）、股份回购（用于股权激励）、非公开发行 A 股（深交所）、非公开发行股票预案（未在本底稿核实到实施完成的公开信息）。 其中规模较大的两笔为：2020-03-03非公开发行股票预案（未在本底稿核实到实施完成的公开信息）（拟募集资金总额不超过 21.5 亿元）；2015-01-14非公开发行 A 股（深交所）（发行 7,335.60 万股，发行价 17.50 元/股，募……）。 股本方面，近三年披露股本变动 3 次；总股本由 2021 年末的 12.43 亿股变为 12.34 亿股（-0.78\%）。 分红方面，近三年现金分红 4 次、累计每股派现 0.08 元（最近一次 2025-06-30）——在利润承压的阶段仍维持了分红。 近三年“再融资”类公告 1 条，涉及定向增发。\footnote{\texttt{http://vip.stock.finance.sina.com.cn/corp/view/vISSUE\_RaiseExplanationDetail.php}}\footnote{\texttt{http://money.finance.sina.com.cn/corp/view/vISSUE\_MarketBulletinDetail.php}}\footnote{\texttt{https://pdf.dfcfw.com/pdf/H3\_AP202003041375833368\_1.pdf}}

\pfl{未来潜在融资需求分析} 2026 年 6 月末货币资金 3.84 亿元，而短期债务 8.95 亿元（现金短债比 0.43 倍），2026 年上半年经营现金流净额 1.27 亿元。按“短期债务减去账面货币资金”的滚动偿付口径，静态缺口约 5.11 亿元；若再计入上半年亏损的年化额（0.00 亿元），维持一年正常经营所需的外部资金约 5.11---5.11 亿元。 公司 2026 年 6 月末资产负债率 44.8\%、2026 年上半年 ROE 1.4\%（未年化），资产负债表尚有余量，融资需求不是“补血型”的刚性需求，而是围绕并购整合、项目投入与营运资金的主动性需求。 公开披露的下一步融资线索包括：2020-03 拟定增募资不超过 21.5 亿元（含补充流动资金 5.3 亿元），本底稿未查得该次发行的完成或终止公告；2026-09 控股股东及一致行动人质押股份未来半年内到期部分对应融资余额 5,000 万元，存在滚动质押融资需求。\footnote{\texttt{https://pdf.dfcfw.com/pdf/H3\_AP202003041375833368\_1.pdf}}\footnote{\texttt{https://finance.sina.com.cn/roll/2026-09-22/doc-inisrnpr7417131.shtml}}


\subsubsection{永顺泰（001338）}
\label{co:001338}

\pfl{公司基本情况} 1987 年设立至今，永顺泰广东省属国企粤海控股旗下的啤酒麦芽制造商，国内产销规模最大的麦芽制造企业，产能规模居亚洲第一、世界第四；2022年11月登陆深交所主板，业绩受大麦成本与啤酒行业景气度影响显著。 公司于 2022-11-16 在深交所主板首发上市，注册资本 5.02 亿元，总股本 5.02 亿股。 截至 2026-09-24收盘价 9.10 元，流通市值 45.7 亿元，近一年-19.8\%，流通市值在三级行业“其他农产品加工”的 11 家公司中排第 \zhnumber{4} 位，近一年涨跌幅低于全样本中位数（-10.1\%）9.7 个百分点。 按 2026 年半年报披露的行业构成，收入占比前三位为啤酒制造 95.2\%；其他(补充) 4.8\%。 发展过程中的关键节点包括：成立宁麦公司（宁波基地）（1996）；成立秦麦公司、昌麦公司（2005）；2017年12月5日粤海控股董事会同意设立广东粤海永顺泰麦芽有限公司，注册资本3,500……（2017）。\footnote{\texttt{https://file.finance.qq.com/finance/hs/pdf/2024/05/16/1220074266.PDF}}\footnote{\texttt{http://pdf.dfcfw.com/pdf/H2\_AN202210311579684054\_1.pdf}}

\pfl{历史股价} 自 2022 年 11 月至 2026 年 9 月，股价由 21.30 元变为 9.10 元，累计 -57.3\%，区间最高 21.30 元（2022 年 11 月）、最低 8.39 元（2026 年 6 月）；当前价相当于区间最高价的 42.7\%，为区间最低价的 1.08 倍。 月度收益的年化波动率 28.5\%，低于全样本中位数 37.4\%，在三级行业“其他农产品加工”的 11 家公司中波动率排第 \zhnumber{10} 位。 从事件看，2022 年 11 月 至 2024 年 8 月的下跌（-55.9\%）与公司2022年9月1日取得证监许可[2022]2009号核准批复；募集资金于20……（2022-11-16）在时间上重合；2024 年 8 月 至 2026 年 4 月的上涨（+41.2\%）与公司2025年营业收入39.27亿元，同比下降8.29\%；归属于上市公司股东的净……（2025-12-31）在时间上重合——股价的拐点多数能在公司自身的资本运作与风险事件上找到对应。\footnote{\texttt{https://epaper.stcn.com/pic/202604/21/aeeabe2acff0ed422e6db66f0be72de7.pdf}}\footnote{\texttt{https://vip.stock.finance.sina.com.cn/corp/view/vCB\_AllBulletinDetail.php}}

\begin{center}
\begin{tikzpicture}
\begin{axis}[
  width=12.6cm, height=3.9cm, scale only axis,
  xmin=2022.755, xmax=2026.888, ymin=7.80, ymax=22.79,
  xtick={2022,2023,2024,2025,2026},
  yearticks,
  yticklabel style={font=\scriptsize},
  ylabel={元/股}, ylabel style={font=\scriptsize},
  grid=major, grid style={LightGray, line width=0.3pt},
  axis line style={InkGray, line width=0.5pt},
  every axis plot/.append style={line width=0.9pt},
]
\addplot[AgriGreen] table[x=x,y=y]{data/clean/profiles/px_001338.dat};
\addplot[SoftRed, dashed, line width=0.7pt] table[x=x,y=y]{data/clean/profiles/pxzz_001338.dat};
\addplot[only marks, mark=*, mark size=1.1pt, SoftRed] table[x=x,y=y]{data/clean/profiles/pxzz_001338.dat};
\end{axis}
\end{tikzpicture}

\captionof{figure}[永顺泰（001338）股价与周期骨架]{永顺泰（001338）月度收盘价（元/股）与 ZigZag 周期骨架（阈值 25\%，圆点为周期转折点）}
\end{center}

\pfl{过去周期分析} 以 25\% 的 ZigZag 阈值划分，样本区间内股价形成 1 个主升段与 2 个主跌段。 最大主跌段为 2022 年 11 月 至 2024 年 8 月，21 个月-55.9\%；最大主升段出现在 2024 年 8 月 至 2026 年 4 月，20 个月+41.2\%。 最近一次转折出现在 2026 年 6 月（下跌段自 2026 年 4 月起，2 个月-36.8\%），当前处于该段的延续之中。 公司股价与申万农林牧渔指数几乎同步见底，个股并未走出独立行情。 该子行业横跨糖、番茄制品、添加剂等多个品种，各自的供给周期与政策（进口配额、关税）是主要变量，公司 2026 年上半年实现盈利、半年 ROE 1.4\%（未年化），好于全样本中位数（0.75\%），下行周期对其报表的冲击小于同业。

\pfl{盈利情况} 2026 年上半年营业收入 19.15 亿元（同比 -13.9\%），归母净利润 0.51 亿元，同比 -69.6\%。 以年报为趋势对照，2023---2025 年营业收入由 48.40 亿元变为 39.27 亿元（两年年均复合增速 -9.9\%），归母净利润由 1.74 亿元变为 2.37 亿元。 按半年报口径，2026 年上半年的 ROE 为 1.4\%（未年化）、销售净利率 2.7\%、毛利率 8.3\%，ROE 在“其他农产品加工”11 家公司中排第 \zhnumber{5} 位。 同期全样本半年 ROE 中位数 0.75\%，公司与之相差 0.68 个百分点（略好于中位数公司）。 结合公司披露，2025年营业收入39.27亿元，同比下降8.29\%；归属于上市公司股东的净利润2.37亿元（236,7……；2025年麦芽产品主营业务收入37.16亿元，同比下降8.62\%；毛利率12.69\%，同比增长0.09个……。\footnote{\texttt{https://vip.stock.finance.sina.com.cn/corp/view/vCB\_AllBulletinDetail.php}}\footnote{\texttt{https://epaper.stcn.com/pic/202604/21/aeeabe2acff0ed422e6db66f0be72de7.pdf}}

\begin{center}\small\setlength{\tabcolsep}{4pt}
\captionof{table}[永顺泰（001338）关键财务指标]{永顺泰（001338）关键财务指标（合并报表口径；两个半年报期均未年化；现金短债比 $=$ 货币资金 $\div$ 短期债务）}
\begin{tabular}{@{}lrrrrr@{}}
\toprule
指标 & 2023 年 & 2024 年 & 2025 年 & 2025 年半年报 & 2026 年半年报 \\
\midrule
营业收入（亿元） & 48.40 & 42.82 & 39.27 & 22.23 & 19.15 \\
归母净利润（亿元） & 1.74 & 2.99 & 2.37 & 1.69 & 0.51 \\
毛利率（\%） & 8.4 & 11.9 & 12.1 & 12.0 & 8.3 \\
ROE（\%） & 5.3 & 8.6 & 6.5 & 4.8 & 1.4 \\
资产负债率（\%） & 25.3 & 14.8 & 10.0 & 15.0 & 21.1 \\
经营活动现金流净额（亿元） & 30.85 & 24.44 & 15.61 & 11.56 & 2.18 \\
货币资金（亿元） & 5.53 & 8.95 & 9.75 & 9.65 & 14.14 \\
有息负债合计（亿元） & 2.20 & 0.64 & 0.04 & 2.28 & 7.41 \\
现金短债比（倍） & 2.52 & 16.51 & 1,606.40 & 4.36 & 1.92 \\
\bottomrule
\end{tabular}
\end{center}

\begin{center}
\begin{tikzpicture}
\begin{groupplot}[
  group style={group size=2 by 1, horizontal sep=0.95cm},
  width=7.2cm, height=4.6cm,
  tick label style={font=\tiny},
  title style={font=\scriptsize\heiti},
  yticklabel style={font=\tiny},
  grid=major, grid style={LightGray, line width=0.3pt},
  axis line style={InkGray, line width=0.5pt},
  enlarge x limits={abs=0.8},
  legend style={font=\scriptsize, at={(0.5,-0.22)}, anchor=north,
                legend columns=2, draw=none, fill=none, column sep=6pt},
]
\nextgroupplot[
  ybar, bar width=4pt,
  ymin=-4.84, ymax=54.21,
  xtick={0,1,2,3,4,5.7},
  xticklabels={2021,2022,2023,2024,2025,2026H1},
  ylabel={亿元},
]
\addplot[fill=AgriGreen!75, draw=AgriGreen, bar shift=-2.6pt] table[x=i, y=rev]{data/clean/profiles/fin_001338.dat};
\addplot[fill=SoftRed!60, draw=SoftRed, bar shift=2.6pt] table[x=i, y=ni]{data/clean/profiles/fin_001338.dat};
\addplot[fill=AgriGreen!30, draw=AgriGreen, pattern=north east lines, bar shift=-2.6pt] table[x=x, y=rev]{data/clean/profiles/finh_001338.dat};
\addplot[fill=SoftRed!20, draw=SoftRed, pattern=north east lines, bar shift=2.6pt] table[x=x, y=ni]{data/clean/profiles/finh_001338.dat};
\legend{营业收入（年/半年）, 归母净利润（年/半年）}
\nextgroupplot[
  ymin=-2.2, ymax=29.6,
  xtick={0,1,2,3,4,5.7},
  xticklabels={2021,2022,2023,2024,2025,2026H1},
  ylabel={\%},
]
\addplot[AgriGold, mark=*, mark size=1.3pt] table[x=i, y=roe]{data/clean/profiles/fin_001338.dat};
\addplot[OilBlue, mark=square*, mark size=1.3pt] table[x=i, y=debt]{data/clean/profiles/fin_001338.dat};
\addplot[only marks, AgriGold, mark=*, mark size=2.0pt] table[x=x, y=roe]{data/clean/profiles/finh_001338.dat};
\addplot[only marks, OilBlue, mark=square*, mark size=2.0pt] table[x=x, y=debt]{data/clean/profiles/finh_001338.dat};
\legend{ROE, 资产负债率}
\end{groupplot}
\end{tikzpicture}
\begin{tikzpicture}
\begin{axis}[
  width=15.2cm, height=3.1cm, scale only axis,
  ymin=-3.08, ymax=34.55,
  xtick={0,1,2,3,4,5.7},
  xticklabels={2021,2022,2023,2024,2025,2026H1},
  tick label style={font=\tiny},
  yticklabel style={font=\tiny},
  ylabel={亿元}, ylabel style={font=\scriptsize},
  ybar, bar width=4pt,
  grid=major, grid style={LightGray, line width=0.3pt},
  axis line style={InkGray, line width=0.5pt},
  enlarge x limits={abs=0.8},
  legend style={font=\scriptsize, at={(0.5,-0.30)}, anchor=north,
                legend columns=3, draw=none, fill=none, column sep=10pt},
]
\addplot[fill=AgriGreen!55, draw=AgriGreen, bar shift=-3.0pt] table[x=i, y=cash]{data/clean/profiles/fin_001338.dat};
\addplot[fill=SoftRed!45, draw=SoftRed, bar shift=3.0pt] table[x=i, y=ibd]{data/clean/profiles/fin_001338.dat};
\addplot[fill=AgriGreen!25, draw=AgriGreen, pattern=north east lines, bar shift=-3.0pt] table[x=x, y=cash]{data/clean/profiles/finh_001338.dat};
\addplot[fill=SoftRed!20, draw=SoftRed, pattern=north east lines, bar shift=3.0pt] table[x=x, y=ibd]{data/clean/profiles/finh_001338.dat};
\addplot[OilBlue, mark=*, mark size=1.1pt, line width=0.8pt] table[x=i, y=ocf]{data/clean/profiles/fin_001338.dat};
\addplot[only marks, OilBlue, mark=*, mark size=1.8pt] table[x=x, y=ocf]{data/clean/profiles/finh_001338.dat};
\legend{货币资金, 有息负债, 经营活动现金流净额}
\end{axis}
\end{tikzpicture}

\captionof{figure}[永顺泰（001338）近年主要财务指标]{永顺泰（001338）近年主要财务指标（左上：营业收入与归母净利润；右上：ROE 与资产负债率；下：货币资金、有息负债与经营活动现金流净额。
2021---2025 年为年报口径，2026H1 为 2026 年半年报/半年末，与其前一年报之间留出间隔、以斜纹或空心点区分；半年报数据未年化）}
\end{center}

\pfl{财务分析} 2026 年 6 月末资产负债率 21.1\%，较 2025 年末 上升 11.1 个百分点（样本中位数 52.2\%，所处三级行业内排第 \zhnumber{2} 位）。 2026 年 6 月末货币资金 14.14 亿元、短期债务（短期借款与一年内到期非流动负债）7.38 亿元、长期债务 0.03 亿元，有息负债合计 7.41 亿元，现金短债比 1.92 倍——覆盖充分。 净有息负债 -6.73 亿元，相当于其 2026 年上半年营业收入的 -35\%。 流动比率 3.43、速动比率 2.39，短期指标尚可。 2026 年上半年经营活动现金流净额 2.18 亿元，经营现金流与利润同向为正，内生资金可以支撑运营。 2026 年 6 月末在建工程 1.33 亿元，较上年末增加 0.05 亿元，仍有在建投入。

\pfl{重大战略转型} 从公开披露看，公司的战略推进包括：2019-2020 陆续开展广麦4期扩建项目和年产13万吨中高档啤酒麦芽项目等产能扩张项目；2025-06-30 强化产购销协同，2025年上半年麦芽产品产能56.92万吨（同比+8.27\%），销量62.48万吨（同比+19.……；2025-2026 推进宝麦公司新建5万吨/年特制麦芽生产线项目，已开工建设/试生产，推动产品结构向高端化、定制化升级；2026-06-30 威士忌麦芽和特种麦芽销量创半年度历史新高，新业态拓展成为亮点。 公告口径上，近三年（2023 年 9 月---2026 年 9 月）公司披露重大重组与并购类公告 0 条、对外投资与转型类公告 0 条、回购与股权激励类公告 0 条，合计公告 281 条。 第一大行业仍是“啤酒制造”，收入占比 95.2\%，与 2021 年年报相比没有方向性变化。 综合判断，公司在报告期内有明确的外延扩张或业务结构调整动作，转型方向与融资安排之间存在直接关联。\footnote{\texttt{https://file.finance.qq.com/finance/hs/pdf/2024/05/16/1220074266.PDF}}\footnote{\texttt{https://stock.finance.sina.com.cn/stock/go.php/vReport\_Show/kind/search/rptid/8095468648}}\footnote{\texttt{https://money.finance.sina.com.cn/corp/view/vCB\_AllBulletinDetail.php}}

\begin{center}\small\setlength{\tabcolsep}{4pt}
\captionof{table}[永顺泰（001338）历史融资与资本运作]{永顺泰（001338）历史融资与资本运作（据公司公告与公开报道整理；金额为披露值，含首次公开发行、再融资、债券、银行借款与重整投资等）}
\begin{tabular}{@{}p{1.5cm}p{3.2cm}p{2.4cm}p{6.3cm}@{}}
\toprule
时间 & 方式 & 规模 & 用途与说明 \\
\midrule
2022-11 & IPO募投项目明细 & 广麦4期扩建项目投资总额50,305.00万元、拟投入募集资金424,534,039.08元；年产13万吨中高档啤酒麦芽项目投资总额52,204.00万元、拟投入募集资金369,711,382.37元；合计投资总额102,509.00万元、拟投入募集资金794,245,421.45元 & 产能扩张；年产13万吨中高档啤酒麦芽项目实施地点为江苏省扬州市宝应县安宜工业园…… \\
2022-11-16 & IPO（深交所主板） & 发行125,432,708股，发行价6.82元/股，募集资金总额8.55亿元（855,451,068.56元），扣除发行费用61,205,647.11元后募集资金净额794,245,421.45元 & 广麦4期扩建项目、年产13万吨中高档啤酒麦芽项目…… \\
2023-12-08 & 募投项目结项 & 首次公开发行募投项目（广麦4期扩建项目、年产13万吨中高档啤酒麦芽项目）结项，节余募集资金永久补充流动资金（节余金额在来源页面被截断，不予引用） & 补充流动资金；第二届董事会第三次会议、第二届监事会第三次会议审议通过（节余金额原文被截断…… \\
2026-06-30 & 闲置募集资金现金管理 & 不超过人民币1,450万元 & 购买现金管理产品；使用期限为自董事会审议通过之日起12个月内，可循环滚动使用…… \\
\bottomrule
\end{tabular}
\end{center}

\pfl{历史融投资情况} 从现金流量表看，2025 年公司取得借款收到的现金 0.20 亿元、偿还债务支付的现金 11.19 亿元、吸收权益性投资 — 亿元，筹资活动现金流净额 -12.14 亿元（2023 年为取得借款 16.02 亿元、偿还债务 45.79 亿元）。 2026 年上半年筹资活动现金流净额为 2.37 亿元（取得借款 — 亿元、偿还债务 2.28 亿元，未年化），仍处于净融入状态。 2025 年分配股利、利润或偿付利息支付的现金合计 1.57 亿元。 公开信息中与公司直接相关的融资与资本运作包括：IPO募投项目明细、IPO（深交所主板）、募投项目结项、闲置募集资金现金管理。 其中规模较大的两笔为：2026-06-30闲置募集资金现金管理（不超过人民币1,450万元）；2022-11IPO募投项目明细（广麦4期扩建项目投资总额50,305.00万元、拟投入募集资……）。 股本方面，近三年披露股本变动 4 次（限售股份上市 1 次）；总股本由 2021 年末的 5.02 亿股变为 5.02 亿股（+0.00\%）。 分红方面，近三年现金分红 4 次、累计每股派现 0.71 元（最近一次 2025-12-31）——在利润承压的阶段仍维持了分红。 近三年“再融资”类公告 1 条，涉及定向增发。\footnote{\texttt{https://money.finance.sina.com.cn/corp/view/vCB\_AllBulletinDetail.php}}\footnote{\texttt{https://epaper.stcn.com/pic/202604/21/aeeabe2acff0ed422e6db66f0be72de7.pdf}}\footnote{\texttt{https://money.finance.sina.com.cn/corp/view/vCB\_AllBulletinDetail.php}}

\pfl{未来潜在融资需求分析} 2026 年 6 月末货币资金 14.14 亿元，而短期债务 7.38 亿元（现金短债比 1.92 倍），2026 年上半年经营现金流净额 2.18 亿元。账面货币资金可以覆盖短期债务，缺口不体现在静态偿付层面。 公司 2026 年 6 月末资产负债率 21.1\%、2026 年上半年 ROE 1.4\%（未年化），资产负债表尚有余量，融资需求不是“补血型”的刚性需求，而是围绕并购整合、项目投入与营运资金的主动性需求。 公开披露的下一步融资线索包括：2026-06-30 公司（含子公司）使用不超过人民币1,450万元的暂时闲置募集资金进行现金管理，反映账上募集资金已基本使用完毕；2026-06-30 2026年6月30日资产负债率仅21.09\%，归母所有者权益36.36亿元，资产负债表留有加杠杆空间（未查得公开的授信或……。 资金需求还要叠加在建工程：期末在建工程 1.33 亿元、2026 年上半年购建固定资产等支出 0.35 亿元（未年化），这部分支出在周期底部通常会被主动放缓。\footnote{\texttt{https://money.finance.sina.com.cn/corp/view/vCB\_AllBulletinDetail.php}}\footnote{\texttt{http://money.finance.sina.com.cn/corp/view/vCB\_AllBulletinDetail.php}}


\subsubsection{晨光生物（300138）}
\label{co:300138}

\pfl{公司基本情况} 晨光生物成立于 2000 年，河北曲周起家的天然植物提取物龙头，辣椒红色素、辣椒精、叶黄素三大主力产品位居世界第一或前列，棉籽类业务规模大但毛利薄；2010年创业板上市，成本领先与全球份额是核心逻辑。 公司于 2010-11-05 在深交所创业板首发上市，注册资本 4.83 亿元，总股本 4.83 亿股。 截至 2026-09-24收盘价 11.20 元，流通市值 44.6 亿元，近一年-11.9\%，流通市值在三级行业“其他农产品加工”的 11 家公司中排第 \zhnumber{5} 位，近一年涨跌幅低于全样本中位数（-10.1\%）1.9 个百分点。 按 2025 年年报披露的行业构成，收入占比前三位为天然植物提取行业 100.0\%。 发展过程中的关键节点包括：2010年10月25日发行，2010年11月5日在深交所创业板上市，证券代码300138……（2010）；2020年6月17日公开发行630万张可转债（“晨光转债”），2020年7月13日在深交……（2020）；“晨光转债”触发有条件赎回并于2021年3月5日赎回摘牌，累计转股增加股本51,212,……（2021）。\footnote{\texttt{http://file.finance.sina.com.cn/211.154.219.97:9494/MRGG/CNSESZ\_STOCK/2011/2011-3/2011-0}}\footnote{\texttt{http://vip.stock.finance.sina.com.cn/corp/view/vISSUE\_MarketBulletinDetail.php}}\footnote{\texttt{http://static.cninfo.com.cn/finalpage/2024-08-27/1220981353.PDF}}

\pfl{历史股价} 自 2021 年 1 月至 2026 年 9 月，股价由 16.50 元变为 11.20 元，累计 -32.1\%，区间最高 18.53 元（2023 年 1 月）、最低 7.52 元（2024 年 8 月）；当前价相当于区间最高价的 60.4\%，为区间最低价的 1.49 倍。 月度收益的年化波动率 33.1\%，低于全样本中位数 37.4\%，在三级行业“其他农产品加工”的 11 家公司中波动率排第 \zhnumber{9} 位。 从事件看，2024 年 8 月 至 2026 年 2 月的上涨（+92.7\%）与公司可转债募集资金已累计使用38,867.84万元（截至2023年12月31日口……（2025-12-31）在时间上重合；2023 年 1 月 至 2024 年 8 月的下跌（-59.4\%）与公司控股子公司新疆晨光生物科技股份有限公司决定终止向不特定合格投资者公开发行并在……（2024-06-30）在时间上重合——股价的拐点多数能在公司自身的资本运作与风险事件上找到对应。\footnote{\texttt{https://stockn.xueqiu.com/SZ300138/20240417682702.pdf}}\footnote{\texttt{http://static.cninfo.com.cn/finalpage/2024-08-27/1220981353.PDF}}

\begin{center}
\begin{tikzpicture}
\begin{axis}[
  width=12.6cm, height=3.9cm, scale only axis,
  xmin=2020.922, xmax=2026.888, ymin=6.99, ymax=19.83,
  xtick={2021,2022,2023,2024,2025,2026},
  yearticks,
  yticklabel style={font=\scriptsize},
  ylabel={元/股}, ylabel style={font=\scriptsize},
  grid=major, grid style={LightGray, line width=0.3pt},
  axis line style={InkGray, line width=0.5pt},
  every axis plot/.append style={line width=0.9pt},
]
\addplot[AgriGreen] table[x=x,y=y]{data/clean/profiles/px_300138.dat};
\addplot[SoftRed, dashed, line width=0.7pt] table[x=x,y=y]{data/clean/profiles/pxzz_300138.dat};
\addplot[only marks, mark=*, mark size=1.1pt, SoftRed] table[x=x,y=y]{data/clean/profiles/pxzz_300138.dat};
\end{axis}
\end{tikzpicture}

\captionof{figure}[晨光生物（300138）股价与周期骨架]{晨光生物（300138）月度收盘价（元/股）与 ZigZag 周期骨架（阈值 25\%，圆点为周期转折点）}
\end{center}

\pfl{过去周期分析} 以 25\% 的 ZigZag 阈值划分，样本区间内股价形成 3 个主升段与 3 个主跌段。 最大主跌段为 2023 年 1 月 至 2024 年 8 月，19 个月-59.4\%；最大主升段出现在 2024 年 8 月 至 2026 年 2 月，18 个月+92.7\%。 最近一次转折出现在 2026 年 7 月（上涨段自 2026 年 6 月起，1 个月+29.3\%），当前处于该段之后的震荡之中。 公司股价较申万农林牧渔指数提前约 21 个月见底，是板块中较早定价周期反转的公司之一。 该子行业横跨糖、番茄制品、添加剂等多个品种，各自的供给周期与政策（进口配额、关税）是主要变量，公司 2026 年上半年实现盈利、半年 ROE 5.2\%（未年化），好于全样本中位数（0.75\%），下行周期对其报表的冲击小于同业。

\pfl{盈利情况} 2026 年上半年营业收入 31.59 亿元（同比 -13.6\%），归母净利润 1.79 亿元，同比 -17.0\%。 以年报为趋势对照，2023---2025 年营业收入由 68.72 亿元变为 65.59 亿元（两年年均复合增速 -2.3\%），归母净利润由 4.80 亿元变为 3.69 亿元。 按半年报口径，2026 年上半年的 ROE 为 5.2\%（未年化）、销售净利率 5.8\%、毛利率 13.5\%，ROE 在“其他农产品加工”11 家公司中排第 \zhnumber{2} 位。 行业同口径半年 ROE 的中位数为 0.75\%，该公司高于中位数 4.40 个百分点。 结合公司披露，2025年棉籽类业务减少贸易业务量、收入同比下降，但随行情回暖产品利润率回升、同比扭亏为盈；植提类业务收……；受原材料成本下降影响，多数植物提取产品价格同比下降，行业价格竞争明显。\footnote{\texttt{https://static.cninfo.com.cn/finalpage/2026-02-27/1224988258.PDF}}\footnote{\texttt{https://money.finance.sina.com.cn/corp/view/vCB\_AllBulletinDetail.php}}

\begin{center}\small\setlength{\tabcolsep}{4pt}
\captionof{table}[晨光生物（300138）关键财务指标]{晨光生物（300138）关键财务指标（合并报表口径；两个半年报期均未年化；现金短债比 $=$ 货币资金 $\div$ 短期债务）}
\begin{tabular}{@{}lrrrrr@{}}
\toprule
指标 & 2023 年 & 2024 年 & 2025 年 & 2025 年半年报 & 2026 年半年报 \\
\midrule
营业收入（亿元） & 68.72 & 69.94 & 65.59 & 36.58 & 31.59 \\
归母净利润（亿元） & 4.80 & 0.94 & 3.69 & 2.15 & 1.79 \\
毛利率（\%） & 11.6 & 8.0 & 13.1 & 13.9 & 13.5 \\
ROE（\%） & 14.3 & 2.8 & 11.4 & 6.7 & 5.2 \\
资产负债率（\%） & 58.6 & 65.0 & 63.8 & 64.4 & 64.1 \\
经营活动现金流净额（亿元） & -5.88 & -2.23 & 2.13 & 18.30 & 12.18 \\
货币资金（亿元） & 14.22 & 18.02 & 14.59 & 36.20 & 28.02 \\
有息负债合计（亿元） & 42.36 & 49.31 & 51.23 & 53.01 & 53.17 \\
现金短债比（倍） & 0.42 & 0.47 & 0.29 & 0.71 & 0.53 \\
\bottomrule
\end{tabular}
\end{center}

\begin{center}
\begin{tikzpicture}
\begin{groupplot}[
  group style={group size=2 by 1, horizontal sep=0.95cm},
  width=7.2cm, height=4.6cm,
  tick label style={font=\tiny},
  title style={font=\scriptsize\heiti},
  yticklabel style={font=\tiny},
  grid=major, grid style={LightGray, line width=0.3pt},
  axis line style={InkGray, line width=0.5pt},
  enlarge x limits={abs=0.8},
  legend style={font=\scriptsize, at={(0.5,-0.22)}, anchor=north,
                legend columns=2, draw=none, fill=none, column sep=6pt},
]
\nextgroupplot[
  ybar, bar width=4pt,
  ymin=-6.99, ymax=78.33,
  xtick={0,1,2,3,4,5.7},
  xticklabels={2021,2022,2023,2024,2025,2026H1},
  ylabel={亿元},
]
\addplot[fill=AgriGreen!75, draw=AgriGreen, bar shift=-2.6pt] table[x=i, y=rev]{data/clean/profiles/fin_300138.dat};
\addplot[fill=SoftRed!60, draw=SoftRed, bar shift=2.6pt] table[x=i, y=ni]{data/clean/profiles/fin_300138.dat};
\addplot[fill=AgriGreen!30, draw=AgriGreen, pattern=north east lines, bar shift=-2.6pt] table[x=x, y=rev]{data/clean/profiles/finh_300138.dat};
\addplot[fill=SoftRed!20, draw=SoftRed, pattern=north east lines, bar shift=2.6pt] table[x=x, y=ni]{data/clean/profiles/finh_300138.dat};
\legend{营业收入（年/半年）, 归母净利润（年/半年）}
\nextgroupplot[
  ymin=-5.2, ymax=71.5,
  xtick={0,1,2,3,4,5.7},
  xticklabels={2021,2022,2023,2024,2025,2026H1},
  ylabel={\%},
]
\addplot[AgriGold, mark=*, mark size=1.3pt] table[x=i, y=roe]{data/clean/profiles/fin_300138.dat};
\addplot[OilBlue, mark=square*, mark size=1.3pt] table[x=i, y=debt]{data/clean/profiles/fin_300138.dat};
\addplot[only marks, AgriGold, mark=*, mark size=2.0pt] table[x=x, y=roe]{data/clean/profiles/finh_300138.dat};
\addplot[only marks, OilBlue, mark=square*, mark size=2.0pt] table[x=x, y=debt]{data/clean/profiles/finh_300138.dat};
\legend{ROE, 资产负债率}
\end{groupplot}
\end{tikzpicture}
\begin{tikzpicture}
\begin{axis}[
  width=15.2cm, height=3.1cm, scale only axis,
  ymin=-11.79, ymax=60.26,
  xtick={0,1,2,3,4,5.7},
  xticklabels={2021,2022,2023,2024,2025,2026H1},
  tick label style={font=\tiny},
  yticklabel style={font=\tiny},
  ylabel={亿元}, ylabel style={font=\scriptsize},
  ybar, bar width=4pt,
  grid=major, grid style={LightGray, line width=0.3pt},
  axis line style={InkGray, line width=0.5pt},
  enlarge x limits={abs=0.8},
  legend style={font=\scriptsize, at={(0.5,-0.30)}, anchor=north,
                legend columns=3, draw=none, fill=none, column sep=10pt},
]
\addplot[fill=AgriGreen!55, draw=AgriGreen, bar shift=-3.0pt] table[x=i, y=cash]{data/clean/profiles/fin_300138.dat};
\addplot[fill=SoftRed!45, draw=SoftRed, bar shift=3.0pt] table[x=i, y=ibd]{data/clean/profiles/fin_300138.dat};
\addplot[fill=AgriGreen!25, draw=AgriGreen, pattern=north east lines, bar shift=-3.0pt] table[x=x, y=cash]{data/clean/profiles/finh_300138.dat};
\addplot[fill=SoftRed!20, draw=SoftRed, pattern=north east lines, bar shift=3.0pt] table[x=x, y=ibd]{data/clean/profiles/finh_300138.dat};
\addplot[OilBlue, mark=*, mark size=1.1pt, line width=0.8pt] table[x=i, y=ocf]{data/clean/profiles/fin_300138.dat};
\addplot[only marks, OilBlue, mark=*, mark size=1.8pt] table[x=x, y=ocf]{data/clean/profiles/finh_300138.dat};
\legend{货币资金, 有息负债, 经营活动现金流净额}
\end{axis}
\end{tikzpicture}

\captionof{figure}[晨光生物（300138）近年主要财务指标]{晨光生物（300138）近年主要财务指标（左上：营业收入与归母净利润；右上：ROE 与资产负债率；下：货币资金、有息负债与经营活动现金流净额。
2021---2025 年为年报口径，2026H1 为 2026 年半年报/半年末，与其前一年报之间留出间隔、以斜纹或空心点区分；半年报数据未年化）}
\end{center}

\pfl{财务分析} 2026 年 6 月末资产负债率 64.1\%，较 2025 年末 上升 0.2 个百分点（样本中位数 52.2\%，所处三级行业内排第 \zhnumber{10} 位）。 2026 年 6 月末货币资金 28.02 亿元、短期债务（短期借款与一年内到期非流动负债）52.71 亿元、长期债务 0.46 亿元，有息负债合计 53.17 亿元，现金短债比 0.53 倍——覆盖偏紧。 净有息负债 25.16 亿元，相当于其 2026 年上半年营业收入的 80\%。 流动比率 1.28、速动比率 0.73，短期指标尚可。 2026 年上半年经营活动现金流净额 12.18 亿元，经营现金流与利润同向为正，内生资金可以支撑运营。 2026 年 6 月末在建工程 0.94 亿元，较上年末增加 0.21 亿元，仍有在建投入。

\pfl{重大战略转型} 从公开披露看，公司的战略推进包括：2023-2025 推进“营养药用综合提取项目（一期）”与“赞比亚土地开发及配套设施建设项目”，在境外布局原料基地并进行万寿菊深加工；2025-06-30 植提类产品收入增长9.47\%、毛利额同比增长32\%，毛利率同比提高3.43个百分点，盈利能力显著恢复；2025-12-31 2025年全年实现收入65.59亿元，归属于上市公司股东的净利润3.69亿元，同比增长292.21\%；扣非净利润……；2025-12-31 产品从单一辣椒红色素发展到天然色素、天然香辛料提取物和精油、天然甜味剂、天然营养及药用提取物、保健食品、油脂和蛋……。 公告口径上，近三年（2023 年 9 月---2026 年 9 月）公司披露重大重组与并购类公告 0 条、对外投资与转型类公告 1 条、回购与股权激励类公告 45 条，合计公告 614 条。 第一大行业仍是“天然植物提取行业”，收入占比 100.0\%，与 2021 年年报相比没有方向性变化。 综合判断，公司在报告期内有明确的外延扩张或业务结构调整动作，转型方向与融资安排之间存在直接关联。\footnote{\texttt{http://static.cninfo.com.cn/finalpage/2025-08-26/1224566487.PDF}}\footnote{\texttt{https://money.finance.sina.com.cn/corp/view/vCB\_AllBulletinDetail.php}}\footnote{\texttt{http://xinpi.zqrb.cn/pdf/CYB/300138/2026/0331/a5fe1f3e559a7209d36a469f741ea913.pdf}}

\begin{center}\small\setlength{\tabcolsep}{4pt}
\captionof{table}[晨光生物（300138）历史融资与资本运作]{晨光生物（300138）历史融资与资本运作（据公司公告与公开报道整理；金额为披露值，含首次公开发行、再融资、债券、银行借款与重整投资等）}
\begin{tabular}{@{}p{1.5cm}p{3.2cm}p{2.4cm}p{6.3cm}@{}}
\toprule
时间 & 方式 & 规模 & 用途与说明 \\
\midrule
2010-11-05 & IPO（深交所创业板） & 发行2,300万股，发行价30.00元/股，募集资金净额64,908.51万元 & 招股说明书披露的募投项目；证监许可[2010]1411号核准…… \\
2020-06-17 & 公开发行可转换公司债券（创业板） & 面值总额6.30亿元（630万张），募集资金总额6.30亿元，扣除发行费用10,955,660.36元后净额619,044,339.64元 & 天然植物综合提取一体化项目（一期）45,000.00万元…… \\
2021-02-24 & 可转债提前赎回与转股 & 2020年12月23日至停止交易日共6,273,758张转股，转股价12.25元/股；未转股26,242张按100.35元/张赎回 & 转股增加股本，赎回剩余债券；2021年1月13日董事会通过提前赎回议案…… \\
2023-03-28 & 对全资子公司增资（自有资金） & “赞比亚土地开发及配套设施建设项目”建设总投资1,500万美元 & 配套增建菊花加工颗粒车间、萃取车间、菊花存储池，用于自产万寿菊深加工…… \\
2023-05-29 & 变更部分可转债募集资金用途 & 将“天然植物综合提取一体化项目（一期）”部分募集资金6,500万元变更用于“营养药用综合提取项目（一期）”，1,100万美元（按2023年5月29日汇率7.0575折算为7,763.25万元人民币）变更用于“赞比亚土地开发及配套设施建设项目” & 营养药用综合提取项目（一期）、赞比亚土地开发及配套设施建设项目…… \\
2024-06-30 & 子公司首次公开发行终止 & 未查得具体金额 & 控股子公司新疆晨光生物科技股份有限公司决定终止向不特定合格投资者公开发行并在北交所上市…… \\
\bottomrule
\end{tabular}
\end{center}

\pfl{历史融投资情况} 从现金流量表看，2025 年公司取得借款收到的现金 59.76 亿元、偿还债务支付的现金 57.66 亿元、吸收权益性投资 — 亿元，筹资活动现金流净额 0.69 亿元（2023 年为取得借款 46.78 亿元、偿还债务 33.80 亿元）。 2026 年上半年筹资活动现金流净额为 -0.42 亿元（取得借款 27.24 亿元、偿还债务 25.28 亿元，未年化），已转为净流出，处于净还债状态。 2025 年分配股利、利润或偿付利息支付的现金合计 1.60 亿元。 公开信息中与公司直接相关的融资与资本运作包括：IPO（深交所创业板）、公开发行可转换公司债券（创业板）、变更部分可转债募集资金用途、可转债提前赎回与转股、子公司首次公开发行终止、对全资子公司增资（自有资金）。 其中规模较大的两笔为：2023-05-29变更部分可转债募集资金用途（将“天然植物综合提取一体化项目（一期）”部分募集资金6,50……）；2020-06-17公开发行可转换公司债券（创业板）（面值总额6.30亿元（630万张），募集资金总额6.30亿元……）。 股本方面，近三年披露股本变动 5 次（其他 1 次、股份回购 1 次）；总股本由 2021 年末的 5.33 亿股变为 4.83 亿股（-9.32\%）。 分红方面，近三年现金分红 4 次、累计每股派现 0.72 元（最近一次 2025-12-31）——在利润承压的阶段仍维持了分红。\footnote{\texttt{http://file.finance.sina.com.cn/211.154.219.97:9494/MRGG/CNSESZ\_STOCK/2011/2011-3/2011-0}}\footnote{\texttt{http://file.finance.sina.com.cn/211.154.219.97:9494/MRGG/CNSESZ\_STOCK/2022/2022-3/2022-0}}\footnote{\texttt{https://money.finance.sina.com.cn/corp/view/vCB\_AllBulletinDetail.php}}

\pfl{未来潜在融资需求分析} 2026 年 6 月末货币资金 28.02 亿元，而短期债务 52.71 亿元（现金短债比 0.53 倍），2026 年上半年经营现金流净额 12.18 亿元。按“短期债务减去账面货币资金”的滚动偿付口径，静态缺口约 24.70 亿元；若再计入上半年亏损的年化额（0.00 亿元），维持一年正常经营所需的外部资金约 24.70---24.70 亿元。 按第 \ref{sec:financing-need} 节的三档划分，公司属于第三档“扩张型”：2026 年上半年 ROE 5.2\%（未年化，折合约 10.3\%）、6 月末资产负债率 64.1\%，资产负债表仍具备加杠杆空间；融资用途以产能扩张、品类并购与海外布局为主，单笔规模通常在 5---20 亿元量级，会被市场定价为成长而非补血。 公开披露的下一步融资线索包括：2025-12-31 可转债募集资金已累计使用38,867.84万元（截至2023年12月31日口径披露），后续项目建设主要依靠自有资金（未查……；2026-08-25 2026年上半年不派发现金红利、不送红股、不以公积金转增股本，资金优先用于经营与项目投入。 资金需求还要叠加在建工程：期末在建工程 0.94 亿元、2026 年上半年购建固定资产等支出 0.60 亿元（未年化），这部分支出在周期底部通常会被主动放缓。\footnote{\texttt{https://stockn.xueqiu.com/SZ300138/20240417682702.pdf}}\footnote{\texttt{http://static.cninfo.com.cn/finalpage/2026-08-25/1225495892.PDF}}


\subsubsection{华资实业（600191）}
\label{co:600191}

\pfl{公司基本情况} 1998 年设立至今，华资实业包头起家的甜菜制糖老牌上市公司，曾是全国甜菜制糖行业唯一上市公司，长期依靠参股华夏银行、恒泰证券（现金融街证券）的投资收益支撑业绩；近年主业转为谷朊粉、食用酒精、黄原胶等粮食深加工与生物发酵，2025年由山东中裕食品通过增资控股股东间接入主。 公司于 1998-12-10 在上交所首发上市，注册资本 4.85 亿元，总股本 4.85 亿股。 截至 2026-09-24收盘价 9.06 元，流通市值 43.9 亿元，近一年+14.5\%，流通市值在三级行业“其他农产品加工”的 11 家公司中排第 \zhnumber{6} 位，近一年涨跌幅高于全样本中位数（-10.1\%）24.6 个百分点。 按 2025 年年报披露的行业构成，收入占比前三位为食品制造业 91.3\%；农副食品加工业 7.7\%；其他(补充) 1.0\%。 与 2021 年年报相比，第一大行业由“制糖业”变为“食品制造业”，主业重心已经迁移。 发展过程中的关键节点包括：经证监发字（1998）276、277号文批准向社会公开发行A股7,000万股，1998年……（1998）；以上市为契机，“草原”商标成为我国驰名商标；肖建华控制的包头创业、包头北普分别为第二、第……（1998）；包头市国资委将控股股东包头草原糖业（集团）100\%国有产权在内蒙古产权交易中心整体挂牌转……（2012）。\footnote{\texttt{https://basic.10jqka.com.cn/600191/company.html}}\footnote{\texttt{https://finance.sina.cn/stock/ssgs/2018-01-22/detail-ifyqwiqi1640400.d.html}}\footnote{\texttt{http://pdf.dfcfw.com/pdf/H2\_AN201211250005615016\_1.pdf}}

\pfl{历史股价} 本报告样本区间（2021 年 1 月 至 2026 年 9 月）内，股价由 4.82 元变为 9.06 元，累计 +88.0\%，区间最高 13.86 元（2026 年 1 月）、最低 2.72 元（2021 年 5 月）；当前价相当于区间最高价的 65.4\%，为区间最低价的 3.33 倍。 月度收益的年化波动率 44.3\%，高于全样本中位数 37.4\%，在三级行业“其他农产品加工”的 11 家公司中波动率排第 \zhnumber{2} 位。 从事件看，2024 年 8 月 至 2026 年 1 月的上涨（+188.1\%）与公司控股股东获得13亿元现金增资用于取得控制权，未披露对上市公司的资金注入安排（2025-04-15）在时间上重合；2021 年 5 月 至 2022 年 8 月的上涨（+182.4\%）与公司撤销退市风险警示恢复“华资实业”；同期“明天系”资本版图收缩（包商银行破产清……（2022-08-16）在时间上重合——股价的拐点多数能在公司自身的资本运作与风险事件上找到对应。\footnote{\texttt{https://www.stcn.com/article/detail/1665420.html}}\footnote{\texttt{https://www.sohu.com/a/562574227\_616823}}

\begin{center}
\begin{tikzpicture}
\begin{axis}[
  width=12.6cm, height=3.9cm, scale only axis,
  xmin=2020.922, xmax=2026.888, ymin=2.53, ymax=14.83,
  xtick={2021,2022,2023,2024,2025,2026},
  yearticks,
  yticklabel style={font=\scriptsize},
  ylabel={元/股}, ylabel style={font=\scriptsize},
  grid=major, grid style={LightGray, line width=0.3pt},
  axis line style={InkGray, line width=0.5pt},
  every axis plot/.append style={line width=0.9pt},
]
\addplot[AgriGreen] table[x=x,y=y]{data/clean/profiles/px_600191.dat};
\addplot[SoftRed, dashed, line width=0.7pt] table[x=x,y=y]{data/clean/profiles/pxzz_600191.dat};
\addplot[only marks, mark=*, mark size=1.1pt, SoftRed] table[x=x,y=y]{data/clean/profiles/pxzz_600191.dat};
\end{axis}
\end{tikzpicture}

\captionof{figure}[华资实业（600191）股价与周期骨架]{华资实业（600191）月度收盘价（元/股）与 ZigZag 周期骨架（阈值 25\%，圆点为周期转折点）}
\end{center}

\pfl{过去周期分析} 以 25\% 的 ZigZag 阈值划分，样本区间内股价形成 3 个主升段与 4 个主跌段。 最大主升段出现在 2024 年 8 月 至 2026 年 1 月，17 个月+188.1\%；最大主跌段为 2021 年 1 月 至 2021 年 5 月，4 个月-43.6\%。 最近一次转折出现在 2026 年 9 月（下跌段自 2026 年 1 月起，8 个月-34.6\%），当前处于该段的延续之中。 公司股价较申万农林牧渔指数提前约 60 个月见底，是板块中较早定价周期反转的公司之一。 该子行业横跨糖、番茄制品、添加剂等多个品种，各自的供给周期与政策（进口配额、关税）是主要变量，公司 2026 年上半年实现盈利、半年 ROE 2.1\%（未年化），好于全样本中位数（0.75\%），下行周期对其报表的冲击小于同业。

\pfl{盈利情况} 2026 年上半年营业收入 3.49 亿元（同比 +58.0\%），归母净利润 0.36 亿元，同比 +722.6\%。 以年报为趋势对照，2023---2025 年营业收入由 6.37 亿元变为 6.29 亿元（两年年均复合增速 -0.6\%），归母净利润由 0.19 亿元变为 0.59 亿元。 按半年报口径，2026 年上半年的 ROE 为 2.1\%（未年化）、销售净利率 10.4\%、毛利率 4.6\%，ROE 在“其他农产品加工”11 家公司中排第 \zhnumber{4} 位。 同期全样本半年 ROE 中位数 0.75\%，公司与之相差 1.37 个百分点（略好于中位数公司）。 结合公司披露，2025年度母公司存在未弥补亏损的相关情况（对分红等事项产生影响，公司已提示）。\footnote{\texttt{https://static.cninfo.com.cn/finalpage/2026-04-21/1225126275.PDF}}

\begin{center}\small\setlength{\tabcolsep}{4pt}
\captionof{table}[华资实业（600191）关键财务指标]{华资实业（600191）关键财务指标（合并报表口径；两个半年报期均未年化；现金短债比 $=$ 货币资金 $\div$ 短期债务）}
\begin{tabular}{@{}lrrrrr@{}}
\toprule
指标 & 2023 年 & 2024 年 & 2025 年 & 2025 年半年报 & 2026 年半年报 \\
\midrule
营业收入（亿元） & 6.37 & 5.16 & 6.29 & 2.21 & 3.49 \\
归母净利润（亿元） & 0.19 & 0.25 & 0.59 & 0.04 & 0.36 \\
毛利率（\%） & 8.6 & 8.3 & 7.5 & 6.0 & 4.6 \\
ROE（\%） & 1.3 & 1.6 & 3.5 & 0.3 & 2.1 \\
资产负债率（\%） & 20.6 & 25.5 & 27.7 & 27.1 & 23.8 \\
经营活动现金流净额（亿元） & 1.10 & 1.39 & -0.06 & 0.19 & 0.31 \\
货币资金（亿元） & 0.47 & 0.42 & 0.90 & 0.46 & 0.46 \\
有息负债合计（亿元） & 2.70 & 1.68 & 2.68 & 2.93 & 2.53 \\
现金短债比（倍） & 0.27 & 0.50 & 0.59 & 0.28 & 0.33 \\
\bottomrule
\end{tabular}
\end{center}

\begin{center}
\begin{tikzpicture}
\begin{groupplot}[
  group style={group size=2 by 1, horizontal sep=0.95cm},
  width=7.2cm, height=4.6cm,
  tick label style={font=\tiny},
  title style={font=\scriptsize\heiti},
  yticklabel style={font=\tiny},
  grid=major, grid style={LightGray, line width=0.3pt},
  axis line style={InkGray, line width=0.5pt},
  enlarge x limits={abs=0.8},
  legend style={font=\scriptsize, at={(0.5,-0.22)}, anchor=north,
                legend columns=2, draw=none, fill=none, column sep=6pt},
]
\nextgroupplot[
  ybar, bar width=4pt,
  ymin=-2.45, ymax=7.34,
  xtick={0,1,2,3,4,5.7},
  xticklabels={2021,2022,2023,2024,2025,2026H1},
  ylabel={亿元},
]
\addplot[fill=AgriGreen!75, draw=AgriGreen, bar shift=-2.6pt] table[x=i, y=rev]{data/clean/profiles/fin_600191.dat};
\addplot[fill=SoftRed!60, draw=SoftRed, bar shift=2.6pt] table[x=i, y=ni]{data/clean/profiles/fin_600191.dat};
\addplot[fill=AgriGreen!30, draw=AgriGreen, pattern=north east lines, bar shift=-2.6pt] table[x=x, y=rev]{data/clean/profiles/finh_600191.dat};
\addplot[fill=SoftRed!20, draw=SoftRed, pattern=north east lines, bar shift=2.6pt] table[x=x, y=ni]{data/clean/profiles/finh_600191.dat};
\legend{营业收入（年/半年）, 归母净利润（年/半年）}
\nextgroupplot[
  ymin=-13.5, ymax=31.5,
  xtick={0,1,2,3,4,5.7},
  xticklabels={2021,2022,2023,2024,2025,2026H1},
  ylabel={\%},
]
\addplot[AgriGold, mark=*, mark size=1.3pt] table[x=i, y=roe]{data/clean/profiles/fin_600191.dat};
\addplot[OilBlue, mark=square*, mark size=1.3pt] table[x=i, y=debt]{data/clean/profiles/fin_600191.dat};
\addplot[only marks, AgriGold, mark=*, mark size=2.0pt] table[x=x, y=roe]{data/clean/profiles/finh_600191.dat};
\addplot[only marks, OilBlue, mark=square*, mark size=2.0pt] table[x=x, y=debt]{data/clean/profiles/finh_600191.dat};
\legend{ROE, 资产负债率}
\end{groupplot}
\end{tikzpicture}
\begin{tikzpicture}
\begin{axis}[
  width=15.2cm, height=3.1cm, scale only axis,
  ymin=-2.76, ymax=3.51,
  xtick={0,1,2,3,4,5.7},
  xticklabels={2021,2022,2023,2024,2025,2026H1},
  tick label style={font=\tiny},
  yticklabel style={font=\tiny},
  ylabel={亿元}, ylabel style={font=\scriptsize},
  ybar, bar width=4pt,
  grid=major, grid style={LightGray, line width=0.3pt},
  axis line style={InkGray, line width=0.5pt},
  enlarge x limits={abs=0.8},
  legend style={font=\scriptsize, at={(0.5,-0.30)}, anchor=north,
                legend columns=3, draw=none, fill=none, column sep=10pt},
]
\addplot[fill=AgriGreen!55, draw=AgriGreen, bar shift=-3.0pt] table[x=i, y=cash]{data/clean/profiles/fin_600191.dat};
\addplot[fill=SoftRed!45, draw=SoftRed, bar shift=3.0pt] table[x=i, y=ibd]{data/clean/profiles/fin_600191.dat};
\addplot[fill=AgriGreen!25, draw=AgriGreen, pattern=north east lines, bar shift=-3.0pt] table[x=x, y=cash]{data/clean/profiles/finh_600191.dat};
\addplot[fill=SoftRed!20, draw=SoftRed, pattern=north east lines, bar shift=3.0pt] table[x=x, y=ibd]{data/clean/profiles/finh_600191.dat};
\addplot[OilBlue, mark=*, mark size=1.1pt, line width=0.8pt] table[x=i, y=ocf]{data/clean/profiles/fin_600191.dat};
\addplot[only marks, OilBlue, mark=*, mark size=1.8pt] table[x=x, y=ocf]{data/clean/profiles/finh_600191.dat};
\legend{货币资金, 有息负债, 经营活动现金流净额}
\end{axis}
\end{tikzpicture}

\captionof{figure}[华资实业（600191）近年主要财务指标]{华资实业（600191）近年主要财务指标（左上：营业收入与归母净利润；右上：ROE 与资产负债率；下：货币资金、有息负债与经营活动现金流净额。
2021---2025 年为年报口径，2026H1 为 2026 年半年报/半年末，与其前一年报之间留出间隔、以斜纹或空心点区分；半年报数据未年化）}
\end{center}

\pfl{财务分析} 2026 年 6 月末资产负债率 23.8\%，较 2025 年末 下降 3.9 个百分点（样本中位数 52.2\%，所处三级行业内排第 \zhnumber{4} 位）。 2026 年 6 月末货币资金 0.46 亿元、短期债务（短期借款与一年内到期非流动负债）1.38 亿元、长期债务 1.14 亿元，有息负债合计 2.53 亿元，现金短债比 0.33 倍——覆盖偏紧。 净有息负债 2.07 亿元，相当于其 2026 年上半年营业收入的 59\%。 流动比率 0.61、速动比率 0.36，短期偿付压力较大。 2026 年上半年经营活动现金流净额 0.31 亿元，经营现金流与利润同向为正，内生资金可以支撑运营。 公司披露的财务与合规风险包括：撤销退市风险警示恢复“华资实业”；同期“明天系”资本版图收缩（包商银行破产清算、新时代信……；因2020年度归母净利润-12,530.24万元、营业收入3,863.94万元低于1亿元……。\footnote{\texttt{https://www.sohu.com/a/562574227\_616823}}\footnote{\texttt{https://pdf.dfcfw.com/pdf/H2\_AN202204291562353761\_1.pdf}}

\pfl{重大战略转型} 近三年公司的战略动作集中在以下几个方面：2025-04 年产2万吨黄原胶项目正式投产（2025年4月30日披露《关于黄原胶项目投产的公告》，编号临2025-023），2……；2025-04 中裕食品入主后从管理、产品、渠道、中长期等维度对公司赋能，形成“中裕赋能”的业务协同逻辑；2025-12-31 与中科院环境生物技术重点实验室共建合作平台；自治区首家微生物改造技术领域院士工作站落地公司；2026-06-30 2026年上半年实现营业收入34,918.85万元、同比增长58.00\%，归属于上市公司股东的净利润3,616.……。 公告口径上，近三年（2023 年 9 月---2026 年 9 月）公司披露重大重组与并购类公告 2 条、对外投资与转型类公告 2 条、回购与股权激励类公告 0 条，合计公告 316 条。 主营结构的迁移已经落到报表上：2021 年年报的第一大行业为“制糖业”（84.2\%），2025 年年报变为“食品制造业”（91.3\%）。 综合判断，公司当前的战略重心不是扩张而是化解债务与完成重整，业务层面的调整让位于司法重整程序。\footnote{\texttt{https://finance.sina.com.cn/stock/aiassist/ggsp/2026-08-31/doc-iniqexup7840667.shtml}}\footnote{\texttt{http://file.finance.sina.com.cn/211.154.219.97:9494/MRGG/CNSESH\_STOCK/2025/2025-8/2025-0}}

\begin{center}\small\setlength{\tabcolsep}{4pt}
\captionof{table}[华资实业（600191）历史融资与资本运作]{华资实业（600191）历史融资与资本运作（据公司公告与公开报道整理；金额为披露值，含首次公开发行、再融资、债券、银行借款与重整投资等）}
\begin{tabular}{@{}p{1.5cm}p{3.2cm}p{2.4cm}p{6.3cm}@{}}
\toprule
时间 & 方式 & 规模 & 用途与说明 \\
\midrule
1998-11-02 & 首发A股（上交所） & 发行7,000万股，发行价4.30元/股，实际募资3.01亿元；发行市盈率13.20倍，首日开盘价7.80元 & 新浪财经报道口径为上市募集资金2.89亿元…… \\
2015-2016 & 非公开发行股份募集资金（拟用于增资华夏人寿） & 拟募集不超过316.8亿元，增资完成后持有华夏人寿股权比例不超过51\% & 对华夏人寿增资；2016年7月证监会发审委审核通过…… \\
2025-04-15 & 控股股东层面增资（间接控股，非上市公司融资） & 中裕科技投资以现金130,000万元对盛泰创发增资，取得增资后55\%股权，通过盛泰创发间接持有上市公司1.45亿股、占公司总股本29.90\% & 取得上市公司控制权；控股股东仍为盛泰创发，中裕科技投资未直接持有上市公司股份…… \\
\bottomrule
\end{tabular}
\end{center}

\pfl{历史融投资情况} 从现金流量表看，2025 年公司取得借款收到的现金 1.60 亿元、偿还债务支付的现金 0.44 亿元、吸收权益性投资 — 亿元，筹资活动现金流净额 2.36 亿元（2023 年为取得借款 0.70 亿元、偿还债务 0.79 亿元）。 2026 年上半年筹资活动现金流净额为 -0.58 亿元（取得借款 0.88 亿元、偿还债务 0.98 亿元，未年化），已转为净流出，处于净还债状态。 2025 年分配股利、利润或偿付利息支付的现金合计 0.07 亿元。 公开信息中与公司直接相关的融资与资本运作包括：控股股东层面增资（间接控股，非上市公司融资）、非公开发行股份募集资金（拟用于增资华夏人寿）、首发A股（上交所）。 其中规模较大的两笔为：2015-2016非公开发行股份募集资金（拟用于增资华夏人寿）（拟募集不超过316.8亿元，增资完成后持有华夏人寿股权比例不……）；1998-11-02首发A股（上交所）（发行7,000万股，发行价4.30元/股，实际募资3.01亿……）。 股本方面，近三年披露股本变动 3 次；总股本由 2021 年末的 4.85 亿股变为 4.85 亿股（+0.00\%）。 分红方面，近三年现金分红 1 次、累计每股派现 0.02 元（最近一次 2025-12-31）——分红规模与利润的下滑幅度相比仍然有限。 近三年“再融资”类公告 7 条，涉及定向增发、募集资金使用。\footnote{\texttt{https://basic.10jqka.com.cn/600191/company.html}}\footnote{\texttt{https://www.chinawznews.net/news/china/2017-02-06/19418.html}}\footnote{\texttt{https://www.stcn.com/article/detail/1665420.html}}

\pfl{未来潜在融资需求分析} 2026 年 6 月末货币资金 0.46 亿元，而短期债务 1.38 亿元（现金短债比 0.33 倍），2026 年上半年经营现金流净额 0.31 亿元。按“短期债务减去账面货币资金”的滚动偿付口径，静态缺口约 0.92 亿元；若再计入上半年亏损的年化额（0.00 亿元），维持一年正常经营所需的外部资金约 0.92---0.92 亿元。 公司 2026 年 6 月末资产负债率 23.8\%、2026 年上半年 ROE 2.1\%（未年化），资产负债表尚有余量，融资需求不是“补血型”的刚性需求，而是围绕并购整合、项目投入与营运资金的主动性需求。 公开披露的下一步融资线索包括：2025-04-15 控股股东获得13亿元现金增资用于取得控制权，未披露对上市公司的资金注入安排；2026-08-29 2026年半年度无利润分配预案或公积金转增股本预案（未查得公开的定增、发债或授信额度信息）。\footnote{\texttt{https://www.stcn.com/article/detail/1665420.html}}\footnote{\texttt{https://finance.sina.cn/7x24/2026-08-28/detail-inipwhrr8976025.d.html}}


\subsubsection{京粮控股（000505）}
\label{co:000505}

\pfl{公司基本情况} 京粮控股成立于 1992 年，北京市国资委系统（首农食品集团—京粮集团）旗下油脂油料加工与食品加工上市平台，由原海南珠江控股2017年重大资产置换并注入京粮食品而来；油脂贸易体量大但毛利薄，2025年因贸易收缩与原料成本上升转为大额亏损。 公司于 1992-12-21 在深交所主板首发上市，注册资本 7.27 亿元，总股本 7.27 亿股。 截至 2026-09-24收盘价 6.36 元，流通市值 42.1 亿元，近一年-1.7\%，流通市值在三级行业“粮油加工”的 7 家公司中排第 \zhnumber{3} 位，近一年涨跌幅高于全样本中位数（-10.1\%）8.4 个百分点。 按 2026 年半年报披露的行业构成，收入占比前三位为油脂 89.2\%；食品 10.5\%；其他(补充) 0.3\%。 发展过程中的关键节点包括：1992年12月21日在深圳证券交易所上市，证券代码000505，B股代码200505；……（1992）；1992年至1999年控股股东为广州珠江实业集团有限公司；1999年至2016年控股股东……（1992-2016）；2017年7月28日取得证监会《关于核准海南珠江控股股份有限公司向北京粮食集团有限责任公……（2017）。\footnote{\texttt{https://q.stock.sohu.com/cn/000505/index.shtml}}\footnote{\texttt{https://pdf.dfcfw.com/pdf/H2\_AN202203301556090886\_1.pdf}}\footnote{\texttt{https://money.finance.sina.com.cn/corp/view/vCB\_AllBulletinDetail.php}}

\pfl{历史股价} 自 2021 年 1 月至 2026 年 9 月，股价由 7.12 元变为 6.36 元，累计 -10.7\%，区间最高 12.31 元（2022 年 3 月）、最低 4.97 元（2024 年 8 月）；当前价相当于区间最高价的 51.7\%，为区间最低价的 1.28 倍。 月度收益的年化波动率 39.9\%，高于全样本中位数 37.4\%，在三级行业“粮油加工”的 7 家公司中波动率排第 \zhnumber{2} 位。 从事件看，2024 年 8 月 至 2025 年 12 月的上涨（+71.2\%）与公司2025年营业收入78.59亿元，同比下降31.28\%；归属于上市公司股东的……（2025-12-31）在时间上重合——股价的拐点多数能在公司自身的资本运作与风险事件上找到对应。\footnote{\texttt{https://finance.sina.com.cn/stock/aigc/stockfs/2026-03-27/doc-inhsmnzv8306228.shtml}}

\begin{center}
\begin{tikzpicture}
\begin{axis}[
  width=12.6cm, height=3.9cm, scale only axis,
  xmin=2020.922, xmax=2026.888, ymin=4.62, ymax=13.17,
  xtick={2021,2022,2023,2024,2025,2026},
  yearticks,
  yticklabel style={font=\scriptsize},
  ylabel={元/股}, ylabel style={font=\scriptsize},
  grid=major, grid style={LightGray, line width=0.3pt},
  axis line style={InkGray, line width=0.5pt},
  every axis plot/.append style={line width=0.9pt},
]
\addplot[AgriGreen] table[x=x,y=y]{data/clean/profiles/px_000505.dat};
\addplot[SoftRed, dashed, line width=0.7pt] table[x=x,y=y]{data/clean/profiles/pxzz_000505.dat};
\addplot[only marks, mark=*, mark size=1.1pt, SoftRed] table[x=x,y=y]{data/clean/profiles/pxzz_000505.dat};
\end{axis}
\end{tikzpicture}

\captionof{figure}[京粮控股（000505）股价与周期骨架]{京粮控股（000505）月度收盘价（元/股）与 ZigZag 周期骨架（阈值 25\%，圆点为周期转折点）}
\end{center}

\pfl{过去周期分析} 以 25\% 的 ZigZag 阈值划分，样本区间内股价形成 3 个主升段与 2 个主跌段。 最大主跌段为 2022 年 3 月 至 2024 年 8 月，29 个月-59.6\%；最大主升段出现在 2021 年 1 月 至 2022 年 3 月，14 个月+72.9\%。 最近一次转折出现在 2026 年 9 月（上涨段自 2026 年 6 月起，3 个月+26.4\%），当前处于该段之后的震荡之中。 公司股价较申万农林牧渔指数提前约 21 个月见底，是板块中较早定价周期反转的公司之一。 粮油加工是“成本加成”型生意（第 \ref{sec:grain} 章），利润来自原料与成品的价差及规模效率，收入规模大而净利率薄，公司 2026 年 6 月末资产负债率 43.5\%、半年 ROE 0.3\%（未年化），在本轮下行中属于随行业同步调整的一类。

\pfl{盈利情况} 2026 年上半年营业收入 31.15 亿元（同比 -26.0\%），归母净利润 0.08 亿元，同比 -55.3\%。 以年报为趋势对照，2023---2025 年营业收入由 116.02 亿元变为 78.59 亿元（两年年均复合增速 -17.7\%），归母净利润由 1.01 亿元变为 -2.66 亿元。 按半年报口径，2026 年上半年的 ROE 为 0.3\%（未年化）、销售净利率 0.3\%、毛利率 5.6\%，ROE 在“粮油加工”7 家公司中排第 \zhnumber{6} 位。 同期全样本半年 ROE 中位数 0.75\%，公司与之相差 0.47 个百分点（弱于中位数公司）。 结合公司披露，2026年上半年营业收入31.15亿元，同比下降25.98\%；归属于上市公司股东的净利润803.24万元……；净利润下降主要受子公司新增食品加工生产线尚处建设期及休闲食品线下渠道受低价竞争冲击等因素影响；报告期内品……。\footnote{\texttt{https://stock.stockstar.com/notice/SN2026082600041982.shtml}}\footnote{\texttt{https://static.weeklyonstock.com/26/0826/ABB2616093891967.html}}

\begin{center}\small\setlength{\tabcolsep}{4pt}
\captionof{table}[京粮控股（000505）关键财务指标]{京粮控股（000505）关键财务指标（合并报表口径；两个半年报期均未年化；现金短债比 $=$ 货币资金 $\div$ 短期债务）}
\begin{tabular}{@{}lrrrrr@{}}
\toprule
指标 & 2023 年 & 2024 年 & 2025 年 & 2025 年半年报 & 2026 年半年报 \\
\midrule
营业收入（亿元） & 116.02 & 114.35 & 78.59 & 42.08 & 31.15 \\
归母净利润（亿元） & 1.01 & 0.26 & -2.66 & 0.18 & 0.08 \\
毛利率（\%） & 3.0 & 4.5 & 4.5 & 5.6 & 5.6 \\
ROE（\%） & 3.3 & 0.8 & -8.9 & 0.6 & 0.3 \\
资产负债率（\%） & 45.0 & 47.6 & 47.5 & 53.2 & 43.5 \\
经营活动现金流净额（亿元） & 1.09 & -1.09 & 6.47 & 2.22 & 3.90 \\
货币资金（亿元） & 15.43 & 14.17 & 18.22 & 17.61 & 16.90 \\
有息负债合计（亿元） & 21.11 & 22.05 & 22.25 & 23.49 & 17.91 \\
现金短债比（倍） & 1.15 & 0.76 & 1.20 & 1.19 & 1.55 \\
\bottomrule
\end{tabular}
\end{center}

\begin{center}
\begin{tikzpicture}
\begin{groupplot}[
  group style={group size=2 by 1, horizontal sep=0.95cm},
  width=7.2cm, height=4.6cm,
  tick label style={font=\tiny},
  title style={font=\scriptsize\heiti},
  yticklabel style={font=\tiny},
  grid=major, grid style={LightGray, line width=0.3pt},
  axis line style={InkGray, line width=0.5pt},
  enlarge x limits={abs=0.8},
  legend style={font=\scriptsize, at={(0.5,-0.22)}, anchor=north,
                legend columns=2, draw=none, fill=none, column sep=6pt},
]
\nextgroupplot[
  ybar, bar width=4pt,
  ymin=-15.78, ymax=144.33,
  xtick={0,1,2,3,4,5.7},
  xticklabels={2021,2022,2023,2024,2025,2026H1},
  ylabel={亿元},
]
\addplot[fill=AgriGreen!75, draw=AgriGreen, bar shift=-2.6pt] table[x=i, y=rev]{data/clean/profiles/fin_000505.dat};
\addplot[fill=SoftRed!60, draw=SoftRed, bar shift=2.6pt] table[x=i, y=ni]{data/clean/profiles/fin_000505.dat};
\addplot[fill=AgriGreen!30, draw=AgriGreen, pattern=north east lines, bar shift=-2.6pt] table[x=x, y=rev]{data/clean/profiles/finh_000505.dat};
\addplot[fill=SoftRed!20, draw=SoftRed, pattern=north east lines, bar shift=2.6pt] table[x=x, y=ni]{data/clean/profiles/finh_000505.dat};
\legend{营业收入（年/半年）, 归母净利润（年/半年）}
\nextgroupplot[
  ymin=-13.4, ymax=53.3,
  xtick={0,1,2,3,4,5.7},
  xticklabels={2021,2022,2023,2024,2025,2026H1},
  ylabel={\%},
]
\addplot[AgriGold, mark=*, mark size=1.3pt] table[x=i, y=roe]{data/clean/profiles/fin_000505.dat};
\addplot[OilBlue, mark=square*, mark size=1.3pt] table[x=i, y=debt]{data/clean/profiles/fin_000505.dat};
\addplot[only marks, AgriGold, mark=*, mark size=2.0pt] table[x=x, y=roe]{data/clean/profiles/finh_000505.dat};
\addplot[only marks, OilBlue, mark=square*, mark size=2.0pt] table[x=x, y=debt]{data/clean/profiles/finh_000505.dat};
\legend{ROE, 资产负债率}
\end{groupplot}
\end{tikzpicture}
\begin{tikzpicture}
\begin{axis}[
  width=15.2cm, height=3.1cm, scale only axis,
  ymin=-8.09, ymax=25.55,
  xtick={0,1,2,3,4,5.7},
  xticklabels={2021,2022,2023,2024,2025,2026H1},
  tick label style={font=\tiny},
  yticklabel style={font=\tiny},
  ylabel={亿元}, ylabel style={font=\scriptsize},
  ybar, bar width=4pt,
  grid=major, grid style={LightGray, line width=0.3pt},
  axis line style={InkGray, line width=0.5pt},
  enlarge x limits={abs=0.8},
  legend style={font=\scriptsize, at={(0.5,-0.30)}, anchor=north,
                legend columns=3, draw=none, fill=none, column sep=10pt},
]
\addplot[fill=AgriGreen!55, draw=AgriGreen, bar shift=-3.0pt] table[x=i, y=cash]{data/clean/profiles/fin_000505.dat};
\addplot[fill=SoftRed!45, draw=SoftRed, bar shift=3.0pt] table[x=i, y=ibd]{data/clean/profiles/fin_000505.dat};
\addplot[fill=AgriGreen!25, draw=AgriGreen, pattern=north east lines, bar shift=-3.0pt] table[x=x, y=cash]{data/clean/profiles/finh_000505.dat};
\addplot[fill=SoftRed!20, draw=SoftRed, pattern=north east lines, bar shift=3.0pt] table[x=x, y=ibd]{data/clean/profiles/finh_000505.dat};
\addplot[OilBlue, mark=*, mark size=1.1pt, line width=0.8pt] table[x=i, y=ocf]{data/clean/profiles/fin_000505.dat};
\addplot[only marks, OilBlue, mark=*, mark size=1.8pt] table[x=x, y=ocf]{data/clean/profiles/finh_000505.dat};
\legend{货币资金, 有息负债, 经营活动现金流净额}
\end{axis}
\end{tikzpicture}

\captionof{figure}[京粮控股（000505）近年主要财务指标]{京粮控股（000505）近年主要财务指标（左上：营业收入与归母净利润；右上：ROE 与资产负债率；下：货币资金、有息负债与经营活动现金流净额。
2021---2025 年为年报口径，2026H1 为 2026 年半年报/半年末，与其前一年报之间留出间隔、以斜纹或空心点区分；半年报数据未年化）}
\end{center}

\pfl{财务分析} 2026 年 6 月末资产负债率 43.5\%，较 2025 年末 下降 4.0 个百分点（样本中位数 52.2\%，所处三级行业内排第 \zhnumber{4} 位）。 2026 年 6 月末货币资金 16.90 亿元、短期债务（短期借款与一年内到期非流动负债）10.88 亿元、长期债务 7.03 亿元，有息负债合计 17.91 亿元，现金短债比 1.55 倍——覆盖充分。 净有息负债 1.01 亿元，相当于其 2026 年上半年营业收入的 3\%。 流动比率 2.26、速动比率 1.32，短期指标尚可。 2026 年上半年经营活动现金流净额 3.90 亿元，经营现金流与利润同向为正，内生资金可以支撑运营。 2026 年 6 月末在建工程 1.31 亿元，较上年末增加 0.42 亿元，仍有在建投入。 2026 年 6 月末商誉 1.26 亿元，占归母净资产的 3.9\%，是净资产中最脆弱的一块。

\pfl{重大战略转型} 从公开披露看，公司的战略推进包括：2025-12-31 2025年度主动优化业务结构、推动贸易业务减量发展，有序压降贸易业务规模；营业收入下降主要因贸易业务收缩；2026-06-30 2026年6月30日加权平均净资产收益率0.28\%（上年同期0.57\%）；2026-08-26 2026年上半年营业收入31.15亿元，同比下降25.98\%；归属于上市公司股东的净利润803.24万元，同比下……；2026-08-26 净利润下降主要受子公司新增食品加工生产线尚处建设期及休闲食品线下渠道受低价竞争冲击等因素影响；报告期内品牌建设成……。 公告口径上，近三年（2023 年 9 月---2026 年 9 月）公司披露重大重组与并购类公告 2 条、对外投资与转型类公告 0 条、回购与股权激励类公告 0 条，合计公告 335 条。 第一大行业仍是“油脂”，收入占比 89.2\%，与 2021 年年报相比没有方向性变化。 综合判断，公司在报告期内有明确的外延扩张或业务结构调整动作，转型方向与融资安排之间存在直接关联。\footnote{\texttt{https://static.weeklyonstock.com/26/0826/ABB2616093891967.html}}\footnote{\texttt{https://stock.stockstar.com/notice/SN2026082600041982.shtml}}

\begin{center}\small\setlength{\tabcolsep}{4pt}
\captionof{table}[京粮控股（000505）历史融资与资本运作]{京粮控股（000505）历史融资与资本运作（据公司公告与公开报道整理；金额为披露值，含首次公开发行、再融资、债券、银行借款与重整投资等）}
\begin{tabular}{@{}p{1.5cm}p{3.2cm}p{2.4cm}p{6.3cm}@{}}
\toprule
时间 & 方式 & 规模 & 用途与说明 \\
\midrule
2017 & 发行股份募集配套资金 & 非公开发行不超过48,965,408股，发行价8.82元/股，募集配套资金总额不超过43,187.49万元；实际募集资金总额431,874,898.56元，扣除发行费用45,000,000.00元 & 配套资金（招股/重组方案约定用途）；股东京粮集团以货币资金认购 \\
2017-09-14 & 重大资产置换＋发行股份购买资产（借壳上市） & 京粮股份（北京京粮食品有限公司）100\%股权作价230,852.72万元，置出资产作价60,898.36万元，交易价格差额169,954.36万元；公司发行210,079,552股、发行价8.09元/股向京粮食品原股东发行股份购买差额部分；同时向京粮集团非公开发行48,965,408股新股募集配套资金 & 资产置换、购买京粮食品100\%股权、配套资金；新增股份上市日期为2017年11月15日…… \\
2023-08-21 & 公开发行公司债券（第一期） & 发行总额3亿元（30,000万元），期限3年，票面利率2.88\%，债券简称“23京粮01”，代码148434.SZ & 偿还到期债务和补充公司流动资金；起息日2023-08-22，兑付日2026-08-22…… \\
2023-11-21 & 变更募集资金用途 & 未查得具体金额 & 变更募集资金用途（2023年11月22日披露《关于变更募集资金用途的公告》…… \\
\bottomrule
\end{tabular}
\end{center}

\pfl{历史融投资情况} 从现金流量表看，2025 年公司取得借款收到的现金 41.83 亿元、偿还债务支付的现金 42.42 亿元、吸收权益性投资 0.01 亿元，筹资活动现金流净额 -1.79 亿元（2023 年为取得借款 33.59 亿元、偿还债务 27.96 亿元）。 2026 年上半年筹资活动现金流净额为 -4.57 亿元（取得借款 7.24 亿元、偿还债务 11.54 亿元，未年化），已转为净流出，处于净还债状态。 2025 年分配股利、利润或偿付利息支付的现金合计 0.73 亿元。 公开信息中与公司直接相关的融资与资本运作包括：公开发行公司债券（第一期）、发行股份募集配套资金、变更募集资金用途、重大资产置换＋发行股份购买资产（借壳上市）。 其中规模较大的两笔为：2023-08-21公开发行公司债券（第一期）（发行总额3亿元（30,000万元），期限3年，票面利率2.8……）；2017发行股份募集配套资金（非公开发行不超过48,965,408股，发行价8.82元/股……）。 股本方面，近三年披露股本变动 4 次（限售股份上市 1 次）；总股本由 2021 年末的 7.27 亿股变为 7.27 亿股（+0.00\%）。 分红方面，近三年现金分红 2 次、累计每股派现 0.09 元（最近一次 2024-12-31）——在利润承压的阶段仍维持了分红。\footnote{\texttt{http://file.finance.sina.com.cn/211.154.219.97:9494/MRGG/CNSESZ\_STOCK/2024/2024-8/2024-0}}\footnote{\texttt{https://money.finance.sina.com.cn/corp/view/vCB\_AllBulletinDetail.php}}\footnote{\texttt{https://disc.static.szse.cn/disc/disk03/finalpage/2024-06-28/e443ce3a-ccf2-4f81-9554-685}}

\pfl{未来潜在融资需求分析} 2026 年 6 月末货币资金 16.90 亿元，而短期债务 10.88 亿元（现金短债比 1.55 倍），2026 年上半年经营现金流净额 3.90 亿元。账面货币资金可以覆盖短期债务，缺口不体现在静态偿付层面。 公司 2026 年 6 月末资产负债率 43.5\%、2026 年上半年 ROE 0.3\%（未年化），资产负债表尚有余量，融资需求不是“补血型”的刚性需求，而是围绕并购整合、项目投入与营运资金的主动性需求。 公开披露的下一步融资线索包括：2026-05-27 联合资信出具2026年跟踪评级报告，跟踪期内“23京粮01”已在付息日正常付息，募集资金使用符合募集说明书约定用途；2026-08-22 “23京粮01”（余额3亿元，票面利率2.88\%）于2026年8月22日到期兑付，存在债券接续融资需求；2026-08-26 子公司新增食品加工生产线仍在建设期，2026年上半年经营活动现金流净额3.90亿元，后续仍需资本开支（未查得公开的定增预……。 资金需求还要叠加在建工程：期末在建工程 1.31 亿元、2026 年上半年购建固定资产等支出 0.76 亿元（未年化），这部分支出在周期底部通常会被主动放缓。\footnote{\texttt{https://disc.static.szse.cn/disc/disk03/finalpage/2026-05-27/f5d98f71-40f3-49f4-9e03-4b0}}\footnote{\texttt{https://static.weeklyonstock.com/26/0826/ABB2616093891967.html}}


\subsubsection{东方海洋（002086）}
\label{co:002086}

\pfl{公司基本情况} 2001 年设立至今，东方海洋山东烟台的海水苗种繁育与养殖、水产品加工企业，2006年深交所上市；2022—2023年因债务危机、资金占用与违规担保进入破产重整，2023年底重整计划执行完毕、2024年6月摘帽，现推进“海洋产业＋大健康（体外诊断）”双主业，但仍持续亏损。 公司于 2006-11-28 在深交所主板首发上市，注册资本 19.59 亿元，总股本 19.59 亿股。 截至 2026-09-24收盘价 2.04 元，流通市值 32.3 亿元，近一年-22.7\%，流通市值在三级行业“其他农产品加工”的 11 家公司中排第 \zhnumber{7} 位，近一年涨跌幅低于全样本中位数（-10.1\%）12.6 个百分点。 按 2026 年半年报披露的行业构成，收入占比前三位为水产品加工 68.4\%；体外诊断 12.9\%；海水养殖 12.7\%。 发展过程中的关键节点包括：经证监发行字[2006]109号文核准，公开发行人民币普通股3,450万股（网下配售69……（2006）；2022年8月8日债权人烟台市春达工贸有限公司以不能清偿到期债务且明显缺乏清偿能力为由向……（2022）；2023年11月25日烟台市中级人民法院裁定受理北京汉业科技有限公司对公司的破产重整申请……（2023）。\footnote{\texttt{https://vip.stock.finance.sina.com.cn/corp/view/vISSUE\_MarketBulletinDetail.php}}\footnote{\texttt{http://xinpi.zqrb.cn/pdf/SZ/002086/2023/1217/7d0f7cd5b53d7627a36408838e476595.pdf}}\footnote{\texttt{http://notice.10jqka.com.cn/api/pdf/c347f66ab65caabf.pdf}}

\pfl{历史股价} 以 2021 年 1 月为起点、2026 年 9 月为终点，股价由 1.66 元变为 2.04 元，累计 +22.9\%，区间最高 3.51 元（2024 年 11 月）、最低 1.58 元（2021 年 4 月）；当前价相当于区间最高价的 58.1\%，为区间最低价的 1.29 倍。 月度收益的年化波动率 43.8\%，高于全样本中位数 37.4\%，在三级行业“其他农产品加工”的 11 家公司中波动率排第 \zhnumber{3} 位。 从事件看，2021 年 1 月 至 2021 年 12 月的上涨（+85.5\%）与公司截至2021年12月31日，控股股东及其关联方资金占用、违规担保、违规使用募……（2021-12-31）在时间上重合——股价的拐点多数能在公司自身的资本运作与风险事件上找到对应。\footnote{\texttt{http://money.finance.sina.com.cn/corp/view/vCB\_AllBulletinDetail.php}}

\begin{center}
\begin{tikzpicture}
\begin{axis}[
  width=12.6cm, height=3.9cm, scale only axis,
  xmin=2020.922, xmax=2026.888, ymin=1.47, ymax=3.76,
  xtick={2021,2022,2023,2024,2025,2026},
  yearticks,
  yticklabel style={font=\scriptsize},
  ylabel={元/股}, ylabel style={font=\scriptsize},
  grid=major, grid style={LightGray, line width=0.3pt},
  axis line style={InkGray, line width=0.5pt},
  every axis plot/.append style={line width=0.9pt},
]
\addplot[AgriGreen] table[x=x,y=y]{data/clean/profiles/px_002086.dat};
\addplot[SoftRed, dashed, line width=0.7pt] table[x=x,y=y]{data/clean/profiles/pxzz_002086.dat};
\addplot[only marks, mark=*, mark size=1.1pt, SoftRed] table[x=x,y=y]{data/clean/profiles/pxzz_002086.dat};
\end{axis}
\end{tikzpicture}

\captionof{figure}[东方海洋（002086）股价与周期骨架]{东方海洋（002086）月度收盘价（元/股）与 ZigZag 周期骨架（阈值 25\%，圆点为周期转折点）}
\end{center}

\pfl{过去周期分析} 以 25\% 的 ZigZag 阈值划分，样本区间内股价形成 3 个主升段与 3 个主跌段。 最大主跌段为 2024 年 11 月 至 2026 年 6 月，19 个月-53.3\%；最大主升段出现在 2021 年 1 月 至 2021 年 12 月，11 个月+85.5\%。 最近一次转折出现在 2026 年 6 月（下跌段自 2024 年 11 月起，19 个月-53.3\%），当前处于该段的延续之中。 公司股价较申万农林牧渔指数提前约 61 个月见底，是板块中较早定价周期反转的公司之一。 该子行业横跨糖、番茄制品、添加剂等多个品种，各自的供给周期与政策（进口配额、关税）是主要变量，公司 2026 年 6 月末资产负债率 19.6\%、半年 ROE -3.8\%（未年化），在本轮下行中属于随行业同步调整的一类。

\pfl{盈利情况} 2026 年上半年营业收入 1.57 亿元（同比 +0.9\%），归母净利润 -0.48 亿元，亏损同比收窄 0.23 亿元。 以年报为趋势对照，2023---2025 年营业收入由 4.37 亿元变为 3.56 亿元（两年年均复合增速 -9.8\%），归母净利润由 17.50 亿元变为 -1.53 亿元。 按半年报口径，2026 年上半年的 ROE 为 -3.8\%（未年化）、销售净利率 -30.7\%、毛利率 0.2\%，ROE 在“其他农产品加工”11 家公司中排第 \zhnumber{11} 位。 行业同口径半年 ROE 的中位数为 0.75\%，该公司低于中位数 4.59 个百分点。 结合公司披露，2026年上半年营业收入157,397,250.94元、同比增长0.87\%；归属于上市公司股东的净利润-……；2025年度业绩预告/年报提及优质新资产注入是重整后仍面临的挑战之一，公司表示将通过外延并购及内生增长方……。\footnote{\texttt{https://finance.sina.com.cn/stock/zqgd/2026-08-25/doc-inipprnk9668806.shtml}}\footnote{\texttt{https://www.stcn.com/article/detail/1204232.html}}

\begin{center}\small\setlength{\tabcolsep}{4pt}
\captionof{table}[东方海洋（002086）关键财务指标]{东方海洋（002086）关键财务指标（合并报表口径；两个半年报期均未年化；现金短债比 $=$ 货币资金 $\div$ 短期债务）}
\begin{tabular}{@{}lrrrrr@{}}
\toprule
指标 & 2023 年 & 2024 年 & 2025 年 & 2025 年半年报 & 2026 年半年报 \\
\midrule
营业收入（亿元） & 4.37 & 3.40 & 3.56 & 1.56 & 1.57 \\
归母净利润（亿元） & 17.50 & -1.90 & -1.53 & -0.71 & -0.48 \\
毛利率（\%） & 13.7 & -3.4 & 0.2 & -0.9 & 0.2 \\
ROE（\%） & 0.0 & -12.5 & -11.3 & -5.1 & -3.8 \\
资产负债率（\%） & 39.9 & 20.9 & 20.4 & 18.5 & 19.6 \\
经营活动现金流净额（亿元） & -0.46 & -1.93 & -1.25 & -0.62 & 0.08 \\
货币资金（亿元） & 12.93 & 4.05 & 2.35 & 3.03 & 2.43 \\
有息负债合计（亿元） & 5.26 & 0.63 & 0.53 & 0.63 & 0.49 \\
现金短债比（倍） & 3.17 & 27.25 & 21.44 & 21.92 & 37.76 \\
\bottomrule
\end{tabular}
\end{center}

\begin{center}
\begin{tikzpicture}
\begin{groupplot}[
  group style={group size=2 by 1, horizontal sep=0.95cm},
  width=7.2cm, height=4.6cm,
  tick label style={font=\tiny},
  title style={font=\scriptsize\heiti},
  yticklabel style={font=\tiny},
  grid=major, grid style={LightGray, line width=0.3pt},
  axis line style={InkGray, line width=0.5pt},
  enlarge x limits={abs=0.8},
  legend style={font=\scriptsize, at={(0.5,-0.22)}, anchor=north,
                legend columns=2, draw=none, fill=none, column sep=6pt},
]
\nextgroupplot[
  ybar, bar width=4pt,
  ymin=-19.18, ymax=21.50,
  xtick={0,1,2,3,4,5.7},
  xticklabels={2021,2022,2023,2024,2025,2026H1},
  ylabel={亿元},
]
\addplot[fill=AgriGreen!75, draw=AgriGreen, bar shift=-2.6pt] table[x=i, y=rev]{data/clean/profiles/fin_002086.dat};
\addplot[fill=SoftRed!60, draw=SoftRed, bar shift=2.6pt] table[x=i, y=ni]{data/clean/profiles/fin_002086.dat};
\addplot[fill=AgriGreen!30, draw=AgriGreen, pattern=north east lines, bar shift=-2.6pt] table[x=x, y=rev]{data/clean/profiles/finh_002086.dat};
\addplot[fill=SoftRed!20, draw=SoftRed, pattern=north east lines, bar shift=2.6pt] table[x=x, y=ni]{data/clean/profiles/finh_002086.dat};
\legend{营业收入（年/半年）, 归母净利润（年/半年）}
\nextgroupplot[
  ymin=-132.1, ymax=184.5,
  xtick={0,1,2,3,4,5.7},
  xticklabels={2021,2022,2023,2024,2025,2026H1},
  ylabel={\%},
]
\addplot[AgriGold, mark=*, mark size=1.3pt] table[x=i, y=roe]{data/clean/profiles/fin_002086.dat};
\addplot[OilBlue, mark=square*, mark size=1.3pt] table[x=i, y=debt]{data/clean/profiles/fin_002086.dat};
\addplot[only marks, AgriGold, mark=*, mark size=2.0pt] table[x=x, y=roe]{data/clean/profiles/finh_002086.dat};
\addplot[only marks, OilBlue, mark=square*, mark size=2.0pt] table[x=x, y=debt]{data/clean/profiles/finh_002086.dat};
\legend{ROE, 资产负债率}
\end{groupplot}
\end{tikzpicture}
\begin{tikzpicture}
\begin{axis}[
  width=15.2cm, height=3.1cm, scale only axis,
  ymin=-3.42, ymax=14.71,
  xtick={0,1,2,3,4,5.7},
  xticklabels={2021,2022,2023,2024,2025,2026H1},
  tick label style={font=\tiny},
  yticklabel style={font=\tiny},
  ylabel={亿元}, ylabel style={font=\scriptsize},
  ybar, bar width=4pt,
  grid=major, grid style={LightGray, line width=0.3pt},
  axis line style={InkGray, line width=0.5pt},
  enlarge x limits={abs=0.8},
  legend style={font=\scriptsize, at={(0.5,-0.30)}, anchor=north,
                legend columns=3, draw=none, fill=none, column sep=10pt},
]
\addplot[fill=AgriGreen!55, draw=AgriGreen, bar shift=-3.0pt] table[x=i, y=cash]{data/clean/profiles/fin_002086.dat};
\addplot[fill=SoftRed!45, draw=SoftRed, bar shift=3.0pt] table[x=i, y=ibd]{data/clean/profiles/fin_002086.dat};
\addplot[fill=AgriGreen!25, draw=AgriGreen, pattern=north east lines, bar shift=-3.0pt] table[x=x, y=cash]{data/clean/profiles/finh_002086.dat};
\addplot[fill=SoftRed!20, draw=SoftRed, pattern=north east lines, bar shift=3.0pt] table[x=x, y=ibd]{data/clean/profiles/finh_002086.dat};
\addplot[OilBlue, mark=*, mark size=1.1pt, line width=0.8pt] table[x=i, y=ocf]{data/clean/profiles/fin_002086.dat};
\addplot[only marks, OilBlue, mark=*, mark size=1.8pt] table[x=x, y=ocf]{data/clean/profiles/finh_002086.dat};
\legend{货币资金, 有息负债, 经营活动现金流净额}
\end{axis}
\end{tikzpicture}

\captionof{figure}[东方海洋（002086）近年主要财务指标]{东方海洋（002086）近年主要财务指标（左上：营业收入与归母净利润；右上：ROE 与资产负债率；下：货币资金、有息负债与经营活动现金流净额。
2021---2025 年为年报口径，2026H1 为 2026 年半年报/半年末，与其前一年报之间留出间隔、以斜纹或空心点区分；半年报数据未年化）}
\end{center}

\pfl{财务分析} 2026 年 6 月末资产负债率 19.6\%，较 2025 年末 下降 0.8 个百分点（样本中位数 52.2\%，所处三级行业内排第 \zhnumber{1} 位）。 2026 年 6 月末货币资金 2.43 亿元、短期债务（短期借款与一年内到期非流动负债）0.06 亿元、长期债务 0.42 亿元，有息负债合计 0.49 亿元，现金短债比 37.76 倍——覆盖充分。 净有息负债 -1.94 亿元，相当于其 2026 年上半年营业收入的 -123\%。 流动比率 3.71、速动比率 2.05，短期指标尚可。 2026 年上半年经营活动现金流净额 0.08 亿元，净利润为负而经营现金流为正，差额主要来自折旧摊销与营运资本变动，说明亏损尚未完全转化为现金流出。 公司披露的财务与合规风险包括：因2022年净资产为负值和审计意见无法表示两项原因被实施退市风险警示；同时因非经营性资金……；2022年度经审计的期末净资产为负值；2022年度财务会计报告被出具无法表示意见的审计报……。\footnote{\texttt{https://www.stcn.com/article/detail/1204232.html}}\footnote{\texttt{http://money.finance.sina.com.cn/corp/view/vCB\_AllBulletinDetail.php}}

\pfl{重大战略转型} 近三年公司的战略动作集中在以下几个方面：2024-2025 依托重整投资人提供的流动资金支持和业务资源支持，一方面做实做精做强海洋产业，另一方面以高端医疗和精准医疗为导向培……；2025-12-31 2025年度业绩预告/年报提及优质新资产注入是重整后仍面临的挑战之一，公司表示将通过外延并购及内生增长方式发展壮……；2026-03-31 2026年一季度营业收入7,837.01万元、同比增长4.11\%，经营活动现金流净额597.11万元、同比增长1……；2026-06-30 2026年上半年继续推进海洋产业和大健康产业发展，加快产业转型、升级，实现企业双轮驱动的可持续发展；在医疗器械注……。 公告口径上，近三年（2023 年 9 月---2026 年 9 月）公司披露重大重组与并购类公告 0 条、对外投资与转型类公告 4 条、回购与股权激励类公告 0 条，合计公告 463 条。 第一大行业“水产品加工”的收入占比由 2021 年年报的 35.5\% 变为 68.4\%（32.9 个百分点），收入结构出现再平衡。 综合判断，公司当前的战略重心不是扩张而是化解债务与完成重整，业务层面的调整让位于司法重整程序。\footnote{\texttt{https://vip.stock.finance.sina.com.cn/corp/view/vCB\_AllBulletinDetail.php}}\footnote{\texttt{https://www.stcn.com/article/detail/1204232.html}}

\begin{center}\small\setlength{\tabcolsep}{4pt}
\captionof{table}[东方海洋（002086）历史融资与资本运作]{东方海洋（002086）历史融资与资本运作（据公司公告与公开报道整理；金额为披露值，含首次公开发行、再融资、债券、银行借款与重整投资等）}
\begin{tabular}{@{}p{1.5cm}p{3.2cm}p{2.4cm}p{6.3cm}@{}}
\toprule
时间 & 方式 & 规模 & 用途与说明 \\
\midrule
2006-11-28 & IPO（深交所中小板） & 发行3,450万股，发行价7.49元/股，实际募集资金合计25,840.50万元，发行费用1,408.37万元，募集资金净额24,432.13万元 & 全部用于水产品加工贸易基地建设项目；发行市盈率（按发行后总股本）23.77倍…… \\
2023-12 & 破产重整投资款（资本公积金转增股票引入重整投资人） & 截至2023年12月29日重整投资人向管理人指定银行账户共计支付13.54亿元投资款 & 偿付债务、支付重整费用、补充公司流动资金、部分解决非经营性资金占用及违规担保问题…… \\
2023-12 & 债务豁免与债务重组 & 债权人直接豁免东方海洋清偿义务256,217,732.62元；豁免金额相应抵偿原控股股东对上市公司同等金额的非经营性资金占用776,907,937.55元，剩余缺口由重整投资人以现金支付至管理人账户 & 解决原控股股东非经营性资金占用及违规担保；部分债权人与原控股股东等签订《债务重组与化解协…… \\
2023-12-17 & 重整投资协议 & 五矿金通股权投资基金管理有限公司牵头设立专项私募投资基金，联合其他投资人共同以重整投资人身份参与重整投资；联合体包含烟台东晨投资、中国长城资产管理股份有限公司山东省分公司、深圳市招平京资二号投资合伙企业（招商平安）、青岛兴锐投资合伙企业、嘉兴鑫扬、绿叶投资集团等 & 认购公司资本公积转增股本，帮助公司化解债务风险…… \\
\bottomrule
\end{tabular}
\end{center}

\pfl{历史融投资情况} 从现金流量表看，2025 年公司取得借款收到的现金 — 亿元、偿还债务支付的现金 — 亿元、吸收权益性投资 0.71 亿元，筹资活动现金流净额 0.67 亿元（2023 年为取得借款 0.45 亿元、偿还债务 0.66 亿元）。 2026 年上半年筹资活动现金流净额为 -0.00 亿元（取得借款 — 亿元、偿还债务 — 亿元，未年化），已转为净流出，处于净还债状态。 公开信息中与公司直接相关的融资与资本运作包括：IPO（深交所中小板）、债务豁免与债务重组、破产重整投资款（资本公积金转增股票引入重整投资人）、重整投资协议。 其中规模较大的两笔为：2023-12破产重整投资款（资本公积金转增股票引入重整投资人）（截至2023年12月29日重整投资人向管理人指定银行账户共计……）；2006-11-28IPO（深交所中小板）（发行3,450万股，发行价7.49元/股，实际募集资金合计2……）。 股本方面，近三年披露股本变动 5 次（转增 1 次、限售股份上市 1 次）；总股本由 2021 年末的 7.56 亿股变为 19.59 亿股（+159.00\%）。 分红方面，近三年未实施现金分红，与重整期间的资金约束一致。\footnote{\texttt{https://vip.stock.finance.sina.com.cn/corp/go.php/vISSUE\_NewStock/stockid/002086.phtml}}\footnote{\texttt{https://money.finance.sina.com.cn/corp/view/vCB\_AllBulletinDetail.php}}\footnote{\texttt{https://disc.static.szse.cn/download/disc/disk03/finalpage/2024-04-18/93d18291-e227-425c}}

\pfl{未来潜在融资需求分析} 2026 年 6 月末货币资金 2.43 亿元，而短期债务 0.06 亿元（现金短债比 37.76 倍），2026 年上半年经营现金流净额 0.08 亿元。账面货币资金可以覆盖短期债务，缺口不体现在静态偿付层面。 公司 2026 年上半年归母净利润为负（-0.48 亿元），但 6 月末资产负债率 19.6\% 尚未进入第一档（亏损且负债率 $>$65\%）的区间，当前融资需求以补充流动资金、支撑低谷期经营性支出为主。 公开披露的下一步融资线索包括：2023-12 重整投资人合计投入13.54亿元资金，解决非经营性资金占用、违规担保及债务清偿后剩余资金全部用于补充流动资金；2026-06-30 2025年财务费用大幅上升、体外诊断收入大幅下滑，持续经营仍依赖重整投资人的资本与业务资源支持；未查得公开的定增预案、发……。\footnote{\texttt{https://i.ifeng.com/c/8aLoY1uns30}}\footnote{\texttt{http://finance.sina.com.cn/stock/aigc/stockfs/2026-04-29/doc-inhwckrq6802415.shtml}}


\subsubsection{保龄宝（002286）}
\label{co:002286}

\pfl{公司基本情况} 山东禹城的功能糖（益生元、膳食纤维、减糖甜味剂）龙头，国内首家完成赤藓糖醇工业化生产的企业，为元气森林等无糖饮品品牌的核心代糖供应商；2009年上市，控股股东永裕投资持股比例不高且近年发生股份司法冻结，公司正谋求以简易程序小额再融资。 公司于 2009-08-28 在深交所主板首发上市，注册资本 3.80 亿元，总股本 3.70 亿股。 截至 2026-09-24收盘价 8.46 元，流通市值 31.3 亿元，近一年-13.2\%，流通市值在三级行业“其他农产品加工”的 11 家公司中排第 \zhnumber{8} 位，近一年涨跌幅低于全样本中位数（-10.1\%）3.1 个百分点。 按 2026 年半年报披露的行业构成，收入占比前三位为农副食品加工业 100.0\%。 与 2021 年年报相比，第一大行业由“食品制造业”变为“农副食品加工业”，主业重心已经迁移。 发展过程中的关键节点包括：2013年4月完成非公开发行，新增股份于2013年4月17日上市；发行49,428,00……（2013）；控股股东及实际控制人变更：控股股东由刘宗利变更为北京永裕投资管理有限公司，实际控制人由刘……（2017）；2021年12月完成2021年股票期权与限制性股票激励计划之限制性股票的首次授予登记工作……（2021）。\footnote{\texttt{http://money.finance.sina.com.cn/corp/view/vISSUE\_MarketBulletinDetail.php}}\footnote{\texttt{http://pdf.dfcfw.com/pdf/H2\_AN201701200282746951\_01.pdf}}\footnote{\texttt{http://file.finance.sina.com.cn/211.154.219.97:9494/MRGG/CNSESZ\_STOCK/2022/2022-4/2022-0}}

\pfl{历史股价} 自 2021 年 1 月至 2026 年 9 月，股价由 8.42 元变为 8.46 元，累计 +0.5\%，区间最高 14.66 元（2021 年 11 月）、最低 5.51 元（2024 年 6 月）；当前价相当于区间最高价的 57.7\%，为区间最低价的 1.54 倍。 月度收益的年化波动率 35.3\%，低于全样本中位数 37.4\%，在三级行业“其他农产品加工”的 11 家公司中波动率排第 \zhnumber{5} 位。 从事件看，2021 年 1 月 至 2021 年 11 月的上涨（+74.1\%）与公司发行对象为控股股东北京永裕投资管理有限公司、赣州中科永裕投资管理中心（有限合……（2021-07-13）在时间上重合——股价的拐点多数能在公司自身的资本运作与风险事件上找到对应。\footnote{\texttt{https://file.finance.qq.com/finance/hs/pdf/2022/09/30/1214717674.PDF}}

\begin{center}
\begin{tikzpicture}
\begin{axis}[
  width=12.6cm, height=3.9cm, scale only axis,
  xmin=2020.922, xmax=2026.888, ymin=5.12, ymax=15.69,
  xtick={2021,2022,2023,2024,2025,2026},
  yearticks,
  yticklabel style={font=\scriptsize},
  ylabel={元/股}, ylabel style={font=\scriptsize},
  grid=major, grid style={LightGray, line width=0.3pt},
  axis line style={InkGray, line width=0.5pt},
  every axis plot/.append style={line width=0.9pt},
]
\addplot[AgriGreen] table[x=x,y=y]{data/clean/profiles/px_002286.dat};
\addplot[SoftRed, dashed, line width=0.7pt] table[x=x,y=y]{data/clean/profiles/pxzz_002286.dat};
\addplot[only marks, mark=*, mark size=1.1pt, SoftRed] table[x=x,y=y]{data/clean/profiles/pxzz_002286.dat};
\end{axis}
\end{tikzpicture}

\captionof{figure}[保龄宝（002286）股价与周期骨架]{保龄宝（002286）月度收盘价（元/股）与 ZigZag 周期骨架（阈值 25\%，圆点为周期转折点）}
\end{center}

\pfl{过去周期分析} 以 25\% 的 ZigZag 阈值划分，样本区间内股价形成 2 个主升段与 2 个主跌段。 最大主升段出现在 2024 年 6 月 至 2025 年 8 月，14 个月+111.6\%；最大主跌段为 2021 年 11 月 至 2024 年 6 月，31 个月-62.4\%。 最近一次转折出现在 2026 年 7 月（下跌段自 2025 年 8 月起，11 个月-32.0\%），当前处于该段的延续之中。 公司股价较申万农林牧渔指数提前约 23 个月见底，是板块中较早定价周期反转的公司之一。 该子行业横跨糖、番茄制品、添加剂等多个品种，各自的供给周期与政策（进口配额、关税）是主要变量，公司 2026 年上半年实现盈利、半年 ROE 3.6\%（未年化），好于全样本中位数（0.75\%），下行周期对其报表的冲击小于同业。

\pfl{盈利情况} 2026 年上半年营业收入 14.96 亿元（同比 +6.9\%），归母净利润 0.82 亿元，同比 -11.2\%。 以年报为趋势对照，2023---2025 年营业收入由 25.24 亿元变为 27.51 亿元（两年年均复合增速 +4.4\%），归母净利润由 0.54 亿元变为 1.51 亿元。 按半年报口径，2026 年上半年的 ROE 为 3.6\%（未年化）、销售净利率 5.5\%、毛利率 14.7\%，ROE 在“其他农产品加工”11 家公司中排第 \zhnumber{3} 位。 同期全样本半年 ROE 中位数 0.75\%，公司与之相差 2.87 个百分点（略好于中位数公司）。 结合公司披露，2026年上半年营业收入14.96亿元、同比增长6.92\%；归属于上市公司股东的净利润8,232.11万……；2025年膳食纤维产品收入同比下降9.85\%，淀粉糖及其他产品收入同比下降4.94\%，产品结构分化明显。\footnote{\texttt{https://paper.cnstock.com/html/2026-08/27/content\_2260948.htm}}\footnote{\texttt{https://finance.sina.com.cn/stock/aigc/stockfs/2026-04-24/doc-inhvramy4654904.shtml}}

\begin{center}\small\setlength{\tabcolsep}{4pt}
\captionof{table}[保龄宝（002286）关键财务指标]{保龄宝（002286）关键财务指标（合并报表口径；两个半年报期均未年化；现金短债比 $=$ 货币资金 $\div$ 短期债务）}
\begin{tabular}{@{}lrrrrr@{}}
\toprule
指标 & 2023 年 & 2024 年 & 2025 年 & 2025 年半年报 & 2026 年半年报 \\
\midrule
营业收入（亿元） & 25.24 & 24.02 & 27.51 & 13.99 & 14.96 \\
归母净利润（亿元） & 0.54 & 1.11 & 1.51 & 0.93 & 0.82 \\
毛利率（\%） & 8.1 & 11.9 & 13.3 & 13.2 & 14.7 \\
ROE（\%） & 2.7 & 5.5 & 7.1 & 4.4 & 3.6 \\
资产负债率（\%） & 25.4 & 23.7 & 19.7 & 23.6 & 23.1 \\
经营活动现金流净额（亿元） & 2.82 & 2.49 & 2.59 & 0.23 & 1.19 \\
货币资金（亿元） & 2.43 & 2.55 & 2.74 & 2.71 & 3.13 \\
有息负债合计（亿元） & 2.93 & 2.72 & 1.16 & 3.09 & 1.77 \\
现金短债比（倍） & 0.95 & 0.94 & 2.39 & 0.88 & 2.29 \\
\bottomrule
\end{tabular}
\end{center}

\begin{center}
\begin{tikzpicture}
\begin{groupplot}[
  group style={group size=2 by 1, horizontal sep=0.95cm},
  width=7.2cm, height=4.6cm,
  tick label style={font=\tiny},
  title style={font=\scriptsize\heiti},
  yticklabel style={font=\tiny},
  grid=major, grid style={LightGray, line width=0.3pt},
  axis line style={InkGray, line width=0.5pt},
  enlarge x limits={abs=0.8},
  legend style={font=\scriptsize, at={(0.5,-0.22)}, anchor=north,
                legend columns=2, draw=none, fill=none, column sep=6pt},
]
\nextgroupplot[
  ybar, bar width=4pt,
  ymin=-2.76, ymax=30.97,
  xtick={0,1,2,3,4,5.7},
  xticklabels={2021,2022,2023,2024,2025,2026H1},
  ylabel={亿元},
]
\addplot[fill=AgriGreen!75, draw=AgriGreen, bar shift=-2.6pt] table[x=i, y=rev]{data/clean/profiles/fin_002286.dat};
\addplot[fill=SoftRed!60, draw=SoftRed, bar shift=2.6pt] table[x=i, y=ni]{data/clean/profiles/fin_002286.dat};
\addplot[fill=AgriGreen!30, draw=AgriGreen, pattern=north east lines, bar shift=-2.6pt] table[x=x, y=rev]{data/clean/profiles/finh_002286.dat};
\addplot[fill=SoftRed!20, draw=SoftRed, pattern=north east lines, bar shift=2.6pt] table[x=x, y=ni]{data/clean/profiles/finh_002286.dat};
\legend{营业收入（年/半年）, 归母净利润（年/半年）}
\nextgroupplot[
  ymin=-3.2, ymax=44.2,
  xtick={0,1,2,3,4,5.7},
  xticklabels={2021,2022,2023,2024,2025,2026H1},
  ylabel={\%},
]
\addplot[AgriGold, mark=*, mark size=1.3pt] table[x=i, y=roe]{data/clean/profiles/fin_002286.dat};
\addplot[OilBlue, mark=square*, mark size=1.3pt] table[x=i, y=debt]{data/clean/profiles/fin_002286.dat};
\addplot[only marks, AgriGold, mark=*, mark size=2.0pt] table[x=x, y=roe]{data/clean/profiles/finh_002286.dat};
\addplot[only marks, OilBlue, mark=square*, mark size=2.0pt] table[x=x, y=debt]{data/clean/profiles/finh_002286.dat};
\legend{ROE, 资产负债率}
\end{groupplot}
\end{tikzpicture}
\begin{tikzpicture}
\begin{axis}[
  width=15.2cm, height=3.1cm, scale only axis,
  ymin=-0.59, ymax=6.56,
  xtick={0,1,2,3,4,5.7},
  xticklabels={2021,2022,2023,2024,2025,2026H1},
  tick label style={font=\tiny},
  yticklabel style={font=\tiny},
  ylabel={亿元}, ylabel style={font=\scriptsize},
  ybar, bar width=4pt,
  grid=major, grid style={LightGray, line width=0.3pt},
  axis line style={InkGray, line width=0.5pt},
  enlarge x limits={abs=0.8},
  legend style={font=\scriptsize, at={(0.5,-0.30)}, anchor=north,
                legend columns=3, draw=none, fill=none, column sep=10pt},
]
\addplot[fill=AgriGreen!55, draw=AgriGreen, bar shift=-3.0pt] table[x=i, y=cash]{data/clean/profiles/fin_002286.dat};
\addplot[fill=SoftRed!45, draw=SoftRed, bar shift=3.0pt] table[x=i, y=ibd]{data/clean/profiles/fin_002286.dat};
\addplot[fill=AgriGreen!25, draw=AgriGreen, pattern=north east lines, bar shift=-3.0pt] table[x=x, y=cash]{data/clean/profiles/finh_002286.dat};
\addplot[fill=SoftRed!20, draw=SoftRed, pattern=north east lines, bar shift=3.0pt] table[x=x, y=ibd]{data/clean/profiles/finh_002286.dat};
\addplot[OilBlue, mark=*, mark size=1.1pt, line width=0.8pt] table[x=i, y=ocf]{data/clean/profiles/fin_002286.dat};
\addplot[only marks, OilBlue, mark=*, mark size=1.8pt] table[x=x, y=ocf]{data/clean/profiles/finh_002286.dat};
\legend{货币资金, 有息负债, 经营活动现金流净额}
\end{axis}
\end{tikzpicture}

\captionof{figure}[保龄宝（002286）近年主要财务指标]{保龄宝（002286）近年主要财务指标（左上：营业收入与归母净利润；右上：ROE 与资产负债率；下：货币资金、有息负债与经营活动现金流净额。
2021---2025 年为年报口径，2026H1 为 2026 年半年报/半年末，与其前一年报之间留出间隔、以斜纹或空心点区分；半年报数据未年化）}
\end{center}

\pfl{财务分析} 2026 年 6 月末资产负债率 23.1\%，较 2025 年末 上升 3.4 个百分点（样本中位数 52.2\%，所处三级行业内排第 \zhnumber{3} 位）。 2026 年 6 月末货币资金 3.13 亿元、短期债务（短期借款与一年内到期非流动负债）1.37 亿元、长期债务 0.41 亿元，有息负债合计 1.77 亿元，现金短债比 2.29 倍——覆盖充分。 净有息负债 -1.35 亿元，相当于其 2026 年上半年营业收入的 -9\%。 流动比率 1.74、速动比率 1.20，短期指标尚可。 2026 年上半年经营活动现金流净额 1.19 亿元，经营现金流与利润同向为正，内生资金可以支撑运营。 2026 年 6 月末在建工程 2.30 亿元，较上年末增加 0.99 亿元，仍有在建投入。 公司披露的财务与合规风险包括：控股股东北京永裕投资管理有限公司所持公司部分股份被法院司法冻结；本次冻结前永裕投资持有的……；公司总质押股份数量1,879.39万股，质押股份占A股总股本比例4.95\%。\footnote{\texttt{https://finance.sina.com.cn/roll/2026-06-05/doc-iniahtff2147753.shtml}}\footnote{\texttt{https://basic.10jqka.com.cn/002286/}}

\pfl{重大战略转型} 从公开披露看，公司的战略推进包括：2025-12-31 2025年度实现营业收入275,102.42万元、同比增长14.54\%；实现扣除股份支付费用影响后的归属于上市公……；2025-12-31 2025年减糖甜味剂实现收入6.9亿元、同比增长34.5\%，毛利率提高约8个百分点；益生元收入4亿元、同比增长2……；2026-05-13 披露《关于2025年限制性股票激励计划首次授予限制性股票第一个解除限售期解除限售条件成就的公告》；2026-08-26 2026年上半年拟向全体股东每10股派发现金红利0.3元（含税）；益生元等销量同比增长，海外业务高速增长，核心产……。 公告口径上，近三年（2023 年 9 月---2026 年 9 月）公司披露重大重组与并购类公告 0 条、对外投资与转型类公告 3 条、回购与股权激励类公告 34 条，合计公告 356 条。 主营结构的迁移已经落到报表上：2021 年年报的第一大行业为“食品制造业”（99.5\%），2026 年半年报变为“农副食品加工业”（100.0\%）。 综合判断，公司在报告期内有明确的外延扩张或业务结构调整动作，转型方向与融资安排之间存在直接关联。\footnote{\texttt{https://www.stcn.com/article/detail/3431212.html}}\footnote{\texttt{http://money.finance.sina.com.cn/corp/view/vCB\_AllBulletinDetail.php}}\footnote{\texttt{https://finance.eastmoney.com/a/202605023727532285.html}}

\begin{center}\small\setlength{\tabcolsep}{4pt}
\captionof{table}[保龄宝（002286）历史融资与资本运作]{保龄宝（002286）历史融资与资本运作（据公司公告与公开报道整理；金额为披露值，含首次公开发行、再融资、债券、银行借款与重整投资等）}
\begin{tabular}{@{}p{1.5cm}p{3.2cm}p{2.4cm}p{6.3cm}@{}}
\toprule
时间 & 方式 & 规模 & 用途与说明 \\
\midrule
2009-08-28 & IPO（深交所中小板） & 发行2,000万股，发行价20.56元/股，预计募集资金总额41,120万元、预计募集资金净额38,634万元（同花顺F10登记实际募集资金净额3.85亿元） & 低聚果糖和水溶性膳食纤维产业化项目；补充流动资金3,400万元…… \\
2013-04 & 非公开发行（定向增发） & 发行49,428,000股A股，发行价12.60元/股，募集资金总额622,792,800.00元，募集资金净额599,992,789.05元 & 未查得（本次发行股票全部以现金认购，不涉及资产过户）…… \\
2021-07-13 & 非公开发行A股股票（预案） & 拟发行股票数量不超过9,000万股，发行价格7.91元/股，募资不超过71,190.00万元；后续预案二次修订将拟募集资金调整为不超过70,650.00万元（反馈意见回复口径为拟募集资金7.11亿元） & 年产3万吨赤藓糖醇晶体项目25,285.55万元…… \\
2022-11-24 & 取得非公开发行核准批复 & 《关于核准保龄宝生物股份有限公司非公开发行股票的批复》（证监许可[2022]2944号） & 年产3万吨赤藓糖醇晶体项目等；公司于2022年6月14日收到证监会行政许可申请受理单（受…… \\
2023-11 & 非公开发行终止（批复到期未实施） & 原拟募资不超过70,650万元 & 2023年年度报告披露“非公开发行项目终止，本承诺至2023年11月17日结束”…… \\
\bottomrule
\end{tabular}
\end{center}

\pfl{历史融投资情况} 从现金流量表看，2025 年公司取得借款收到的现金 2.75 亿元、偿还债务支付的现金 4.32 亿元、吸收权益性投资 0.41 亿元，筹资活动现金流净额 -1.50 亿元（2023 年为取得借款 2.90 亿元、偿还债务 4.97 亿元）。 2026 年上半年筹资活动现金流净额为 0.57 亿元（取得借款 1.90 亿元、偿还债务 1.27 亿元，未年化），仍处于净融入状态。 2025 年分配股利、利润或偿付利息支付的现金合计 0.32 亿元。 公开信息中与公司直接相关的融资与资本运作包括：IPO（深交所中小板）、取得非公开发行核准批复、非公开发行A股股票（预案）、非公开发行终止（批复到期未实施）、非公开发行（定向增发）。 其中规模较大的两笔为：2009-08-28IPO（深交所中小板）（发行2,000万股，发行价20.56元/股，预计募集资金总额……）；2023-11非公开发行终止（批复到期未实施）（原拟募资不超过70,650万元）。 股本方面，近三年披露股本变动 5 次（股份回购 2 次）；总股本由 2021 年末的 3.72 亿股变为 3.70 亿股（-0.58\%）。 分红方面，近三年现金分红 6 次、累计每股派现 0.40 元（最近一次 2026-06-30）——在利润承压的阶段仍维持了分红。 近三年“再融资”类公告 9 条，涉及定向增发。\footnote{\texttt{https://pdf.dfcfw.com/pdf/H2\_AN201202290004462416\_1.pdf}}\footnote{\texttt{http://money.finance.sina.com.cn/corp/view/vISSUE\_MarketBulletinDetail.php}}\footnote{\texttt{https://file.finance.qq.com/finance/hs/pdf/2022/09/30/1214717674.PDF}}

\pfl{未来潜在融资需求分析} 2026 年 6 月末货币资金 3.13 亿元，而短期债务 1.37 亿元（现金短债比 2.29 倍），2026 年上半年经营现金流净额 1.19 亿元。账面货币资金可以覆盖短期债务，缺口不体现在静态偿付层面。 公司 2026 年 6 月末资产负债率 23.1\%、2026 年上半年 ROE 3.6\%（未年化），资产负债表尚有余量，融资需求不是“补血型”的刚性需求，而是围绕并购整合、项目投入与营运资金的主动性需求。 公开披露的下一步融资线索包括：2025-04-24 预计2025年度与山东禹城农村商业银行股份有限公司发生的关联交易业务总额不超过1.002亿元，另银行日最高存款余额不超过……；2026-04-23 第六届董事会第十六次会议审议通过《关于提请股东会授权董事会办理以简易程序向特定对象发行股票的议案》，拟以简易程序向特定对……；2026-05-29 2026年第一次临时股东会审议通过向激励对象定向发行本公司人民币A股普通股作为股票来源的议案。 资金需求还要叠加在建工程：期末在建工程 2.30 亿元、2026 年上半年购建固定资产等支出 1.32 亿元（未年化），这部分支出在周期底部通常会被主动放缓。\footnote{\texttt{https://paper.cnstock.com/html/2025-04/25/content\_2055157.htm}}\footnote{\texttt{https://disc.static.szse.cn/download/disc/disk03/finalpage/2026-04-25/9f005cea-6551-49f0}}


\subsubsection{中基健康（000972）}
\label{co:000972}

\pfl{公司基本情况} 1994 年设立至今，中基健康新疆兵团第六师国资控股的番茄制品企业，主营大包装番茄酱、小包装番茄制品与番茄红素；2024年因净资产为负、扣除后营收不足3亿元被实施退市风险警示，2025年7月进入预重整，2025年底靠政府补助与债务豁免净资产转正并于2026年6月摘帽，但2026年中期净资产再度转负。 公司于 2000-09-26 在深交所主板首发上市，注册资本 7.71 亿元，总股本 7.71 亿股。 截至 2026-09-24收盘价 3.96 元，流通市值 30.5 亿元，近一年+1.0\%，流通市值在三级行业“果蔬加工”的 5 家公司中排第 \zhnumber{4} 位，近一年涨跌幅高于全样本中位数（-10.1\%）11.1 个百分点。 按 2026 年半年报披露的行业构成，收入占比前三位为制造业 95.3\%；其他 4.5\%；租赁服务 0.3\%。 发展过程中的关键节点包括：1999年4月26日更名为“新疆中基实业股份有限公司”（1999）；1999年3月15日股东大会决议同意向社会公众发行4,500.00万股普通股；2000年……（2000）；2016年5月9日更名为“中基健康产业股份有限公司”；母公司为新疆生产建设兵团第六师国有……（2016）。\footnote{\texttt{http://money.finance.sina.com.cn/corp/view/vCB\_AllBulletinDetail.php}}\footnote{\texttt{http://static.cninfo.com.cn/finalpage/2004-02-14/13630567.PDF}}

\pfl{历史股价} 本报告样本区间（2021 年 1 月 至 2026 年 9 月）内，股价由 1.80 元变为 3.96 元，累计 +120.0\%，区间最高 4.45 元（2025 年 8 月）、最低 1.68 元（2021 年 2 月）；当前价相当于区间最高价的 89.0\%，为区间最低价的 2.36 倍。 月度收益的年化波动率 43.9\%，高于全样本中位数 37.4\%，在三级行业“果蔬加工”的 5 家公司中波动率排第 \zhnumber{4} 位。 从事件看，2022 年 4 月 至 2023 年 1 月的上涨（+139.0\%）与公司构成关联交易及重大资产重组，不构成重组上市，不导致实际控制人变更；2023年……（2022-08）在时间上重合；2024 年 6 月 至 2025 年 8 月的上涨（+99.6\%）与公司临时管理人公开招募重整投资人，分别引入产业投资人和财务投资人；招募条件明确意……（2025-08-05）在时间上重合——股价的拐点多数能在公司自身的资本运作与风险事件上找到对应。\footnote{\texttt{http://static.cninfo.com.cn/finalpage/2023-07-25/1217377184.PDF}}\footnote{\texttt{https://stcn.com/article/detail/2932670.html}}

\begin{center}
\begin{tikzpicture}
\begin{axis}[
  width=12.6cm, height=3.9cm, scale only axis,
  xmin=2020.922, xmax=2026.888, ymin=1.56, ymax=4.76,
  xtick={2021,2022,2023,2024,2025,2026},
  yearticks,
  yticklabel style={font=\scriptsize},
  ylabel={元/股}, ylabel style={font=\scriptsize},
  grid=major, grid style={LightGray, line width=0.3pt},
  axis line style={InkGray, line width=0.5pt},
  every axis plot/.append style={line width=0.9pt},
]
\addplot[AgriGreen] table[x=x,y=y]{data/clean/profiles/px_000972.dat};
\addplot[SoftRed, dashed, line width=0.7pt] table[x=x,y=y]{data/clean/profiles/pxzz_000972.dat};
\addplot[only marks, mark=*, mark size=1.1pt, SoftRed] table[x=x,y=y]{data/clean/profiles/pxzz_000972.dat};
\end{axis}
\end{tikzpicture}

\captionof{figure}[中基健康（000972）股价与周期骨架]{中基健康（000972）月度收盘价（元/股）与 ZigZag 周期骨架（阈值 25\%，圆点为周期转折点）}
\end{center}

\pfl{过去周期分析} 以 25\% 的 ZigZag 阈值划分，样本区间内股价形成 4 个主升段与 3 个主跌段。 最大主跌段为 2023 年 1 月 至 2024 年 6 月，17 个月-47.3\%；最大主升段出现在 2022 年 4 月 至 2023 年 1 月，9 个月+139.0\%。 最近一次转折出现在 2026 年 9 月（上涨段自 2026 年 6 月起，3 个月+26.5\%），当前处于该段之后的震荡之中。 公司股价较申万农林牧渔指数提前约 63 个月见底，是板块中较早定价周期反转的公司之一。 果蔬加工的外向型特征明显，出口订单与原料产地价格决定利润率，公司 2026 年 6 月末资产负债率 104.1\%，高于全样本中位数 52.2\%，其股价因此更多反映资金面与信用风险，而不是价格弹性本身。

\pfl{盈利情况} 2026 年上半年营业收入 1.97 亿元（同比 -19.6\%），归母净利润 -0.66 亿元，亏损同比收窄 0.09 亿元。 以年报为趋势对照，2023---2025 年营业收入由 5.76 亿元变为 4.92 亿元（两年年均复合增速 -7.6\%），归母净利润由 1.08 亿元变为 -0.46 亿元。 按半年报口径，2026 年上半年的 ROE 为 —\%（未年化）、销售净利率 -34.6\%、毛利率 -19.8\%。 结合公司披露，2026年上半年营业收入1.97亿元、同比下降19.6\%；归母净利润-6,592万元（上年同期-7,47……；截至二季度末公司总资产9.96亿元、较上年度末下降18.47\%；归属于上市公司股东的净资产-39,796……。\footnote{\texttt{https://finance.sina.com.cn/stock/zqgd/2026-08-28/doc-inipwnxp9011778.shtml}}\footnote{\texttt{https://stock.stockstar.com/notice/SN2026082800049854.shtml}}

\begin{center}\small\setlength{\tabcolsep}{4pt}
\captionof{table}[中基健康（000972）关键财务指标]{中基健康（000972）关键财务指标（合并报表口径；两个半年报期均未年化；现金短债比 $=$ 货币资金 $\div$ 短期债务）}
\begin{tabular}{@{}lrrrrr@{}}
\toprule
指标 & 2023 年 & 2024 年 & 2025 年 & 2025 年半年报 & 2026 年半年报 \\
\midrule
营业收入（亿元） & 5.76 & 3.08 & 4.92 & 2.45 & 1.97 \\
归母净利润（亿元） & 1.08 & -2.31 & -0.46 & -0.75 & -0.66 \\
毛利率（\%） & 25.3 & 5.0 & -15.7 & -13.8 & -19.8 \\
ROE（\%） & 72.4 & 0.0 & --- & 0.0 & --- \\
资产负债率（\%） & 82.9 & 101.0 & 97.8 & 106.8 & 104.1 \\
经营活动现金流净额（亿元） & -1.82 & -0.98 & 2.47 & -1.01 & 1.27 \\
货币资金（亿元） & 1.11 & 1.60 & 3.13 & 0.66 & 3.01 \\
有息负债合计（亿元） & 4.28 & 6.58 & 6.77 & 7.15 & 5.54 \\
现金短债比（倍） & 0.31 & 0.27 & 0.50 & 0.10 & 0.60 \\
\bottomrule
\end{tabular}
\end{center}

\begin{center}
\begin{tikzpicture}
\begin{groupplot}[
  group style={group size=2 by 1, horizontal sep=0.95cm},
  width=7.2cm, height=4.6cm,
  tick label style={font=\tiny},
  title style={font=\scriptsize\heiti},
  yticklabel style={font=\tiny},
  grid=major, grid style={LightGray, line width=0.3pt},
  axis line style={InkGray, line width=0.5pt},
  enlarge x limits={abs=0.8},
  legend style={font=\scriptsize, at={(0.5,-0.22)}, anchor=north,
                legend columns=2, draw=none, fill=none, column sep=6pt},
]
\nextgroupplot[
  ybar, bar width=4pt,
  ymin=-3.13, ymax=6.88,
  xtick={0,1,2,3,4,5.7},
  xticklabels={2021,2022,2023,2024,2025,2026H1},
  ylabel={亿元},
]
\addplot[fill=AgriGreen!75, draw=AgriGreen, bar shift=-2.6pt] table[x=i, y=rev]{data/clean/profiles/fin_000972.dat};
\addplot[fill=SoftRed!60, draw=SoftRed, bar shift=2.6pt] table[x=i, y=ni]{data/clean/profiles/fin_000972.dat};
\addplot[fill=AgriGreen!30, draw=AgriGreen, pattern=north east lines, bar shift=-2.6pt] table[x=x, y=rev]{data/clean/profiles/finh_000972.dat};
\addplot[fill=SoftRed!20, draw=SoftRed, pattern=north east lines, bar shift=2.6pt] table[x=x, y=ni]{data/clean/profiles/finh_000972.dat};
\legend{营业收入（年/半年）, 归母净利润（年/半年）}
\nextgroupplot[
  ymin=-102.4, ymax=123.3,
  xtick={0,1,2,3,4,5.7},
  xticklabels={2021,2022,2023,2024,2025,2026H1},
  ylabel={\%},
]
\addplot[AgriGold, mark=*, mark size=1.3pt] table[x=i, y=roe]{data/clean/profiles/fin_000972.dat};
\addplot[OilBlue, mark=square*, mark size=1.3pt] table[x=i, y=debt]{data/clean/profiles/fin_000972.dat};
\addplot[only marks, AgriGold, mark=*, mark size=2.0pt] table[x=x, y=roe]{data/clean/profiles/finh_000972.dat};
\addplot[only marks, OilBlue, mark=square*, mark size=2.0pt] table[x=x, y=debt]{data/clean/profiles/finh_000972.dat};
\legend{ROE, 资产负债率}
\end{groupplot}
\end{tikzpicture}
\begin{tikzpicture}
\begin{axis}[
  width=15.2cm, height=3.1cm, scale only axis,
  ymin=-2.75, ymax=7.81,
  xtick={0,1,2,3,4,5.7},
  xticklabels={2021,2022,2023,2024,2025,2026H1},
  tick label style={font=\tiny},
  yticklabel style={font=\tiny},
  ylabel={亿元}, ylabel style={font=\scriptsize},
  ybar, bar width=4pt,
  grid=major, grid style={LightGray, line width=0.3pt},
  axis line style={InkGray, line width=0.5pt},
  enlarge x limits={abs=0.8},
  legend style={font=\scriptsize, at={(0.5,-0.30)}, anchor=north,
                legend columns=3, draw=none, fill=none, column sep=10pt},
]
\addplot[fill=AgriGreen!55, draw=AgriGreen, bar shift=-3.0pt] table[x=i, y=cash]{data/clean/profiles/fin_000972.dat};
\addplot[fill=SoftRed!45, draw=SoftRed, bar shift=3.0pt] table[x=i, y=ibd]{data/clean/profiles/fin_000972.dat};
\addplot[fill=AgriGreen!25, draw=AgriGreen, pattern=north east lines, bar shift=-3.0pt] table[x=x, y=cash]{data/clean/profiles/finh_000972.dat};
\addplot[fill=SoftRed!20, draw=SoftRed, pattern=north east lines, bar shift=3.0pt] table[x=x, y=ibd]{data/clean/profiles/finh_000972.dat};
\addplot[OilBlue, mark=*, mark size=1.1pt, line width=0.8pt] table[x=i, y=ocf]{data/clean/profiles/fin_000972.dat};
\addplot[only marks, OilBlue, mark=*, mark size=1.8pt] table[x=x, y=ocf]{data/clean/profiles/finh_000972.dat};
\legend{货币资金, 有息负债, 经营活动现金流净额}
\end{axis}
\end{tikzpicture}

\captionof{figure}[中基健康（000972）近年主要财务指标]{中基健康（000972）近年主要财务指标（左上：营业收入与归母净利润；右上：ROE 与资产负债率；下：货币资金、有息负债与经营活动现金流净额。
2021---2025 年为年报口径，2026H1 为 2026 年半年报/半年末，与其前一年报之间留出间隔、以斜纹或空心点区分；半年报数据未年化）}
\end{center}

\pfl{财务分析} 2026 年 6 月末资产负债率 104.1\%，较 2025 年末 上升 6.3 个百分点（样本中位数 52.2\%，所处三级行业内排第 \zhnumber{5} 位）。 2026 年 6 月末货币资金 3.01 亿元、短期债务（短期借款与一年内到期非流动负债）5.01 亿元、长期债务 0.53 亿元，有息负债合计 5.54 亿元，现金短债比 0.60 倍——覆盖偏紧。 净有息负债 2.54 亿元，相当于其 2026 年上半年营业收入的 129\%。 流动比率 0.64、速动比率 0.50，短期偿付压力较大。 2026 年上半年经营活动现金流净额 1.27 亿元，净利润为负而经营现金流为正，差额主要来自折旧摊销与营运资本变动，说明亏损尚未完全转化为现金流出。 公司披露的财务与合规风险包括：向深圳证券交易所申请撤销退市风险警示，撤销情况以深交所审核意见为准，是否存在不确定性；审……；披露2024年年度报告及《关于公司股票交易被实施退市风险警示暨停牌的公告》（公告编号20……。\footnote{\texttt{https://disc.static.szse.cn/download/disc/disk03/finalpage/2026-04-17/33cb6f06-744e-4bc0}}\footnote{\texttt{https://paper.cnstock.com/html/2025-07/29/content\_2098842.htm}}

\pfl{重大战略转型} 公司在报告期内的战略动作可归为 4 项：2025-09-12 确定重整产业投资人为新疆新业国有资产经营（集团）有限责任公司，并确定7家财务投资人，完成重整投资人招募工作；公司……；2025-12-31 2025年度实现营业总收入49,211.56万元；归属母公司所有者的净利润-4,623.18万元；归属于上市公司……；2025-12-31 2025年度全资子公司红色番茄累计获得31,678.56万元政府补助，属与收益相关的政府补助，计入“其他收益”2……；2026-06-29 撤销退市风险警示并复牌，简称恢复为“中基健康”，日涨跌幅限制由5\%变更为10\%。 公告口径上，近三年（2023 年 9 月---2026 年 9 月）公司披露重大重组与并购类公告 16 条、对外投资与转型类公告 5 条、回购与股权激励类公告 8 条，合计公告 497 条。 第一大行业仍是“制造业”，收入占比 95.3\%，与 2021 年年报相比没有方向性变化。 综合判断，公司当前的战略重心不是扩张而是化解债务与完成重整，业务层面的调整让位于司法重整程序。\footnote{\texttt{https://paper.cnstock.com/html/2025-07/29/content\_2098842.htm}}\footnote{\texttt{https://stcn.com/article/detail/2932670.html}}\footnote{\texttt{http://file.finance.sina.com.cn/211.154.219.97:9494/MRGG/CNSESZ\_STOCK/2026/2026-4/2026-0}}

\begin{center}\small\setlength{\tabcolsep}{4pt}
\captionof{table}[中基健康（000972）历史融资与资本运作]{中基健康（000972）历史融资与资本运作（据公司公告与公开报道整理；金额为披露值，含首次公开发行、再融资、债券、银行借款与重整投资等）}
\begin{tabular}{@{}p{1.5cm}p{3.2cm}p{2.4cm}p{6.3cm}@{}}
\toprule
时间 & 方式 & 规模 & 用途与说明 \\
\midrule
2000-09-26 & IPO（深交所） & 公开发行人民币普通股4,500.00万股，发行价5.00元/股（募集资金总额按发行价计算约2.25亿元，来源未直接给出总额口径） & 2000年4月14日证监会证监发行字[2000]33号文批准…… \\
2022-08 & 发行股份购买资产并募集配套资金（预案，后终止） & 拟以发行股份方式向新疆中泰农业发展有限责任公司、新疆粮油集团有限责任公司、新疆艳阳天番茄制品有限责任公司购买其合计持有的新疆新粮艳阳天番茄股份有限公司100.00\%股份，交易对价41,316.56万元、发行股份数量184,448,914股；同时拟向不超过35名特定对象募集配套资金不超过10,000.00万元 & 收购新粮艳阳天100\%股份，配套资金用于相关项目；构成关联交易及重大资产重组…… \\
2024-02-28 & 终止发行股份购买资产并募集配套资金 & 原交易对价41,316.56万元 & 第九届董事会第四十次临时会议、第九届监事会第二十二次临时会议审议通过终止议案…… \\
2025-09-12 & 重整投资人招募与《重整投资协议》签署 & 产业投资人新疆新业国有资产经营（集团）有限责任公司；财务投资人7家（含新疆兵金建投资有限公司；广州昊智科技投资合伙企业（有限合伙）及信风投资管理有限公司联合体；新疆兵投资产管理有限责任公司及青岛鲁创私募基金管理有限公司联合体；北京博雅春芽投资有限公司及长城（天津）股权投资基金管理有限责任公司联合体等）。具体投资金额未查得 & 以重整投资人投资款解决债务、优化资产负债结构与股本结构…… \\
2025-12-29 & 债务豁免（子公司红色番茄） & 新疆生产建设兵团第六师国有资产经营有限责任公司豁免债务4,248万元（含借款本金3,942.44万元、利息305.56万元）；新疆国恒投资发展集团有限公司豁免债务5,752万元；合计增加子公司红色番茄资本公积1亿元 & 改善子公司资产负债结构；2025年12月29日与六师国资公司签订《债务豁免协议》…… \\
\bottomrule
\end{tabular}
\end{center}

\pfl{历史融投资情况} 从现金流量表看，2025 年公司取得借款收到的现金 2.19 亿元、偿还债务支付的现金 2.18 亿元、吸收权益性投资 — 亿元，筹资活动现金流净额 -0.48 亿元（2023 年为取得借款 5.18 亿元、偿还债务 1.04 亿元）。 2026 年上半年筹资活动现金流净额为 -1.29 亿元（取得借款 — 亿元、偿还债务 1.31 亿元，未年化），已转为净流出，处于净还债状态。 2025 年分配股利、利润或偿付利息支付的现金合计 0.23 亿元。 公开信息中与公司直接相关的融资与资本运作包括：IPO（深交所）、债务豁免（子公司红色番茄）、发行股份购买资产并募集配套资金（预案，后终止）、终止发行股份购买资产并募集配套资金、重整投资人招募与《重整投资协议》签署。 其中规模较大的两笔为：2024-02-28终止发行股份购买资产并募集配套资金（原交易对价41,316.56万元）；2022-08发行股份购买资产并募集配套资金（预案，后终止）（拟以发行股份方式向新疆中泰农业发展有限责任公司、新疆粮油集团……）。 股本方面，近三年披露股本变动 3 次；总股本由 2021 年末的 7.71 亿股变为 7.71 亿股（+0.00\%）。 分红方面，近三年未实施现金分红，与重整期间的资金约束一致。 近三年“再融资”类公告 36 条，涉及定向增发、发行股份购买资产。\footnote{\texttt{http://static.cninfo.com.cn/finalpage/2004-02-14/13630567.PDF}}\footnote{\texttt{http://static.cninfo.com.cn/finalpage/2023-07-25/1217377184.PDF}}\footnote{\texttt{https://file.finance.qq.com/finance/hs/pdf/2024/04/30/1219911681.PDF}}

\pfl{未来潜在融资需求分析} 2026 年 6 月末货币资金 3.01 亿元，而短期债务 5.01 亿元（现金短债比 0.60 倍），2026 年上半年经营现金流净额 1.27 亿元。按“短期债务减去账面货币资金”的滚动偿付口径，静态缺口约 2.00 亿元；若再计入上半年亏损的年化额（1.32 亿元），维持一年正常经营所需的外部资金约 2.00---3.32 亿元。 按第 \ref{sec:financing-need} 节的三档划分，公司属于第一档“补血型”：2026 年上半年归母净利润 -0.66 亿元、6 月末资产负债率 104.1\%。负债率已超过 80\%，净资产被亏损侵蚀的程度较深，属于全部样本中资产负债表最紧张的一组。 可行的融资路径按现实性排序为：引入国资或产业投资人（部分公司以控制权让渡为对价）、债务重组或展期（与银行和债券持有人协商）、定向增发（需股价配合，且控股股东需放弃优先认购或同步增资）。 公开披露的下一步融资线索包括：2026-04-28 临时管理人以部分债权审查需要补充证据材料、预重整方案仍在制定过程中、且进行重整需履行有关审批程序为由申请延长预重整期限；2026-06-26 公司披露《关于深圳证券交易所2025年年报问询函的回复》，2025年毛利率为负，经营高度依赖政府补助与债务豁免；若法院裁……；2026-07-27 六师中院通知延长公司及子公司红色番茄预重整期限三个月至2026年10月28日（此前已分别延长至2026年1月28日、20……。\footnote{\texttt{http://file.finance.sina.com.cn/211.154.219.97:9494/MRGG/CNSESZ\_STOCK/2026/2026-4/2026-0}}\footnote{\texttt{https://disc.static.szse.cn/download/disc/disk03/finalpage/2026-04-28/65357dc1-4d21-45db}}\footnote{\texttt{https://disc.static.szse.cn/download/disc/disk03/finalpage/2026-06-26/45c25b7a-1c4a-4111}}


\subsubsection{深粮控股（000019）}
\label{co:000019}

\pfl{公司基本情况} 1991 年设立至今，深粮控股深圳市属国有大型上市粮食企业（深农投旗下），承担深圳市超百万吨粮油储备任务，是深圳“米袋子”和粤港澳大湾区粮油供应保障重要载体；前身为1992年上市的原深圳市深宝实业（深深宝A），2018年发行股份购买深粮集团100\%股权完成重组转型。 公司于 1992-10-12 在深交所主板首发上市，注册资本 11.53 亿元，总股本 11.53 亿股。 截至 2026-09-24收盘价 6.78 元，流通市值 28.2 亿元，近一年+2.7\%，流通市值在三级行业“粮油加工”的 7 家公司中排第 \zhnumber{4} 位，近一年涨跌幅高于全样本中位数（-10.1\%）12.8 个百分点。 按 2026 年半年报披露的行业构成，收入占比前三位为批发与零售业 63.3\%；租赁及商业服务业 20.6\%；制造业 16.1\%。 发展过程中的关键节点包括：证券简称“深深宝A”自1992年10月12日起使用；后续简称变更为G深宝A（2006-0……（1992）；1999年9月10日深圳市投资管理公司与农产品签订《关于深圳市深宝实业股份有限公司股权转……（1999）；2018年4月3日深圳市投资控股有限公司将其持有的公司79,484,302股A股股份全部……（2018）。\footnote{\texttt{https://quote.eastmoney.com/concept/sz000019.html}}\footnote{\texttt{https://vip.stock.finance.sina.com.cn/corp/view/vCB\_AllBulletinDetail.php}}

\pfl{历史股价} 本报告样本区间（2021 年 1 月 至 2026 年 9 月）内，股价由 6.26 元变为 6.78 元，累计 +8.3\%，区间最高 9.05 元（2022 年 5 月）、最低 5.67 元（2026 年 6 月）；当前价相当于区间最高价的 74.9\%，为区间最低价的 1.20 倍。 月度收益的年化波动率 27.4\%，低于全样本中位数 37.4\%，在三级行业“粮油加工”的 7 家公司中波动率排第 \zhnumber{6} 位。

\begin{center}
\begin{tikzpicture}
\begin{axis}[
  width=12.6cm, height=3.9cm, scale only axis,
  xmin=2020.922, xmax=2026.888, ymin=5.27, ymax=9.68,
  xtick={2021,2022,2023,2024,2025,2026},
  yearticks,
  yticklabel style={font=\scriptsize},
  ylabel={元/股}, ylabel style={font=\scriptsize},
  grid=major, grid style={LightGray, line width=0.3pt},
  axis line style={InkGray, line width=0.5pt},
  every axis plot/.append style={line width=0.9pt},
]
\addplot[AgriGreen] table[x=x,y=y]{data/clean/profiles/px_000019.dat};
\addplot[SoftRed, dashed, line width=0.7pt] table[x=x,y=y]{data/clean/profiles/pxzz_000019.dat};
\addplot[only marks, mark=*, mark size=1.1pt, SoftRed] table[x=x,y=y]{data/clean/profiles/pxzz_000019.dat};
\end{axis}
\end{tikzpicture}

\captionof{figure}[深粮控股（000019）股价与周期骨架]{深粮控股（000019）月度收盘价（元/股）与 ZigZag 周期骨架（阈值 25\%，圆点为周期转折点）}
\end{center}

\pfl{过去周期分析} 以 25\% 的 ZigZag 阈值划分，样本区间内股价形成 2 个主升段与 2 个主跌段。 最大主升段出现在 2021 年 1 月 至 2022 年 5 月，16 个月+44.6\%；最大主跌段为 2022 年 5 月 至 2024 年 6 月，25 个月-36.6\%。 最近一次转折出现在 2026 年 6 月（下跌段自 2026 年 3 月起，3 个月-27.2\%），当前处于该段的延续之中。 公司股价与申万农林牧渔指数几乎同步见底，个股并未走出独立行情。 粮油加工是“成本加成”型生意（第 \ref{sec:grain} 章），利润来自原料与成品的价差及规模效率，收入规模大而净利率薄，公司 2026 年上半年实现盈利、半年 ROE 2.2\%（未年化），好于全样本中位数（0.75\%），下行周期对其报表的冲击小于同业。

\pfl{盈利情况} 2026 年上半年营业收入 23.88 亿元（同比 +0.1\%），归母净利润 1.12 亿元，同比 -36.4\%。 以年报为趋势对照，2023---2025 年营业收入由 61.90 亿元变为 55.05 亿元（两年年均复合增速 -5.7\%），归母净利润由 3.48 亿元变为 2.43 亿元。 按半年报口径，2026 年上半年的 ROE 为 2.2\%（未年化）、销售净利率 4.7\%、毛利率 16.2\%，ROE 在“粮油加工”7 家公司中排第 \zhnumber{4} 位。 行业同口径半年 ROE 的中位数为 0.75\%，该公司高于中位数 1.47 个百分点。 结合公司披露，截至2025年12月31日短期借款1,155,754,328.18元，较上年末1,484,605,101……；2025年上半年营业收入2,384,227,437.90元、同比下降1.33\%；归属于上市公司股东的净利……。\footnote{\texttt{https://disc.static.szse.cn/download/disc/disk03/finalpage/2026-04-28/a07805d7-6e55-4604}}\footnote{\texttt{https://epaper.stcn.com/pic/202508/20/996a48baeaadf8624b568c0ce3baf018.pdf}}

\begin{center}\small\setlength{\tabcolsep}{4pt}
\captionof{table}[深粮控股（000019）关键财务指标]{深粮控股（000019）关键财务指标（合并报表口径；两个半年报期均未年化；现金短债比 $=$ 货币资金 $\div$ 短期债务）}
\begin{tabular}{@{}lrrrrr@{}}
\toprule
指标 & 2023 年 & 2024 年 & 2025 年 & 2025 年半年报 & 2026 年半年报 \\
\midrule
营业收入（亿元） & 61.90 & 53.75 & 55.05 & 23.84 & 23.88 \\
归母净利润（亿元） & 3.48 & 3.25 & 2.43 & 1.76 & 1.12 \\
毛利率（\%） & 15.5 & 16.7 & 16.0 & 18.5 & 16.2 \\
ROE（\%） & 7.2 & 6.7 & 4.9 & 3.5 & 2.2 \\
资产负债率（\%） & 34.1 & 36.2 & 32.3 & 30.2 & 38.1 \\
经营活动现金流净额（亿元） & 5.86 & -0.15 & 11.00 & 7.37 & -0.53 \\
货币资金（亿元） & 2.41 & 1.68 & 0.74 & 1.80 & 0.87 \\
有息负债合计（亿元） & 12.84 & 16.06 & 12.12 & 11.86 & 21.43 \\
现金短债比（倍） & 0.19 & 0.11 & 0.06 & 0.16 & 0.05 \\
\bottomrule
\end{tabular}
\end{center}

\begin{center}
\begin{tikzpicture}
\begin{groupplot}[
  group style={group size=2 by 1, horizontal sep=0.95cm},
  width=7.2cm, height=4.6cm,
  tick label style={font=\tiny},
  title style={font=\scriptsize\heiti},
  yticklabel style={font=\tiny},
  grid=major, grid style={LightGray, line width=0.3pt},
  axis line style={InkGray, line width=0.5pt},
  enlarge x limits={abs=0.8},
  legend style={font=\scriptsize, at={(0.5,-0.22)}, anchor=north,
                legend columns=2, draw=none, fill=none, column sep=6pt},
]
\nextgroupplot[
  ybar, bar width=4pt,
  ymin=-10.14, ymax=113.56,
  xtick={0,1,2,3,4,5.7},
  xticklabels={2021,2022,2023,2024,2025,2026H1},
  ylabel={亿元},
]
\addplot[fill=AgriGreen!75, draw=AgriGreen, bar shift=-2.6pt] table[x=i, y=rev]{data/clean/profiles/fin_000019.dat};
\addplot[fill=SoftRed!60, draw=SoftRed, bar shift=2.6pt] table[x=i, y=ni]{data/clean/profiles/fin_000019.dat};
\addplot[fill=AgriGreen!30, draw=AgriGreen, pattern=north east lines, bar shift=-2.6pt] table[x=x, y=rev]{data/clean/profiles/finh_000019.dat};
\addplot[fill=SoftRed!20, draw=SoftRed, pattern=north east lines, bar shift=2.6pt] table[x=x, y=ni]{data/clean/profiles/finh_000019.dat};
\legend{营业收入（年/半年）, 归母净利润（年/半年）}
\nextgroupplot[
  ymin=-3.1, ymax=42.6,
  xtick={0,1,2,3,4,5.7},
  xticklabels={2021,2022,2023,2024,2025,2026H1},
  ylabel={\%},
]
\addplot[AgriGold, mark=*, mark size=1.3pt] table[x=i, y=roe]{data/clean/profiles/fin_000019.dat};
\addplot[OilBlue, mark=square*, mark size=1.3pt] table[x=i, y=debt]{data/clean/profiles/fin_000019.dat};
\addplot[only marks, AgriGold, mark=*, mark size=2.0pt] table[x=x, y=roe]{data/clean/profiles/finh_000019.dat};
\addplot[only marks, OilBlue, mark=square*, mark size=2.0pt] table[x=x, y=debt]{data/clean/profiles/finh_000019.dat};
\legend{ROE, 资产负债率}
\end{groupplot}
\end{tikzpicture}
\begin{tikzpicture}
\begin{axis}[
  width=15.2cm, height=3.1cm, scale only axis,
  ymin=-2.73, ymax=24.07,
  xtick={0,1,2,3,4,5.7},
  xticklabels={2021,2022,2023,2024,2025,2026H1},
  tick label style={font=\tiny},
  yticklabel style={font=\tiny},
  ylabel={亿元}, ylabel style={font=\scriptsize},
  ybar, bar width=4pt,
  grid=major, grid style={LightGray, line width=0.3pt},
  axis line style={InkGray, line width=0.5pt},
  enlarge x limits={abs=0.8},
  legend style={font=\scriptsize, at={(0.5,-0.30)}, anchor=north,
                legend columns=3, draw=none, fill=none, column sep=10pt},
]
\addplot[fill=AgriGreen!55, draw=AgriGreen, bar shift=-3.0pt] table[x=i, y=cash]{data/clean/profiles/fin_000019.dat};
\addplot[fill=SoftRed!45, draw=SoftRed, bar shift=3.0pt] table[x=i, y=ibd]{data/clean/profiles/fin_000019.dat};
\addplot[fill=AgriGreen!25, draw=AgriGreen, pattern=north east lines, bar shift=-3.0pt] table[x=x, y=cash]{data/clean/profiles/finh_000019.dat};
\addplot[fill=SoftRed!20, draw=SoftRed, pattern=north east lines, bar shift=3.0pt] table[x=x, y=ibd]{data/clean/profiles/finh_000019.dat};
\addplot[OilBlue, mark=*, mark size=1.1pt, line width=0.8pt] table[x=i, y=ocf]{data/clean/profiles/fin_000019.dat};
\addplot[only marks, OilBlue, mark=*, mark size=1.8pt] table[x=x, y=ocf]{data/clean/profiles/finh_000019.dat};
\legend{货币资金, 有息负债, 经营活动现金流净额}
\end{axis}
\end{tikzpicture}

\captionof{figure}[深粮控股（000019）近年主要财务指标]{深粮控股（000019）近年主要财务指标（左上：营业收入与归母净利润；右上：ROE 与资产负债率；下：货币资金、有息负债与经营活动现金流净额。
2021---2025 年为年报口径，2026H1 为 2026 年半年报/半年末，与其前一年报之间留出间隔、以斜纹或空心点区分；半年报数据未年化）}
\end{center}

\pfl{财务分析} 2026 年 6 月末资产负债率 38.1\%，较 2025 年末 上升 5.8 个百分点（样本中位数 52.2\%，所处三级行业内排第 \zhnumber{3} 位）。 2026 年 6 月末货币资金 0.87 亿元、短期债务（短期借款与一年内到期非流动负债）18.04 亿元、长期债务 3.39 亿元，有息负债合计 21.43 亿元，现金短债比 0.05 倍——缺口明显。 净有息负债 20.57 亿元，相当于其 2026 年上半年营业收入的 86\%。 流动比率 1.81、速动比率 0.18，短期指标尚可。 2026 年上半年经营活动现金流净额 -0.53 亿元，净利润为正但经营现金流为负，利润的现金含量偏低。 2026 年 6 月末在建工程 0.71 亿元，较上年末增加 0.19 亿元，仍有在建投入。

\pfl{重大战略转型} 从公开披露看，公司的战略推进包括：2018-2019 重组后由茶及天然植物精深加工转型为深圳市属国有粮食储备与粮油流通平台，简称变更为“深粮控股”（2019年2月28……；2025-06-30 2025年上半年营业收入2,384,227,437.90元、同比下降1.33\%；归属于上市公司股东的净利润176……；2025-12-31 2025年切实发挥区域粮食安全“压舱石”作用，扎实完成增储服务；稳步实施产业升级，经营业绩稳中有进；2025-12-31 2025年总资产73.72亿元，资产负债率32.31\%。 公告口径上，近三年（2023 年 9 月---2026 年 9 月）公司披露重大重组与并购类公告 0 条、对外投资与转型类公告 0 条、回购与股权激励类公告 0 条，合计公告 227 条。 第一大行业“批发与零售业”的收入占比由 2021 年年报的 82.5\% 变为 63.3\%（19.2 个百分点），收入结构出现再平衡。 综合判断，公司在报告期内有明确的外延扩张或业务结构调整动作，转型方向与融资安排之间存在直接关联。\footnote{\texttt{https://disc.static.szse.cn/disc/disk01/finalpage/2019-04-27/ef471555-e5d1-4f7e-9f83-3dc}}\footnote{\texttt{https://epaper.stcn.com/pic/202508/20/996a48baeaadf8624b568c0ce3baf018.pdf}}\footnote{\texttt{https://disc.static.szse.cn/download/disc/disk03/finalpage/2026-04-28/a07805d7-6e55-4604}}

\begin{center}\small\setlength{\tabcolsep}{4pt}
\captionof{table}[深粮控股（000019）历史融资与资本运作]{深粮控股（000019）历史融资与资本运作（据公司公告与公开报道整理；金额为披露值，含首次公开发行、再融资、债券、银行借款与重整投资等）}
\begin{tabular}{@{}p{1.5cm}p{3.2cm}p{2.4cm}p{6.3cm}@{}}
\toprule
时间 & 方式 & 规模 & 用途与说明 \\
\midrule
2018-03-23 & 签署重组协议 & 交易双方初步商定标的资产深粮集团100\%股权交易作价为58.59亿元，全部以发行股份的形式支付 & 本次交易完成后深深宝A将持有深粮集团100\%股权…… \\
2018-11-12 & 发行股份购买资产（重大资产重组） & 向福德资本发行655,752,951股A股购买其持有的深粮集团100\%股权；根据中企华评报字（2018）第3558号《资产评估报告》，标的资产在评估基准日2017年9月30日的评估值为587,554.64万元，经协商整体成交金额为587,554.64万元（交易双方最初初步商定作价58.59亿元、预估交易作价585,943.21万元） & 注入深圳市粮食集团有限公司（深粮集团）100\%股权，实现地方粮企借道深深宝上市…… \\
2026-01-22 & 银行综合授信额度计划（董事会决议） & 2026年度向银行申请综合授信额度不超过人民币147.5亿元；其中公司和全资子公司深圳市粮食集团有限公司拟向21家银行申请综合授信总额不超过人民币143亿元，授信均为信用方式、无需提供资产抵押 & 日常经营与粮油储备业务资金需求；同日董事会同意公司（含下属企业）2026年度通过银行流动…… \\
\bottomrule
\end{tabular}
\end{center}

\pfl{历史融投资情况} 从现金流量表看，2025 年公司取得借款收到的现金 16.76 亿元、偿还债务支付的现金 26.04 亿元、吸收权益性投资 0.00 亿元，筹资活动现金流净额 -11.55 亿元（2023 年为取得借款 19.31 亿元、偿还债务 18.72 亿元）。 2026 年上半年筹资活动现金流净额为 1.27 亿元（取得借款 12.85 亿元、偿还债务 9.70 亿元，未年化），仍处于净融入状态。 2025 年分配股利、利润或偿付利息支付的现金合计 2.01 亿元。 公开信息中与公司直接相关的融资与资本运作包括：发行股份购买资产（重大资产重组）、签署重组协议、银行综合授信额度计划（董事会决议）。 其中规模较大的两笔为：2026-01-22银行综合授信额度计划（董事会决议）（2026年度向银行申请综合授信额度不超过人民币147.5亿元……）；2018-11-12发行股份购买资产（重大资产重组）（向福德资本发行655,752,951股A股购买其持有的深粮集……）。 股本方面，近三年披露股本变动 3 次；总股本由 2021 年末的 11.53 亿股变为 11.53 亿股（+0.00\%）。 分红方面，近三年现金分红 4 次、累计每股派现 0.72 元（最近一次 2025-12-31）——在利润承压的阶段仍维持了分红。\footnote{\texttt{https://www.jiemian.com/article/2498284.html}}\footnote{\texttt{https://vip.stock.finance.sina.com.cn/corp/view/vCB\_AllBulletinDetail.php}}\footnote{\texttt{https://paper.cnstock.com/html/2026-01/22/content\_2172819.htm}}

\pfl{未来潜在融资需求分析} 2026 年 6 月末货币资金 0.87 亿元，而短期债务 18.04 亿元（现金短债比 0.05 倍），2026 年上半年经营现金流净额 -0.53 亿元。按“短期债务减去账面货币资金”的滚动偿付口径，静态缺口约 17.18 亿元；若再计入上半年亏损的年化额（0.00 亿元），维持一年正常经营所需的外部资金约 17.18---17.18 亿元。 公司 2026 年 6 月末资产负债率 38.1\%、2026 年上半年 ROE 2.2\%（未年化），资产负债表尚有余量，融资需求不是“补血型”的刚性需求，而是围绕并购整合、项目投入与营运资金的主动性需求。 公开披露的下一步融资线索包括：2026-01-22 2026年度银行综合授信额度不超过147.5亿元（其中公司及深粮集团向21家银行申请不超过143亿元），全部为信用方式；……；2026-04-28 2026年第一季度报告披露短期借款1,088,841,913.93元，较2025年末继续下降；公司主要依靠银行信用借款融……；2026-06-29 全资子公司深粮红荔粮油（深圳）有限公司拟参与竞买深圳市光明区建设用地并投资建设深圳市食用植物油储备库（光明），新增资本开……。 资金需求还要叠加在建工程：期末在建工程 0.71 亿元、2026 年上半年购建固定资产等支出 0.50 亿元（未年化），这部分支出在周期底部通常会被主动放缓。\footnote{\texttt{https://paper.cnstock.com/html/2026-01/22/content\_2172819.htm}}\footnote{\texttt{https://disc.static.szse.cn/disc/disk03/finalpage/2026-04-28/10d495a7-5ca6-47ec-9ebc-209}}\footnote{\texttt{https://quote.eastmoney.com/concept/sz000019.html}}


\subsubsection{朗源股份（300175）}
\label{co:300175}

\pfl{公司基本情况} 主营业务为“果干、坚果果仁、鲜果的加工、生产及销售”，按申万行业分类（2021 版）归属三级行业“果蔬加工”。 公司于 2011-02-15 在深交所创业板首发上市，注册资本 4.71 亿元，总股本 4.71 亿股。 截至 2026-09-24收盘价 5.35 元，流通市值 25.2 亿元，近一年-15.8\%，流通市值在三级行业“果蔬加工”的 5 家公司中排第 \zhnumber{5} 位，近一年涨跌幅低于全样本中位数（-10.1\%）5.7 个百分点。 按 2025 年年报披露的行业构成，收入占比前三位为制造业 100.0\%。

\pfl{历史股价} 以 2021 年 1 月为起点、2026 年 9 月为终点，股价由 4.60 元变为 5.35 元，累计 +16.3\%，区间最高 10.14 元（2024 年 5 月）、最低 3.21 元（2023 年 5 月）；当前价相当于区间最高价的 52.8\%，为区间最低价的 1.67 倍。 月度收益的年化波动率 52.3\%，高于全样本中位数 37.4\%，在三级行业“果蔬加工”的 5 家公司中波动率排第 \zhnumber{2} 位。

\begin{center}
\begin{tikzpicture}
\begin{axis}[
  width=12.6cm, height=3.9cm, scale only axis,
  xmin=2020.922, xmax=2026.888, ymin=2.99, ymax=10.85,
  xtick={2021,2022,2023,2024,2025,2026},
  yearticks,
  yticklabel style={font=\scriptsize},
  ylabel={元/股}, ylabel style={font=\scriptsize},
  grid=major, grid style={LightGray, line width=0.3pt},
  axis line style={InkGray, line width=0.5pt},
  every axis plot/.append style={line width=0.9pt},
]
\addplot[AgriGreen] table[x=x,y=y]{data/clean/profiles/px_300175.dat};
\addplot[SoftRed, dashed, line width=0.7pt] table[x=x,y=y]{data/clean/profiles/pxzz_300175.dat};
\addplot[only marks, mark=*, mark size=1.1pt, SoftRed] table[x=x,y=y]{data/clean/profiles/pxzz_300175.dat};
\end{axis}
\end{tikzpicture}

\captionof{figure}[朗源股份（300175）股价与周期骨架]{朗源股份（300175）月度收盘价（元/股）与 ZigZag 周期骨架（阈值 25\%，圆点为周期转折点）}
\end{center}

\pfl{过去周期分析} 以 25\% 的 ZigZag 阈值划分，样本区间内股价形成 4 个主升段与 4 个主跌段。 最大主跌段为 2024 年 5 月 至 2024 年 8 月，3 个月-47.8\%；最大主升段出现在 2023 年 5 月 至 2024 年 5 月，12 个月+215.9\%。 最近一次转折出现在 2026 年 6 月（下跌段自 2026 年 2 月起，4 个月-31.5\%），当前处于该段的延续之中。 公司股价较申万农林牧渔指数提前约 36 个月见底，是板块中较早定价周期反转的公司之一。 果蔬加工的外向型特征明显，出口订单与原料产地价格决定利润率，公司 2026 年 6 月末资产负债率 8.8\%、半年 ROE -2.7\%（未年化），在本轮下行中属于随行业同步调整的一类。

\pfl{盈利情况} 2026 年上半年营业收入 1.73 亿元（同比 +67.7\%），归母净利润 -0.13 亿元，亏损同比扩大 0.08 亿元。 以年报为趋势对照，2023---2025 年营业收入由 2.20 亿元变为 2.92 亿元（两年年均复合增速 +15.1\%），归母净利润由 -0.54 亿元变为 -0.15 亿元。 按半年报口径，2026 年上半年的 ROE 为 -2.7\%（未年化）、销售净利率 -7.8\%、毛利率 2.4\%，ROE 在“果蔬加工”5 家公司中排第 \zhnumber{4} 位。 行业同口径半年 ROE 的中位数为 0.75\%，该公司低于中位数 3.46 个百分点。

\begin{center}\small\setlength{\tabcolsep}{4pt}
\captionof{table}[朗源股份（300175）关键财务指标]{朗源股份（300175）关键财务指标（合并报表口径；两个半年报期均未年化；现金短债比 $=$ 货币资金 $\div$ 短期债务）}
\begin{tabular}{@{}lrrrrr@{}}
\toprule
指标 & 2023 年 & 2024 年 & 2025 年 & 2025 年半年报 & 2026 年半年报 \\
\midrule
营业收入（亿元） & 2.20 & 2.45 & 2.92 & 1.03 & 1.73 \\
归母净利润（亿元） & -0.54 & -0.36 & -0.15 & -0.05 & -0.13 \\
毛利率（\%） & 16.7 & 11.2 & 8.2 & 11.7 & 2.4 \\
ROE（\%） & -10.2 & -7.6 & -2.9 & -1.0 & -2.7 \\
资产负债率（\%） & 29.6 & 28.5 & 10.1 & 6.2 & 8.8 \\
经营活动现金流净额（亿元） & -0.06 & -0.35 & -0.57 & 0.14 & 0.65 \\
货币资金（亿元） & 0.22 & 0.68 & 0.29 & 1.03 & 0.84 \\
有息负债合计（亿元） & 0.62 & 0.15 & 0.33 & 0.20 & 0.28 \\
现金短债比（倍） & 0.70 & 706.92 & 0.94 & 20.35 & 61.19 \\
\bottomrule
\end{tabular}
\end{center}

\begin{center}
\begin{tikzpicture}
\begin{groupplot}[
  group style={group size=2 by 1, horizontal sep=0.95cm},
  width=7.2cm, height=4.6cm,
  tick label style={font=\tiny},
  title style={font=\scriptsize\heiti},
  yticklabel style={font=\tiny},
  grid=major, grid style={LightGray, line width=0.3pt},
  axis line style={InkGray, line width=0.5pt},
  enlarge x limits={abs=0.8},
  legend style={font=\scriptsize, at={(0.5,-0.22)}, anchor=north,
                legend columns=2, draw=none, fill=none, column sep=6pt},
]
\nextgroupplot[
  ybar, bar width=4pt,
  ymin=-0.88, ymax=3.33,
  xtick={0,1,2,3,4,5.7},
  xticklabels={2021,2022,2023,2024,2025,2026H1},
  ylabel={亿元},
]
\addplot[fill=AgriGreen!75, draw=AgriGreen, bar shift=-2.6pt] table[x=i, y=rev]{data/clean/profiles/fin_300175.dat};
\addplot[fill=SoftRed!60, draw=SoftRed, bar shift=2.6pt] table[x=i, y=ni]{data/clean/profiles/fin_300175.dat};
\addplot[fill=AgriGreen!30, draw=AgriGreen, pattern=north east lines, bar shift=-2.6pt] table[x=x, y=rev]{data/clean/profiles/finh_300175.dat};
\addplot[fill=SoftRed!20, draw=SoftRed, pattern=north east lines, bar shift=2.6pt] table[x=x, y=ni]{data/clean/profiles/finh_300175.dat};
\legend{营业收入（年/半年）, 归母净利润（年/半年）}
\nextgroupplot[
  ymin=-13.4, ymax=33.5,
  xtick={0,1,2,3,4,5.7},
  xticklabels={2021,2022,2023,2024,2025,2026H1},
  ylabel={\%},
]
\addplot[AgriGold, mark=*, mark size=1.3pt] table[x=i, y=roe]{data/clean/profiles/fin_300175.dat};
\addplot[OilBlue, mark=square*, mark size=1.3pt] table[x=i, y=debt]{data/clean/profiles/fin_300175.dat};
\addplot[only marks, AgriGold, mark=*, mark size=2.0pt] table[x=x, y=roe]{data/clean/profiles/finh_300175.dat};
\addplot[only marks, OilBlue, mark=square*, mark size=2.0pt] table[x=x, y=debt]{data/clean/profiles/finh_300175.dat};
\legend{ROE, 资产负债率}
\end{groupplot}
\end{tikzpicture}
\begin{tikzpicture}
\begin{axis}[
  width=15.2cm, height=3.1cm, scale only axis,
  ymin=-0.71, ymax=1.09,
  xtick={0,1,2,3,4,5.7},
  xticklabels={2021,2022,2023,2024,2025,2026H1},
  tick label style={font=\tiny},
  yticklabel style={font=\tiny},
  ylabel={亿元}, ylabel style={font=\scriptsize},
  ybar, bar width=4pt,
  grid=major, grid style={LightGray, line width=0.3pt},
  axis line style={InkGray, line width=0.5pt},
  enlarge x limits={abs=0.8},
  legend style={font=\scriptsize, at={(0.5,-0.30)}, anchor=north,
                legend columns=3, draw=none, fill=none, column sep=10pt},
]
\addplot[fill=AgriGreen!55, draw=AgriGreen, bar shift=-3.0pt] table[x=i, y=cash]{data/clean/profiles/fin_300175.dat};
\addplot[fill=SoftRed!45, draw=SoftRed, bar shift=3.0pt] table[x=i, y=ibd]{data/clean/profiles/fin_300175.dat};
\addplot[fill=AgriGreen!25, draw=AgriGreen, pattern=north east lines, bar shift=-3.0pt] table[x=x, y=cash]{data/clean/profiles/finh_300175.dat};
\addplot[fill=SoftRed!20, draw=SoftRed, pattern=north east lines, bar shift=3.0pt] table[x=x, y=ibd]{data/clean/profiles/finh_300175.dat};
\addplot[OilBlue, mark=*, mark size=1.1pt, line width=0.8pt] table[x=i, y=ocf]{data/clean/profiles/fin_300175.dat};
\addplot[only marks, OilBlue, mark=*, mark size=1.8pt] table[x=x, y=ocf]{data/clean/profiles/finh_300175.dat};
\legend{货币资金, 有息负债, 经营活动现金流净额}
\end{axis}
\end{tikzpicture}

\captionof{figure}[朗源股份（300175）近年主要财务指标]{朗源股份（300175）近年主要财务指标（左上：营业收入与归母净利润；右上：ROE 与资产负债率；下：货币资金、有息负债与经营活动现金流净额。
2021---2025 年为年报口径，2026H1 为 2026 年半年报/半年末，与其前一年报之间留出间隔、以斜纹或空心点区分；半年报数据未年化）}
\end{center}

\pfl{财务分析} 2026 年 6 月末资产负债率 8.8\%，较 2025 年末 下降 1.3 个百分点（样本中位数 52.2\%，所处三级行业内排第 \zhnumber{2} 位）。 2026 年 6 月末货币资金 0.84 亿元、短期债务（短期借款与一年内到期非流动负债）0.01 亿元、长期债务 0.26 亿元，有息负债合计 0.28 亿元，现金短债比 61.19 倍——覆盖充分。 净有息负债 -0.56 亿元，相当于其 2026 年上半年营业收入的 -32\%。 流动比率 14.77、速动比率 8.83，短期指标尚可。 2026 年上半年经营活动现金流净额 0.65 亿元，净利润为负而经营现金流为正，差额主要来自折旧摊销与营运资本变动，说明亏损尚未完全转化为现金流出。

\pfl{重大战略转型} 公告口径上，近三年（2023 年 9 月---2026 年 9 月）公司披露重大重组与并购类公告 2 条、对外投资与转型类公告 3 条、回购与股权激励类公告 0 条，合计公告 294 条。 第一大行业“制造业”的收入占比由 2021 年年报的 85.9\% 变为 100.0\%（14.1 个百分点），收入结构出现再平衡。 综合判断，公司报告期内未发生实质性的战略转型，资本开支与资本运作均以维持现有业务为主。

\pfl{历史融投资情况} 从现金流量表看，2025 年公司取得借款收到的现金 0.44 亿元、偿还债务支付的现金 0.29 亿元、吸收权益性投资 0.00 亿元，筹资活动现金流净额 0.13 亿元（2023 年为取得借款 0.25 亿元、偿还债务 — 亿元）。 2026 年上半年筹资活动现金流净额为 -0.05 亿元（取得借款 0.25 亿元、偿还债务 0.30 亿元，未年化），已转为净流出，处于净还债状态。 2025 年分配股利、利润或偿付利息支付的现金合计 0.01 亿元。 股本方面，近三年披露股本变动 3 次；总股本由 2021 年末的 4.71 亿股变为 4.71 亿股（+0.00\%）。 分红方面，近三年未实施现金分红，在全部样本中属于分红能力偏弱的一端。

\pfl{未来潜在融资需求分析} 2026 年 6 月末货币资金 0.84 亿元，而短期债务 0.01 亿元（现金短债比 61.19 倍），2026 年上半年经营现金流净额 0.65 亿元。账面货币资金可以覆盖短期债务，缺口不体现在静态偿付层面。 公司 2026 年上半年归母净利润为负（-0.13 亿元），但 6 月末资产负债率 8.8\% 尚未进入第一档（亏损且负债率 $>$65\%）的区间，当前融资需求以补充流动资金、支撑低谷期经营性支出为主。


\subsubsection{道道全（002852）}
\label{co:002852}

\pfl{公司基本情况} 主营业务为“食用植物油产品的研发、生产和销售”，按申万行业分类（2021 版）归属三级行业“粮油加工”。 公司于 2017-03-10 在深交所主板首发上市，注册资本 3.44 亿元，总股本 3.44 亿股。 截至 2026-09-24收盘价 8.09 元，流通市值 23.1 亿元，近一年-25.2\%，流通市值在三级行业“粮油加工”的 7 家公司中排第 \zhnumber{5} 位，近一年涨跌幅低于全样本中位数（-10.1\%）15.1 个百分点。 按 2026 年半年报披露的行业构成，收入占比前三位为食品加工 99.6\%；其他行业 0.4\%。

\pfl{历史股价} 本报告样本区间（2021 年 1 月 至 2026 年 9 月）内，股价由 12.59 元变为 8.09 元，累计 -35.7\%，区间最高 16.75 元（2021 年 12 月）、最低 6.45 元（2024 年 8 月）；当前价相当于区间最高价的 48.3\%，为区间最低价的 1.25 倍。 月度收益的年化波动率 39.6\%，高于全样本中位数 37.4\%，在三级行业“粮油加工”的 7 家公司中波动率排第 \zhnumber{3} 位。

\begin{center}
\begin{tikzpicture}
\begin{axis}[
  width=12.6cm, height=3.9cm, scale only axis,
  xmin=2020.922, xmax=2026.888, ymin=6.00, ymax=17.92,
  xtick={2021,2022,2023,2024,2025,2026},
  yearticks,
  yticklabel style={font=\scriptsize},
  ylabel={元/股}, ylabel style={font=\scriptsize},
  grid=major, grid style={LightGray, line width=0.3pt},
  axis line style={InkGray, line width=0.5pt},
  every axis plot/.append style={line width=0.9pt},
]
\addplot[AgriGreen] table[x=x,y=y]{data/clean/profiles/px_002852.dat};
\addplot[SoftRed, dashed, line width=0.7pt] table[x=x,y=y]{data/clean/profiles/pxzz_002852.dat};
\addplot[only marks, mark=*, mark size=1.1pt, SoftRed] table[x=x,y=y]{data/clean/profiles/pxzz_002852.dat};
\end{axis}
\end{tikzpicture}

\captionof{figure}[道道全（002852）股价与周期骨架]{道道全（002852）月度收盘价（元/股）与 ZigZag 周期骨架（阈值 25\%，圆点为周期转折点）}
\end{center}

\pfl{过去周期分析} 以 25\% 的 ZigZag 阈值划分，样本区间内股价形成 4 个主升段与 4 个主跌段。 最大主升段出现在 2024 年 8 月 至 2025 年 7 月，11 个月+88.1\%；最大主跌段为 2023 年 4 月 至 2024 年 8 月，16 个月-56.6\%。 最近一次转折出现在 2026 年 6 月（下跌段自 2025 年 7 月起，11 个月-41.5\%），当前处于该段的延续之中。 公司股价较申万农林牧渔指数提前约 21 个月见底，是板块中较早定价周期反转的公司之一。 粮油加工是“成本加成”型生意（第 \ref{sec:grain} 章），利润来自原料与成品的价差及规模效率，收入规模大而净利率薄，公司 2026 年 6 月末资产负债率 49.1\%、半年 ROE 0.8\%（未年化），在本轮下行中属于随行业同步调整的一类。

\pfl{盈利情况} 2026 年上半年营业收入 28.65 亿元（同比 +2.6\%），归母净利润 0.17 亿元，同比 -90.7\%。 以年报为趋势对照，2023---2025 年营业收入由 70.01 亿元变为 61.95 亿元（两年年均复合增速 -5.9\%），归母净利润由 0.76 亿元变为 2.34 亿元。 按半年报口径，2026 年上半年的 ROE 为 0.8\%（未年化）、销售净利率 0.7\%、毛利率 5.4\%，ROE 在“粮油加工”7 家公司中排第 \zhnumber{5} 位。 行业同口径半年 ROE 的中位数为 0.75\%，该公司低于中位数 0.00 个百分点。

\begin{center}\small\setlength{\tabcolsep}{4pt}
\captionof{table}[道道全（002852）关键财务指标]{道道全（002852）关键财务指标（合并报表口径；两个半年报期均未年化；现金短债比 $=$ 货币资金 $\div$ 短期债务）}
\begin{tabular}{@{}lrrrrr@{}}
\toprule
指标 & 2023 年 & 2024 年 & 2025 年 & 2025 年半年报 & 2026 年半年报 \\
\midrule
营业收入（亿元） & 70.01 & 59.43 & 61.95 & 27.92 & 28.65 \\
归母净利润（亿元） & 0.76 & 1.77 & 2.34 & 1.81 & 0.17 \\
毛利率（\%） & 7.6 & 9.6 & 8.3 & 11.6 & 5.4 \\
ROE（\%） & 3.8 & 8.6 & 10.7 & 8.2 & 0.8 \\
资产负债率（\%） & 57.2 & 46.1 & 44.8 & 48.1 & 49.1 \\
经营活动现金流净额（亿元） & 30.04 & 31.63 & 9.93 & 3.25 & -2.12 \\
货币资金（亿元） & 2.57 & 2.89 & 5.56 & 1.69 & 4.92 \\
有息负债合计（亿元） & 16.11 & 10.08 & 9.96 & 11.43 & 12.69 \\
现金短债比（倍） & 0.16 & 0.29 & 0.56 & 0.15 & 0.39 \\
\bottomrule
\end{tabular}
\end{center}

\begin{center}
\begin{tikzpicture}
\begin{groupplot}[
  group style={group size=2 by 1, horizontal sep=0.95cm},
  width=7.2cm, height=4.6cm,
  tick label style={font=\tiny},
  title style={font=\scriptsize\heiti},
  yticklabel style={font=\tiny},
  grid=major, grid style={LightGray, line width=0.3pt},
  axis line style={InkGray, line width=0.5pt},
  enlarge x limits={abs=0.8},
  legend style={font=\scriptsize, at={(0.5,-0.22)}, anchor=north,
                legend columns=2, draw=none, fill=none, column sep=6pt},
]
\nextgroupplot[
  ybar, bar width=4pt,
  ymin=-11.63, ymax=79.22,
  xtick={0,1,2,3,4,5.7},
  xticklabels={2021,2022,2023,2024,2025,2026H1},
  ylabel={亿元},
]
\addplot[fill=AgriGreen!75, draw=AgriGreen, bar shift=-2.6pt] table[x=i, y=rev]{data/clean/profiles/fin_002852.dat};
\addplot[fill=SoftRed!60, draw=SoftRed, bar shift=2.6pt] table[x=i, y=ni]{data/clean/profiles/fin_002852.dat};
\addplot[fill=AgriGreen!30, draw=AgriGreen, pattern=north east lines, bar shift=-2.6pt] table[x=x, y=rev]{data/clean/profiles/finh_002852.dat};
\addplot[fill=SoftRed!20, draw=SoftRed, pattern=north east lines, bar shift=2.6pt] table[x=x, y=ni]{data/clean/profiles/finh_002852.dat};
\legend{营业收入（年/半年）, 归母净利润（年/半年）}
\nextgroupplot[
  ymin=-25.3, ymax=65.9,
  xtick={0,1,2,3,4,5.7},
  xticklabels={2021,2022,2023,2024,2025,2026H1},
  ylabel={\%},
]
\addplot[AgriGold, mark=*, mark size=1.3pt] table[x=i, y=roe]{data/clean/profiles/fin_002852.dat};
\addplot[OilBlue, mark=square*, mark size=1.3pt] table[x=i, y=debt]{data/clean/profiles/fin_002852.dat};
\addplot[only marks, AgriGold, mark=*, mark size=2.0pt] table[x=x, y=roe]{data/clean/profiles/finh_002852.dat};
\addplot[only marks, OilBlue, mark=square*, mark size=2.0pt] table[x=x, y=debt]{data/clean/profiles/finh_002852.dat};
\legend{ROE, 资产负债率}
\end{groupplot}
\end{tikzpicture}
\begin{tikzpicture}
\begin{axis}[
  width=15.2cm, height=3.1cm, scale only axis,
  ymin=-5.50, ymax=35.68,
  xtick={0,1,2,3,4,5.7},
  xticklabels={2021,2022,2023,2024,2025,2026H1},
  tick label style={font=\tiny},
  yticklabel style={font=\tiny},
  ylabel={亿元}, ylabel style={font=\scriptsize},
  ybar, bar width=4pt,
  grid=major, grid style={LightGray, line width=0.3pt},
  axis line style={InkGray, line width=0.5pt},
  enlarge x limits={abs=0.8},
  legend style={font=\scriptsize, at={(0.5,-0.30)}, anchor=north,
                legend columns=3, draw=none, fill=none, column sep=10pt},
]
\addplot[fill=AgriGreen!55, draw=AgriGreen, bar shift=-3.0pt] table[x=i, y=cash]{data/clean/profiles/fin_002852.dat};
\addplot[fill=SoftRed!45, draw=SoftRed, bar shift=3.0pt] table[x=i, y=ibd]{data/clean/profiles/fin_002852.dat};
\addplot[fill=AgriGreen!25, draw=AgriGreen, pattern=north east lines, bar shift=-3.0pt] table[x=x, y=cash]{data/clean/profiles/finh_002852.dat};
\addplot[fill=SoftRed!20, draw=SoftRed, pattern=north east lines, bar shift=3.0pt] table[x=x, y=ibd]{data/clean/profiles/finh_002852.dat};
\addplot[OilBlue, mark=*, mark size=1.1pt, line width=0.8pt] table[x=i, y=ocf]{data/clean/profiles/fin_002852.dat};
\addplot[only marks, OilBlue, mark=*, mark size=1.8pt] table[x=x, y=ocf]{data/clean/profiles/finh_002852.dat};
\legend{货币资金, 有息负债, 经营活动现金流净额}
\end{axis}
\end{tikzpicture}

\captionof{figure}[道道全（002852）近年主要财务指标]{道道全（002852）近年主要财务指标（左上：营业收入与归母净利润；右上：ROE 与资产负债率；下：货币资金、有息负债与经营活动现金流净额。
2021---2025 年为年报口径，2026H1 为 2026 年半年报/半年末，与其前一年报之间留出间隔、以斜纹或空心点区分；半年报数据未年化）}
\end{center}

\pfl{财务分析} 2026 年 6 月末资产负债率 49.1\%，较 2025 年末 上升 4.3 个百分点（样本中位数 52.2\%，所处三级行业内排第 \zhnumber{5} 位）。 2026 年 6 月末货币资金 4.92 亿元、短期债务（短期借款与一年内到期非流动负债）12.69 亿元、长期债务 0.00 亿元，有息负债合计 12.69 亿元，现金短债比 0.39 倍——覆盖偏紧。 净有息负债 7.77 亿元，相当于其 2026 年上半年营业收入的 27\%。 流动比率 1.07、速动比率 0.40，短期指标尚可。 2026 年上半年经营活动现金流净额 -2.12 亿元，净利润为正但经营现金流为负，利润的现金含量偏低。 2026 年 6 月末在建工程 2.35 亿元，较上年末增加 0.63 亿元，仍有在建投入。

\pfl{重大战略转型} 公告口径上，近三年（2023 年 9 月---2026 年 9 月）公司披露重大重组与并购类公告 0 条、对外投资与转型类公告 0 条、回购与股权激励类公告 0 条，合计公告 300 条。 第一大行业仍是“食品加工”，收入占比 99.6\%，与 2021 年年报相比没有方向性变化。 综合判断，公司报告期内未发生实质性的战略转型，资本开支与资本运作均以维持现有业务为主。

\pfl{历史融投资情况} 从现金流量表看，2025 年公司取得借款收到的现金 29.40 亿元、偿还债务支付的现金 34.11 亿元、吸收权益性投资 — 亿元，筹资活动现金流净额 -6.18 亿元（2023 年为取得借款 29.54 亿元、偿还债务 57.57 亿元）。 2026 年上半年筹资活动现金流净额为 2.19 亿元（取得借款 11.21 亿元、偿还债务 8.48 亿元，未年化），仍处于净融入状态。 2025 年分配股利、利润或偿付利息支付的现金合计 1.46 亿元。 股本方面，近三年披露股本变动 4 次（股份回购 1 次）；总股本由 2021 年末的 3.59 亿股变为 3.44 亿股（-4.19\%）。 分红方面，近三年现金分红 7 次、累计每股派现 0.85 元（最近一次 2026-06-30）——在利润承压的阶段仍维持了分红。 近三年“再融资”类公告 3 条，涉及定向增发、可转债。

\pfl{未来潜在融资需求分析} 2026 年 6 月末货币资金 4.92 亿元，而短期债务 12.69 亿元（现金短债比 0.39 倍），2026 年上半年经营现金流净额 -2.12 亿元。按“短期债务减去账面货币资金”的滚动偿付口径，静态缺口约 7.77 亿元；若再计入上半年亏损的年化额（0.00 亿元），维持一年正常经营所需的外部资金约 7.77---7.77 亿元。 公司 2026 年 6 月末资产负债率 49.1\%、2026 年上半年 ROE 0.8\%（未年化），资产负债表尚有余量，融资需求不是“补血型”的刚性需求，而是围绕并购整合、项目投入与营运资金的主动性需求。 资金需求还要叠加在建工程：期末在建工程 2.35 亿元、2026 年上半年购建固定资产等支出 0.92 亿元（未年化），这部分支出在周期底部通常会被主动放缓。


\subsubsection{嘉华股份（603182）}
\label{co:603182}

\pfl{公司基本情况} 主营业务为“大豆蛋白及相关副产品的研发、生产及销售”，按申万行业分类（2021 版）归属三级行业“粮油加工”。 公司于 2022-09-09 在上交所首发上市，注册资本 1.65 亿元，总股本 1.65 亿股。 截至 2026-09-24收盘价 13.52 元，流通市值 22.2 亿元，近一年-2.9\%，流通市值在三级行业“粮油加工”的 7 家公司中排第 \zhnumber{6} 位，近一年涨跌幅高于全样本中位数（-10.1\%）7.1 个百分点。 按 2025 年年报披露的行业构成，收入占比前三位为农副食品加工业 92.9\%；蒸汽、电力 6.8\%；其他(补充) 0.3\%。

\pfl{历史股价} 自 2022 年 9 月至 2026 年 9 月，股价由 17.23 元变为 13.52 元，累计 -21.5\%，区间最高 19.27 元（2022 年 11 月）、最低 10.28 元（2024 年 8 月）；当前价相当于区间最高价的 70.2\%，为区间最低价的 1.32 倍。 月度收益的年化波动率 25.3\%，低于全样本中位数 37.4\%，在三级行业“粮油加工”的 7 家公司中波动率排第 \zhnumber{7} 位。

\begin{center}
\begin{tikzpicture}
\begin{axis}[
  width=12.6cm, height=3.9cm, scale only axis,
  xmin=2022.588, xmax=2026.888, ymin=9.56, ymax=20.62,
  xtick={2022,2023,2024,2025,2026},
  yearticks,
  yticklabel style={font=\scriptsize},
  ylabel={元/股}, ylabel style={font=\scriptsize},
  grid=major, grid style={LightGray, line width=0.3pt},
  axis line style={InkGray, line width=0.5pt},
  every axis plot/.append style={line width=0.9pt},
]
\addplot[AgriGreen] table[x=x,y=y]{data/clean/profiles/px_603182.dat};
\addplot[SoftRed, dashed, line width=0.7pt] table[x=x,y=y]{data/clean/profiles/pxzz_603182.dat};
\addplot[only marks, mark=*, mark size=1.1pt, SoftRed] table[x=x,y=y]{data/clean/profiles/pxzz_603182.dat};
\end{axis}
\end{tikzpicture}

\captionof{figure}[嘉华股份（603182）股价与周期骨架]{嘉华股份（603182）月度收盘价（元/股）与 ZigZag 周期骨架（阈值 25\%，圆点为周期转折点）}
\end{center}

\pfl{过去周期分析} 以 25\% 的 ZigZag 阈值划分，样本区间内股价形成 2 个主升段与 2 个主跌段。 最大主跌段为 2022 年 11 月 至 2024 年 8 月，21 个月-46.7\%；最大主升段出现在 2024 年 8 月 至 2026 年 2 月，18 个月+58.6\%。 最近一次转折出现在 2026 年 6 月（下跌段自 2026 年 2 月起，4 个月-30.1\%），当前处于该段的延续之中。 公司股价较申万农林牧渔指数提前约 21 个月见底，是板块中较早定价周期反转的公司之一。 粮油加工是“成本加成”型生意（第 \ref{sec:grain} 章），利润来自原料与成品的价差及规模效率，收入规模大而净利率薄，公司 2026 年上半年实现盈利、半年 ROE 5.0\%（未年化），好于全样本中位数（0.75\%），下行周期对其报表的冲击小于同业。

\pfl{盈利情况} 2026 年上半年营业收入 8.56 亿元（同比 +25.7\%），归母净利润 0.57 亿元，同比 -9.0\%。 以年报为趋势对照，2023---2025 年营业收入由 16.95 亿元变为 13.86 亿元（两年年均复合增速 -9.6\%），归母净利润由 1.10 亿元变为 0.95 亿元。 按半年报口径，2026 年上半年的 ROE 为 5.0\%（未年化）、销售净利率 6.6\%、毛利率 12.7\%，ROE 在“粮油加工”7 家公司中排第 \zhnumber{1} 位。 行业同口径半年 ROE 的中位数为 0.75\%，该公司高于中位数 4.25 个百分点。

\begin{center}\small\setlength{\tabcolsep}{4pt}
\captionof{table}[嘉华股份（603182）关键财务指标]{嘉华股份（603182）关键财务指标（合并报表口径；两个半年报期均未年化；现金短债比 $=$ 货币资金 $\div$ 短期债务）}
\begin{tabular}{@{}lrrrrr@{}}
\toprule
指标 & 2023 年 & 2024 年 & 2025 年 & 2025 年半年报 & 2026 年半年报 \\
\midrule
营业收入（亿元） & 16.95 & 14.88 & 13.86 & 6.81 & 8.56 \\
归母净利润（亿元） & 1.10 & 1.09 & 0.95 & 0.62 & 0.57 \\
毛利率（\%） & 12.7 & 12.5 & 13.1 & 16.0 & 12.7 \\
ROE（\%） & 11.1 & 10.4 & 8.6 & 5.6 & 5.0 \\
资产负债率（\%） & 19.9 & 24.9 & 19.6 & 17.9 & 18.8 \\
经营活动现金流净额（亿元） & 0.37 & 3.14 & 0.79 & 0.04 & 0.70 \\
货币资金（亿元） & 3.06 & 1.47 & 0.94 & 0.75 & 1.39 \\
有息负债合计（亿元） & 0.35 & 0.25 & 0.09 & --- & 0.14 \\
现金短债比（倍） & 8.75 & 5.89 & 31.33 & --- & --- \\
\bottomrule
\end{tabular}
\end{center}

\begin{center}
\begin{tikzpicture}
\begin{groupplot}[
  group style={group size=2 by 1, horizontal sep=0.95cm},
  width=7.2cm, height=4.6cm,
  tick label style={font=\tiny},
  title style={font=\scriptsize\heiti},
  yticklabel style={font=\tiny},
  grid=major, grid style={LightGray, line width=0.3pt},
  axis line style={InkGray, line width=0.5pt},
  enlarge x limits={abs=0.8},
  legend style={font=\scriptsize, at={(0.5,-0.22)}, anchor=north,
                legend columns=2, draw=none, fill=none, column sep=6pt},
]
\nextgroupplot[
  ybar, bar width=4pt,
  ymin=-1.69, ymax=18.98,
  xtick={0,1,2,3,4,5.7},
  xticklabels={2021,2022,2023,2024,2025,2026H1},
  ylabel={亿元},
]
\addplot[fill=AgriGreen!75, draw=AgriGreen, bar shift=-2.6pt] table[x=i, y=rev]{data/clean/profiles/fin_603182.dat};
\addplot[fill=SoftRed!60, draw=SoftRed, bar shift=2.6pt] table[x=i, y=ni]{data/clean/profiles/fin_603182.dat};
\addplot[fill=AgriGreen!30, draw=AgriGreen, pattern=north east lines, bar shift=-2.6pt] table[x=x, y=rev]{data/clean/profiles/finh_603182.dat};
\addplot[fill=SoftRed!20, draw=SoftRed, pattern=north east lines, bar shift=2.6pt] table[x=x, y=ni]{data/clean/profiles/finh_603182.dat};
\legend{营业收入（年/半年）, 归母净利润（年/半年）}
\nextgroupplot[
  ymin=-3.4, ymax=46.1,
  xtick={0,1,2,3,4,5.7},
  xticklabels={2021,2022,2023,2024,2025,2026H1},
  ylabel={\%},
]
\addplot[AgriGold, mark=*, mark size=1.3pt] table[x=i, y=roe]{data/clean/profiles/fin_603182.dat};
\addplot[OilBlue, mark=square*, mark size=1.3pt] table[x=i, y=debt]{data/clean/profiles/fin_603182.dat};
\addplot[only marks, AgriGold, mark=*, mark size=2.0pt] table[x=x, y=roe]{data/clean/profiles/finh_603182.dat};
\addplot[only marks, OilBlue, mark=square*, mark size=2.0pt] table[x=x, y=debt]{data/clean/profiles/finh_603182.dat};
\legend{ROE, 资产负债率}
\end{groupplot}
\end{tikzpicture}
\begin{tikzpicture}
\begin{axis}[
  width=15.2cm, height=3.1cm, scale only axis,
  ymin=-0.44, ymax=4.90,
  xtick={0,1,2,3,4,5.7},
  xticklabels={2021,2022,2023,2024,2025,2026H1},
  tick label style={font=\tiny},
  yticklabel style={font=\tiny},
  ylabel={亿元}, ylabel style={font=\scriptsize},
  ybar, bar width=4pt,
  grid=major, grid style={LightGray, line width=0.3pt},
  axis line style={InkGray, line width=0.5pt},
  enlarge x limits={abs=0.8},
  legend style={font=\scriptsize, at={(0.5,-0.30)}, anchor=north,
                legend columns=3, draw=none, fill=none, column sep=10pt},
]
\addplot[fill=AgriGreen!55, draw=AgriGreen, bar shift=-3.0pt] table[x=i, y=cash]{data/clean/profiles/fin_603182.dat};
\addplot[fill=SoftRed!45, draw=SoftRed, bar shift=3.0pt] table[x=i, y=ibd]{data/clean/profiles/fin_603182.dat};
\addplot[fill=AgriGreen!25, draw=AgriGreen, pattern=north east lines, bar shift=-3.0pt] table[x=x, y=cash]{data/clean/profiles/finh_603182.dat};
\addplot[fill=SoftRed!20, draw=SoftRed, pattern=north east lines, bar shift=3.0pt] table[x=x, y=ibd]{data/clean/profiles/finh_603182.dat};
\addplot[OilBlue, mark=*, mark size=1.1pt, line width=0.8pt] table[x=i, y=ocf]{data/clean/profiles/fin_603182.dat};
\addplot[only marks, OilBlue, mark=*, mark size=1.8pt] table[x=x, y=ocf]{data/clean/profiles/finh_603182.dat};
\legend{货币资金, 有息负债, 经营活动现金流净额}
\end{axis}
\end{tikzpicture}

\captionof{figure}[嘉华股份（603182）近年主要财务指标]{嘉华股份（603182）近年主要财务指标（左上：营业收入与归母净利润；右上：ROE 与资产负债率；下：货币资金、有息负债与经营活动现金流净额。
2021---2025 年为年报口径，2026H1 为 2026 年半年报/半年末，与其前一年报之间留出间隔、以斜纹或空心点区分；半年报数据未年化）}
\end{center}

\pfl{财务分析} 2026 年 6 月末资产负债率 18.8\%，较 2025 年末 下降 0.8 个百分点（样本中位数 52.2\%，所处三级行业内排第 \zhnumber{2} 位）。 流动比率 3.27、速动比率 1.98，短期指标尚可。 2026 年上半年经营活动现金流净额 0.70 亿元，经营现金流与利润同向为正，内生资金可以支撑运营。

\pfl{重大战略转型} 公告口径上，近三年（2023 年 9 月---2026 年 9 月）公司披露重大重组与并购类公告 6 条、对外投资与转型类公告 0 条、回购与股权激励类公告 0 条，合计公告 293 条。 第一大行业仍是“农副食品加工业”，收入占比 92.9\%，与 2021 年年报相比没有方向性变化。 综合判断，公司在报告期内有明确的外延扩张或业务结构调整动作，转型方向与融资安排之间存在直接关联。

\pfl{历史融投资情况} 从现金流量表看，2025 年公司取得借款收到的现金 0.09 亿元、偿还债务支付的现金 0.25 亿元、吸收权益性投资 — 亿元，筹资活动现金流净额 -0.82 亿元（2023 年为取得借款 1.25 亿元、偿还债务 0.90 亿元）。 2026 年上半年筹资活动现金流净额为 0.05 亿元（取得借款 0.08 亿元、偿还债务 0.03 亿元，未年化），仍处于净融入状态。 2025 年分配股利、利润或偿付利息支付的现金合计 0.66 亿元。 股本方面，近三年披露股本变动 4 次（限售股份上市 1 次）；总股本由 2021 年末的 1.65 亿股变为 1.65 亿股（+0.00\%）。 分红方面，近三年现金分红 6 次、累计每股派现 1.50 元（最近一次 2026-06-30）——在利润承压的阶段仍维持了分红。 近三年“再融资”类公告 9 条，涉及定向增发、募集资金使用。

\pfl{未来潜在融资需求分析} 公司 2026 年 6 月末资产负债率 18.8\%、2026 年上半年 ROE 5.0\%（未年化），资产负债表稳健、盈利水平平淡，暂不存在刚性融资需求；潜在的融资需求以技改扩产、并购或被并购为主。


\subsubsection{祖名股份（003030）}
\label{co:003030}

\pfl{公司基本情况} 主营业务为“豆制品的研发、生产和销售”，按申万行业分类（2021 版）归属三级行业“其他农产品加工”。 公司于 2021-01-06 在深交所主板首发上市，注册资本 1.25 亿元，总股本 1.25 亿股。 截至 2026-09-24收盘价 26.87 元，流通市值 21.9 亿元，近一年+46.9\%，流通市值在三级行业“其他农产品加工”的 11 家公司中排第 \zhnumber{9} 位，近一年涨跌幅高于全样本中位数（-10.1\%）57.0 个百分点。 按 2026 年半年报披露的产品构成，收入占比前三位为生鲜豆制品 67.3\%；植物蛋白饮品 16.3\%；其他 12.3\%。

\pfl{历史股价} 本报告样本区间（2021 年 1 月 至 2026 年 9 月）内，股价由 42.60 元变为 26.87 元，累计 -36.9\%，区间最高 42.60 元（2021 年 1 月）、最低 12.99 元（2024 年 8 月）；当前价相当于区间最高价的 63.1\%，为区间最低价的 2.07 倍。 月度收益的年化波动率 34.2\%，低于全样本中位数 37.4\%，在三级行业“其他农产品加工”的 11 家公司中波动率排第 \zhnumber{7} 位。

\begin{center}
\begin{tikzpicture}
\begin{axis}[
  width=12.6cm, height=3.9cm, scale only axis,
  xmin=2020.922, xmax=2026.888, ymin=12.08, ymax=45.58,
  xtick={2021,2022,2023,2024,2025,2026},
  yearticks,
  yticklabel style={font=\scriptsize},
  ylabel={元/股}, ylabel style={font=\scriptsize},
  grid=major, grid style={LightGray, line width=0.3pt},
  axis line style={InkGray, line width=0.5pt},
  every axis plot/.append style={line width=0.9pt},
]
\addplot[AgriGreen] table[x=x,y=y]{data/clean/profiles/px_003030.dat};
\addplot[SoftRed, dashed, line width=0.7pt] table[x=x,y=y]{data/clean/profiles/pxzz_003030.dat};
\addplot[only marks, mark=*, mark size=1.1pt, SoftRed] table[x=x,y=y]{data/clean/profiles/pxzz_003030.dat};
\end{axis}
\end{tikzpicture}

\captionof{figure}[祖名股份（003030）股价与周期骨架]{祖名股份（003030）月度收盘价（元/股）与 ZigZag 周期骨架（阈值 25\%，圆点为周期转折点）}
\end{center}

\pfl{过去周期分析} 以 25\% 的 ZigZag 阈值划分，样本区间内股价形成 2 个主升段与 2 个主跌段。 最大主跌段为 2022 年 12 月 至 2024 年 8 月，20 个月-55.4\%；最大主升段出现在 2024 年 8 月 至 2026 年 9 月，25 个月+106.9\%。 最近一次转折出现在 2026 年 9 月（上涨段自 2024 年 8 月起，25 个月+106.9\%），当前处于该段之后的震荡之中。 公司股价较申万农林牧渔指数提前约 21 个月见底，是板块中较早定价周期反转的公司之一。 该子行业横跨糖、番茄制品、添加剂等多个品种，各自的供给周期与政策（进口配额、关税）是主要变量，公司 2026 年 6 月末资产负债率 57.9\%、半年 ROE -2.5\%（未年化），在本轮下行中属于随行业同步调整的一类。

\pfl{盈利情况} 2026 年上半年营业收入 10.13 亿元（同比 +10.0\%），归母净利润 -0.24 亿元，亏损同比扩大 0.15 亿元。 以年报为趋势对照，2023---2025 年营业收入由 14.78 亿元变为 20.25 亿元（两年年均复合增速 +17.0\%），归母净利润由 0.40 亿元变为 0.25 亿元。 按半年报口径，2026 年上半年的 ROE 为 -2.5\%（未年化）、销售净利率 -2.8\%、毛利率 21.6\%，ROE 在“其他农产品加工”11 家公司中排第 \zhnumber{10} 位。 行业同口径半年 ROE 的中位数为 0.75\%，该公司低于中位数 3.20 个百分点。

\begin{center}\small\setlength{\tabcolsep}{4pt}
\captionof{table}[祖名股份（003030）关键财务指标]{祖名股份（003030）关键财务指标（合并报表口径；两个半年报期均未年化；现金短债比 $=$ 货币资金 $\div$ 短期债务）}
\begin{tabular}{@{}lrrrrr@{}}
\toprule
指标 & 2023 年 & 2024 年 & 2025 年 & 2025 年半年报 & 2026 年半年报 \\
\midrule
营业收入（亿元） & 14.78 & 16.64 & 20.25 & 9.21 & 10.13 \\
归母净利润（亿元） & 0.40 & -0.25 & 0.25 & -0.09 & -0.24 \\
毛利率（\%） & 27.3 & 22.7 & 23.8 & 24.1 & 21.6 \\
ROE（\%） & 3.9 & -2.5 & 2.5 & -0.9 & -2.5 \\
资产负债率（\%） & 47.3 & 55.6 & 54.6 & 58.2 & 57.9 \\
经营活动现金流净额（亿元） & 1.76 & 1.07 & 1.25 & 0.55 & 0.58 \\
货币资金（亿元） & 1.95 & 1.42 & 1.18 & 1.96 & 1.80 \\
有息负债合计（亿元） & 5.90 & 8.81 & 8.49 & 9.51 & 10.02 \\
现金短债比（倍） & 0.58 & 0.22 & 0.19 & 0.33 & 0.21 \\
\bottomrule
\end{tabular}
\end{center}

\begin{center}
\begin{tikzpicture}
\begin{groupplot}[
  group style={group size=2 by 1, horizontal sep=0.95cm},
  width=7.2cm, height=4.6cm,
  tick label style={font=\tiny},
  title style={font=\scriptsize\heiti},
  yticklabel style={font=\tiny},
  grid=major, grid style={LightGray, line width=0.3pt},
  axis line style={InkGray, line width=0.5pt},
  enlarge x limits={abs=0.8},
  legend style={font=\scriptsize, at={(0.5,-0.22)}, anchor=north,
                legend columns=2, draw=none, fill=none, column sep=6pt},
]
\nextgroupplot[
  ybar, bar width=4pt,
  ymin=-2.30, ymax=22.71,
  xtick={0,1,2,3,4,5.7},
  xticklabels={2021,2022,2023,2024,2025,2026H1},
  ylabel={亿元},
]
\addplot[fill=AgriGreen!75, draw=AgriGreen, bar shift=-2.6pt] table[x=i, y=rev]{data/clean/profiles/fin_003030.dat};
\addplot[fill=SoftRed!60, draw=SoftRed, bar shift=2.6pt] table[x=i, y=ni]{data/clean/profiles/fin_003030.dat};
\addplot[fill=AgriGreen!30, draw=AgriGreen, pattern=north east lines, bar shift=-2.6pt] table[x=x, y=rev]{data/clean/profiles/finh_003030.dat};
\addplot[fill=SoftRed!20, draw=SoftRed, pattern=north east lines, bar shift=2.6pt] table[x=x, y=ni]{data/clean/profiles/finh_003030.dat};
\legend{营业收入（年/半年）, 归母净利润（年/半年）}
\nextgroupplot[
  ymin=-7.3, ymax=64.0,
  xtick={0,1,2,3,4,5.7},
  xticklabels={2021,2022,2023,2024,2025,2026H1},
  ylabel={\%},
]
\addplot[AgriGold, mark=*, mark size=1.3pt] table[x=i, y=roe]{data/clean/profiles/fin_003030.dat};
\addplot[OilBlue, mark=square*, mark size=1.3pt] table[x=i, y=debt]{data/clean/profiles/fin_003030.dat};
\addplot[only marks, AgriGold, mark=*, mark size=2.0pt] table[x=x, y=roe]{data/clean/profiles/finh_003030.dat};
\addplot[only marks, OilBlue, mark=square*, mark size=2.0pt] table[x=x, y=debt]{data/clean/profiles/finh_003030.dat};
\legend{ROE, 资产负债率}
\end{groupplot}
\end{tikzpicture}
\begin{tikzpicture}
\begin{axis}[
  width=15.2cm, height=3.1cm, scale only axis,
  ymin=-1.00, ymax=11.23,
  xtick={0,1,2,3,4,5.7},
  xticklabels={2021,2022,2023,2024,2025,2026H1},
  tick label style={font=\tiny},
  yticklabel style={font=\tiny},
  ylabel={亿元}, ylabel style={font=\scriptsize},
  ybar, bar width=4pt,
  grid=major, grid style={LightGray, line width=0.3pt},
  axis line style={InkGray, line width=0.5pt},
  enlarge x limits={abs=0.8},
  legend style={font=\scriptsize, at={(0.5,-0.30)}, anchor=north,
                legend columns=3, draw=none, fill=none, column sep=10pt},
]
\addplot[fill=AgriGreen!55, draw=AgriGreen, bar shift=-3.0pt] table[x=i, y=cash]{data/clean/profiles/fin_003030.dat};
\addplot[fill=SoftRed!45, draw=SoftRed, bar shift=3.0pt] table[x=i, y=ibd]{data/clean/profiles/fin_003030.dat};
\addplot[fill=AgriGreen!25, draw=AgriGreen, pattern=north east lines, bar shift=-3.0pt] table[x=x, y=cash]{data/clean/profiles/finh_003030.dat};
\addplot[fill=SoftRed!20, draw=SoftRed, pattern=north east lines, bar shift=3.0pt] table[x=x, y=ibd]{data/clean/profiles/finh_003030.dat};
\addplot[OilBlue, mark=*, mark size=1.1pt, line width=0.8pt] table[x=i, y=ocf]{data/clean/profiles/fin_003030.dat};
\addplot[only marks, OilBlue, mark=*, mark size=1.8pt] table[x=x, y=ocf]{data/clean/profiles/finh_003030.dat};
\legend{货币资金, 有息负债, 经营活动现金流净额}
\end{axis}
\end{tikzpicture}

\captionof{figure}[祖名股份（003030）近年主要财务指标]{祖名股份（003030）近年主要财务指标（左上：营业收入与归母净利润；右上：ROE 与资产负债率；下：货币资金、有息负债与经营活动现金流净额。
2021---2025 年为年报口径，2026H1 为 2026 年半年报/半年末，与其前一年报之间留出间隔、以斜纹或空心点区分；半年报数据未年化）}
\end{center}

\pfl{财务分析} 2026 年 6 月末资产负债率 57.9\%，较 2025 年末 上升 3.3 个百分点（样本中位数 52.2\%，所处三级行业内排第 \zhnumber{9} 位）。 2026 年 6 月末货币资金 1.80 亿元、短期债务（短期借款与一年内到期非流动负债）8.74 亿元、长期债务 1.28 亿元，有息负债合计 10.02 亿元，现金短债比 0.21 倍——缺口明显。 净有息负债 8.22 亿元，相当于其 2026 年上半年营业收入的 81\%。 流动比率 0.54、速动比率 0.44，短期偿付压力较大。 2026 年上半年经营活动现金流净额 0.58 亿元，净利润为负而经营现金流为正，差额主要来自折旧摊销与营运资本变动，说明亏损尚未完全转化为现金流出。

\pfl{重大战略转型} 公告口径上，近三年（2023 年 9 月---2026 年 9 月）公司披露重大重组与并购类公告 3 条、对外投资与转型类公告 2 条、回购与股权激励类公告 27 条，合计公告 348 条。 第一大产品仍是“生鲜豆制品”，收入占比 67.3\%，与 2021 年年报相比没有方向性变化。 综合判断，公司报告期内未发生实质性的战略转型，资本开支与资本运作均以维持现有业务为主。

\pfl{历史融投资情况} 从现金流量表看，2025 年公司取得借款收到的现金 7.04 亿元、偿还债务支付的现金 6.67 亿元、吸收权益性投资 0.00 亿元，筹资活动现金流净额 -0.59 亿元（2023 年为取得借款 5.14 亿元、偿还债务 4.00 亿元）。 2026 年上半年筹资活动现金流净额为 1.20 亿元（取得借款 5.02 亿元、偿还债务 3.94 亿元，未年化），仍处于净融入状态。 2025 年分配股利、利润或偿付利息支付的现金合计 0.48 亿元。 股本方面，近三年披露股本变动 4 次（限售股份上市 1 次）；总股本由 2021 年末的 1.25 亿股变为 1.25 亿股（+0.00\%）。 分红方面，近三年现金分红 4 次、累计每股派现 0.60 元（最近一次 2025-12-31）——在利润承压的阶段仍维持了分红。

\pfl{未来潜在融资需求分析} 2026 年 6 月末货币资金 1.80 亿元，而短期债务 8.74 亿元（现金短债比 0.21 倍），2026 年上半年经营现金流净额 0.58 亿元。按“短期债务减去账面货币资金”的滚动偿付口径，静态缺口约 6.94 亿元；若再计入上半年亏损的年化额（0.49 亿元），维持一年正常经营所需的外部资金约 6.94---7.43 亿元。 公司 2026 年上半年归母净利润为负（-0.24 亿元），但 6 月末资产负债率 57.9\% 尚未进入第一档（亏损且负债率 $>$65\%）的区间，当前融资需求以补充流动资金、支撑低谷期经营性支出为主。


\subsubsection{*ST广糖（000911）}
\label{co:000911}

\pfl{公司基本情况} 主营业务为“机制糖的生产与销售”，按申万行业分类（2021 版）归属三级行业“其他农产品加工”。 公司于 1999-05-27 在深交所主板首发上市，注册资本 4.00 亿元，总股本 4.00 亿股。 截至 2026-09-24收盘价 5.11 元，流通市值 20.5 亿元，近一年-30.6\%，流通市值在三级行业“其他农产品加工”的 11 家公司中排第 \zhnumber{10} 位，近一年涨跌幅低于全样本中位数（-10.1\%）20.5 个百分点。 按 2026 年半年报披露的行业构成，收入占比前三位为工业 88.6\%；物流贸易及其他 11.4\%。

\pfl{历史股价} 自 2021 年 1 月至 2026 年 9 月，股价由 7.70 元变为 5.11 元，累计 -33.6\%，区间最高 13.11 元（2021 年 8 月）、最低 3.83 元（2026 年 6 月）；当前价相当于区间最高价的 39.0\%，为区间最低价的 1.33 倍。 月度收益的年化波动率 40.6\%，高于全样本中位数 37.4\%，在三级行业“其他农产品加工”的 11 家公司中波动率排第 \zhnumber{4} 位。

\begin{center}
\begin{tikzpicture}
\begin{axis}[
  width=12.6cm, height=3.9cm, scale only axis,
  xmin=2020.922, xmax=2026.888, ymin=3.56, ymax=14.03,
  xtick={2021,2022,2023,2024,2025,2026},
  yearticks,
  yticklabel style={font=\scriptsize},
  ylabel={元/股}, ylabel style={font=\scriptsize},
  grid=major, grid style={LightGray, line width=0.3pt},
  axis line style={InkGray, line width=0.5pt},
  every axis plot/.append style={line width=0.9pt},
]
\addplot[AgriGreen] table[x=x,y=y]{data/clean/profiles/px_000911.dat};
\addplot[SoftRed, dashed, line width=0.7pt] table[x=x,y=y]{data/clean/profiles/pxzz_000911.dat};
\addplot[only marks, mark=*, mark size=1.1pt, SoftRed] table[x=x,y=y]{data/clean/profiles/pxzz_000911.dat};
\end{axis}
\end{tikzpicture}

\captionof{figure}[*ST广糖（000911）股价与周期骨架]{*ST广糖（000911）月度收盘价（元/股）与 ZigZag 周期骨架（阈值 25\%，圆点为周期转折点）}
\end{center}

\pfl{过去周期分析} 以 25\% 的 ZigZag 阈值划分，样本区间内股价形成 5 个主升段与 4 个主跌段。 最大主跌段为 2024 年 11 月 至 2026 年 6 月，19 个月-55.4\%；最大主升段出现在 2021 年 1 月 至 2021 年 8 月，7 个月+70.3\%。 最近一次转折出现在 2026 年 9 月（上涨段自 2026 年 6 月起，3 个月+33.4\%），当前处于该段之后的震荡之中。 公司股价与申万农林牧渔指数几乎同步见底，个股并未走出独立行情。 该子行业横跨糖、番茄制品、添加剂等多个品种，各自的供给周期与政策（进口配额、关税）是主要变量，公司 2026 年 6 月末资产负债率 100.2\%，高于全样本中位数 52.2\%，其股价因此更多反映资金面与信用风险，而不是价格弹性本身。

\pfl{盈利情况} 2026 年上半年营业收入 15.82 亿元（同比 +17.4\%），归母净利润 -0.16 亿元，由上年同期的盈利转为亏损。 以年报为趋势对照，2023---2025 年营业收入由 33.42 亿元变为 24.67 亿元（两年年均复合增速 -14.1\%），归母净利润由 0.07 亿元变为 -1.20 亿元。 按半年报口径，2026 年上半年的 ROE 为 0.0\%（未年化）、销售净利率 -1.1\%、毛利率 9.1\%，ROE 在“其他农产品加工”11 家公司中排第 \zhnumber{7} 位。 同期全样本半年 ROE 中位数 0.75\%，公司与之相差 0.75 个百分点（弱于中位数公司）。

\begin{center}\small\setlength{\tabcolsep}{4pt}
\captionof{table}[*ST广糖（000911）关键财务指标]{*ST广糖（000911）关键财务指标（合并报表口径；两个半年报期均未年化；现金短债比 $=$ 货币资金 $\div$ 短期债务）}
\begin{tabular}{@{}lrrrrr@{}}
\toprule
指标 & 2023 年 & 2024 年 & 2025 年 & 2025 年半年报 & 2026 年半年报 \\
\midrule
营业收入（亿元） & 33.42 & 32.75 & 24.67 & 13.48 & 15.82 \\
归母净利润（亿元） & 0.07 & 0.08 & -1.20 & 0.09 & -0.16 \\
毛利率（\%） & 12.8 & 11.9 & 9.7 & 12.1 & 9.1 \\
ROE（\%） & 41.6 & 16.4 & 0.0 & 15.6 & 0.0 \\
资产负债率（\%） & 97.0 & 96.4 & 99.8 & 95.7 & 100.2 \\
经营活动现金流净额（亿元） & 5.51 & 4.54 & 3.61 & -2.04 & -11.74 \\
货币资金（亿元） & 14.27 & 7.61 & 7.93 & 2.72 & 2.29 \\
有息负债合计（亿元） & 37.39 & 26.92 & 25.17 & 25.19 & 31.13 \\
现金短债比（倍） & 0.41 & 0.30 & 0.32 & 0.11 & 0.08 \\
\bottomrule
\end{tabular}
\end{center}

\begin{center}
\begin{tikzpicture}
\begin{groupplot}[
  group style={group size=2 by 1, horizontal sep=0.95cm},
  width=7.2cm, height=4.6cm,
  tick label style={font=\tiny},
  title style={font=\scriptsize\heiti},
  yticklabel style={font=\tiny},
  grid=major, grid style={LightGray, line width=0.3pt},
  axis line style={InkGray, line width=0.5pt},
  enlarge x limits={abs=0.8},
  legend style={font=\scriptsize, at={(0.5,-0.22)}, anchor=north,
                legend columns=2, draw=none, fill=none, column sep=6pt},
]
\nextgroupplot[
  ybar, bar width=4pt,
  ymin=-8.37, ymax=37.97,
  xtick={0,1,2,3,4,5.7},
  xticklabels={2021,2022,2023,2024,2025,2026H1},
  ylabel={亿元},
]
\addplot[fill=AgriGreen!75, draw=AgriGreen, bar shift=-2.6pt] table[x=i, y=rev]{data/clean/profiles/fin_000911.dat};
\addplot[fill=SoftRed!60, draw=SoftRed, bar shift=2.6pt] table[x=i, y=ni]{data/clean/profiles/fin_000911.dat};
\addplot[fill=AgriGreen!30, draw=AgriGreen, pattern=north east lines, bar shift=-2.6pt] table[x=x, y=rev]{data/clean/profiles/finh_000911.dat};
\addplot[fill=SoftRed!20, draw=SoftRed, pattern=north east lines, bar shift=2.6pt] table[x=x, y=ni]{data/clean/profiles/finh_000911.dat};
\legend{营业收入（年/半年）, 归母净利润（年/半年）}
\nextgroupplot[
  ymin=-191.7, ymax=127.2,
  xtick={0,1,2,3,4,5.7},
  xticklabels={2021,2022,2023,2024,2025,2026H1},
  ylabel={\%},
]
\addplot[AgriGold, mark=*, mark size=1.3pt] table[x=i, y=roe]{data/clean/profiles/fin_000911.dat};
\addplot[OilBlue, mark=square*, mark size=1.3pt] table[x=i, y=debt]{data/clean/profiles/fin_000911.dat};
\addplot[only marks, AgriGold, mark=*, mark size=2.0pt] table[x=x, y=roe]{data/clean/profiles/finh_000911.dat};
\addplot[only marks, OilBlue, mark=square*, mark size=2.0pt] table[x=x, y=debt]{data/clean/profiles/finh_000911.dat};
\legend{ROE, 资产负债率}
\end{groupplot}
\end{tikzpicture}
\begin{tikzpicture}
\begin{axis}[
  width=15.2cm, height=3.1cm, scale only axis,
  ymin=-18.22, ymax=60.89,
  xtick={0,1,2,3,4,5.7},
  xticklabels={2021,2022,2023,2024,2025,2026H1},
  tick label style={font=\tiny},
  yticklabel style={font=\tiny},
  ylabel={亿元}, ylabel style={font=\scriptsize},
  ybar, bar width=4pt,
  grid=major, grid style={LightGray, line width=0.3pt},
  axis line style={InkGray, line width=0.5pt},
  enlarge x limits={abs=0.8},
  legend style={font=\scriptsize, at={(0.5,-0.30)}, anchor=north,
                legend columns=3, draw=none, fill=none, column sep=10pt},
]
\addplot[fill=AgriGreen!55, draw=AgriGreen, bar shift=-3.0pt] table[x=i, y=cash]{data/clean/profiles/fin_000911.dat};
\addplot[fill=SoftRed!45, draw=SoftRed, bar shift=3.0pt] table[x=i, y=ibd]{data/clean/profiles/fin_000911.dat};
\addplot[fill=AgriGreen!25, draw=AgriGreen, pattern=north east lines, bar shift=-3.0pt] table[x=x, y=cash]{data/clean/profiles/finh_000911.dat};
\addplot[fill=SoftRed!20, draw=SoftRed, pattern=north east lines, bar shift=3.0pt] table[x=x, y=ibd]{data/clean/profiles/finh_000911.dat};
\addplot[OilBlue, mark=*, mark size=1.1pt, line width=0.8pt] table[x=i, y=ocf]{data/clean/profiles/fin_000911.dat};
\addplot[only marks, OilBlue, mark=*, mark size=1.8pt] table[x=x, y=ocf]{data/clean/profiles/finh_000911.dat};
\legend{货币资金, 有息负债, 经营活动现金流净额}
\end{axis}
\end{tikzpicture}

\captionof{figure}[*ST广糖（000911）近年主要财务指标]{*ST广糖（000911）近年主要财务指标（左上：营业收入与归母净利润；右上：ROE 与资产负债率；下：货币资金、有息负债与经营活动现金流净额。
2021---2025 年为年报口径，2026H1 为 2026 年半年报/半年末，与其前一年报之间留出间隔、以斜纹或空心点区分；半年报数据未年化）}
\end{center}

\pfl{财务分析} 2026 年 6 月末资产负债率 100.2\%，较 2025 年末 上升 0.4 个百分点（样本中位数 52.2\%，所处三级行业内排第 \zhnumber{11} 位）。 2026 年 6 月末货币资金 2.29 亿元、短期债务（短期借款与一年内到期非流动负债）29.24 亿元、长期债务 1.90 亿元，有息负债合计 31.13 亿元，现金短债比 0.08 倍——缺口明显。 净有息负债 28.85 亿元，相当于其 2026 年上半年营业收入的 182\%。 流动比率 0.60、速动比率 0.27，短期偿付压力较大。 2026 年上半年经营活动现金流净额 -11.74 亿元，经营现金流与利润同时为负，日常运营对债务与股东投入的依赖度较高。

\pfl{重大战略转型} 公告口径上，近三年（2023 年 9 月---2026 年 9 月）公司披露重大重组与并购类公告 1 条、对外投资与转型类公告 1 条、回购与股权激励类公告 1 条，合计公告 496 条。 第一大行业仍是“工业”，收入占比 88.6\%，与 2021 年年报相比没有方向性变化。 公司证券简称带风险警示标识，报告期内的资本运作以维持上市地位为主线，属于第 \ref{sec:financing-strategy} 节界定的“跨界与保壳”路径。

\pfl{历史融投资情况} 从现金流量表看，2025 年公司取得借款收到的现金 29.04 亿元、偿还债务支付的现金 30.13 亿元、吸收权益性投资 — 亿元，筹资活动现金流净额 -2.06 亿元（2023 年为取得借款 45.31 亿元、偿还债务 46.64 亿元）。 2026 年上半年筹资活动现金流净额为 6.23 亿元（取得借款 14.75 亿元、偿还债务 9.75 亿元，未年化），仍处于净融入状态。 2025 年分配股利、利润或偿付利息支付的现金合计 0.86 亿元。 股本方面，近三年披露股本变动 3 次；总股本由 2021 年末的 4.00 亿股变为 4.00 亿股（+0.00\%）。 分红方面，近三年未实施现金分红，在全部样本中属于分红能力偏弱的一端。 近三年“再融资”类公告 78 条，涉及定向增发、募集资金使用。

\pfl{未来潜在融资需求分析} 2026 年 6 月末货币资金 2.29 亿元，而短期债务 29.24 亿元（现金短债比 0.08 倍），2026 年上半年经营现金流净额 -11.74 亿元。按“短期债务减去账面货币资金”的滚动偿付口径，静态缺口约 26.95 亿元；若再计入上半年亏损的年化额（0.32 亿元），维持一年正常经营所需的外部资金约 26.95---27.27 亿元。 按第 \ref{sec:financing-need} 节的三档划分，公司属于第一档“补血型”：2026 年上半年归母净利润 -0.16 亿元、6 月末资产负债率 100.2\%。负债率已超过 80\%，净资产被亏损侵蚀的程度较深，属于全部样本中资产负债表最紧张的一组。 可行的融资路径按现实性排序为：债务重组或展期（与银行和债券持有人协商）、定向增发（需股价配合，且控股股东需放弃优先认购或同步增资）、出售非核心资产与参股股权。


\subsubsection{索宝蛋白（603231）}
\label{co:603231}

\pfl{公司基本情况} 主营业务为“大豆蛋白系列产品的研发、生产和销售”，按申万行业分类（2021 版）归属三级行业“粮油加工”。 公司于 2023-12-15 在上交所首发上市，注册资本 1.91 亿元，总股本 1.91 亿股。 截至 2026-09-24收盘价 17.16 元，流通市值 18.2 亿元，近一年-9.4\%，流通市值在三级行业“粮油加工”的 7 家公司中排第 \zhnumber{7} 位，近一年涨跌幅高于全样本中位数（-10.1\%）0.7 个百分点。

\pfl{历史股价} 本报告样本区间（2023 年 12 月 至 2026 年 9 月）内，股价由 23.52 元变为 17.16 元，累计 -27.0\%，区间最高 23.52 元（2023 年 12 月）、最低 14.76 元（2026 年 7 月）；当前价相当于区间最高价的 73.0\%，为区间最低价的 1.16 倍。 月度收益的年化波动率 28.6\%，低于全样本中位数 37.4\%，在三级行业“粮油加工”的 7 家公司中波动率排第 \zhnumber{5} 位。

\begin{center}
\begin{tikzpicture}
\begin{axis}[
  width=12.6cm, height=3.9cm, scale only axis,
  xmin=2023.838, xmax=2026.888, ymin=13.73, ymax=25.17,
  xtick={2023,2024,2025,2026},
  yearticks,
  yticklabel style={font=\scriptsize},
  ylabel={元/股}, ylabel style={font=\scriptsize},
  grid=major, grid style={LightGray, line width=0.3pt},
  axis line style={InkGray, line width=0.5pt},
  every axis plot/.append style={line width=0.9pt},
]
\addplot[AgriGreen] table[x=x,y=y]{data/clean/profiles/px_603231.dat};
\addplot[SoftRed, dashed, line width=0.7pt] table[x=x,y=y]{data/clean/profiles/pxzz_603231.dat};
\addplot[only marks, mark=*, mark size=1.1pt, SoftRed] table[x=x,y=y]{data/clean/profiles/pxzz_603231.dat};
\end{axis}
\end{tikzpicture}

\captionof{figure}[索宝蛋白（603231）股价与周期骨架]{索宝蛋白（603231）月度收盘价（元/股）与 ZigZag 周期骨架（阈值 25\%，圆点为周期转折点）}
\end{center}

\pfl{过去周期分析} 以 25\% 的 ZigZag 阈值划分，样本区间内股价形成 1 个主升段与 2 个主跌段。 最大主跌段为 2023 年 12 月 至 2024 年 8 月，8 个月-35.5\%；最大主升段出现在 2024 年 8 月 至 2025 年 10 月，14 个月+33.8\%。 最近一次转折出现在 2026 年 7 月（下跌段自 2025 年 10 月起，9 个月-27.3\%），当前处于该段的延续之中。 与申万农林牧渔指数的最近一次底部相比，该公司见底晚约 2 个月，在本轮下行中属于调整更慢的一类。 粮油加工是“成本加成”型生意（第 \ref{sec:grain} 章），利润来自原料与成品的价差及规模效率，收入规模大而净利率薄，公司 2026 年上半年实现盈利、半年 ROE 3.4\%（未年化），好于全样本中位数（0.75\%），下行周期对其报表的冲击小于同业。

\pfl{盈利情况} 2026 年上半年营业收入 9.67 亿元（同比 +24.9\%），归母净利润 0.68 亿元，同比 -30.1\%。 以年报为趋势对照，2023---2025 年营业收入由 17.48 亿元变为 17.48 亿元（两年年均复合增速 -0.0\%），归母净利润由 1.46 亿元变为 1.90 亿元。 按半年报口径，2026 年上半年的 ROE 为 3.4\%（未年化）、销售净利率 6.9\%、毛利率 12.9\%，ROE 在“粮油加工”7 家公司中排第 \zhnumber{2} 位。 同期全样本半年 ROE 中位数 0.75\%，公司与之相差 2.62 个百分点（略好于中位数公司）。

\begin{center}\small\setlength{\tabcolsep}{4pt}
\captionof{table}[索宝蛋白（603231）关键财务指标]{索宝蛋白（603231）关键财务指标（合并报表口径；两个半年报期均未年化；现金短债比 $=$ 货币资金 $\div$ 短期债务）}
\begin{tabular}{@{}lrrrrr@{}}
\toprule
指标 & 2023 年 & 2024 年 & 2025 年 & 2025 年半年报 & 2026 年半年报 \\
\midrule
营业收入（亿元） & 17.48 & 15.55 & 17.48 & 7.74 & 9.67 \\
归母净利润（亿元） & 1.46 & 1.21 & 1.90 & 0.97 & 0.68 \\
毛利率（\%） & 13.6 & 12.0 & 16.7 & 18.7 & 12.9 \\
ROE（\%） & 16.3 & 6.4 & 9.8 & 5.1 & 3.4 \\
资产负债率（\%） & 11.7 & 10.1 & 11.6 & 8.7 & 13.6 \\
经营活动现金流净额（亿元） & 1.44 & 1.48 & 3.11 & 0.29 & 0.08 \\
货币资金（亿元） & 10.73 & 3.20 & 2.39 & 2.61 & 2.84 \\
有息负债合计（亿元） & 0.85 & 0.24 & 0.24 & 0.24 & 0.67 \\
现金短债比（倍） & 12.69 & 13.07 & 9.84 & 10.94 & 4.39 \\
\bottomrule
\end{tabular}
\end{center}

\begin{center}
\begin{tikzpicture}
\begin{groupplot}[
  group style={group size=2 by 1, horizontal sep=0.95cm},
  width=7.2cm, height=4.6cm,
  tick label style={font=\tiny},
  title style={font=\scriptsize\heiti},
  yticklabel style={font=\tiny},
  grid=major, grid style={LightGray, line width=0.3pt},
  axis line style={InkGray, line width=0.5pt},
  enlarge x limits={abs=0.8},
  legend style={font=\scriptsize, at={(0.5,-0.22)}, anchor=north,
                legend columns=2, draw=none, fill=none, column sep=6pt},
]
\nextgroupplot[
  ybar, bar width=4pt,
  ymin=-1.85, ymax=20.68,
  xtick={0,1,2,3,4,5.7},
  xticklabels={2021,2022,2023,2024,2025,2026H1},
  ylabel={亿元},
]
\addplot[fill=AgriGreen!75, draw=AgriGreen, bar shift=-2.6pt] table[x=i, y=rev]{data/clean/profiles/fin_603231.dat};
\addplot[fill=SoftRed!60, draw=SoftRed, bar shift=2.6pt] table[x=i, y=ni]{data/clean/profiles/fin_603231.dat};
\addplot[fill=AgriGreen!30, draw=AgriGreen, pattern=north east lines, bar shift=-2.6pt] table[x=x, y=rev]{data/clean/profiles/finh_603231.dat};
\addplot[fill=SoftRed!20, draw=SoftRed, pattern=north east lines, bar shift=2.6pt] table[x=x, y=ni]{data/clean/profiles/finh_603231.dat};
\legend{营业收入（年/半年）, 归母净利润（年/半年）}
\nextgroupplot[
  ymin=-2.4, ymax=32.9,
  xtick={0,1,2,3,4,5.7},
  xticklabels={2021,2022,2023,2024,2025,2026H1},
  ylabel={\%},
]
\addplot[AgriGold, mark=*, mark size=1.3pt] table[x=i, y=roe]{data/clean/profiles/fin_603231.dat};
\addplot[OilBlue, mark=square*, mark size=1.3pt] table[x=i, y=debt]{data/clean/profiles/fin_603231.dat};
\addplot[only marks, AgriGold, mark=*, mark size=2.0pt] table[x=x, y=roe]{data/clean/profiles/finh_603231.dat};
\addplot[only marks, OilBlue, mark=square*, mark size=2.0pt] table[x=x, y=debt]{data/clean/profiles/finh_603231.dat};
\legend{ROE, 资产负债率}
\end{groupplot}
\end{tikzpicture}
\begin{tikzpicture}
\begin{axis}[
  width=15.2cm, height=3.1cm, scale only axis,
  ymin=-1.07, ymax=12.02,
  xtick={0,1,2,3,4,5.7},
  xticklabels={2021,2022,2023,2024,2025,2026H1},
  tick label style={font=\tiny},
  yticklabel style={font=\tiny},
  ylabel={亿元}, ylabel style={font=\scriptsize},
  ybar, bar width=4pt,
  grid=major, grid style={LightGray, line width=0.3pt},
  axis line style={InkGray, line width=0.5pt},
  enlarge x limits={abs=0.8},
  legend style={font=\scriptsize, at={(0.5,-0.30)}, anchor=north,
                legend columns=3, draw=none, fill=none, column sep=10pt},
]
\addplot[fill=AgriGreen!55, draw=AgriGreen, bar shift=-3.0pt] table[x=i, y=cash]{data/clean/profiles/fin_603231.dat};
\addplot[fill=SoftRed!45, draw=SoftRed, bar shift=3.0pt] table[x=i, y=ibd]{data/clean/profiles/fin_603231.dat};
\addplot[fill=AgriGreen!25, draw=AgriGreen, pattern=north east lines, bar shift=-3.0pt] table[x=x, y=cash]{data/clean/profiles/finh_603231.dat};
\addplot[fill=SoftRed!20, draw=SoftRed, pattern=north east lines, bar shift=3.0pt] table[x=x, y=ibd]{data/clean/profiles/finh_603231.dat};
\addplot[OilBlue, mark=*, mark size=1.1pt, line width=0.8pt] table[x=i, y=ocf]{data/clean/profiles/fin_603231.dat};
\addplot[only marks, OilBlue, mark=*, mark size=1.8pt] table[x=x, y=ocf]{data/clean/profiles/finh_603231.dat};
\legend{货币资金, 有息负债, 经营活动现金流净额}
\end{axis}
\end{tikzpicture}

\captionof{figure}[索宝蛋白（603231）近年主要财务指标]{索宝蛋白（603231）近年主要财务指标（左上：营业收入与归母净利润；右上：ROE 与资产负债率；下：货币资金、有息负债与经营活动现金流净额。
2021---2025 年为年报口径，2026H1 为 2026 年半年报/半年末，与其前一年报之间留出间隔、以斜纹或空心点区分；半年报数据未年化）}
\end{center}

\pfl{财务分析} 2026 年 6 月末资产负债率 13.6\%，较 2025 年末 上升 2.0 个百分点（样本中位数 52.2\%，所处三级行业内排第 \zhnumber{1} 位）。 2026 年 6 月末货币资金 2.84 亿元、短期债务（短期借款与一年内到期非流动负债）0.65 亿元、长期债务 0.02 亿元，有息负债合计 0.67 亿元，现金短债比 4.39 倍——覆盖充分。 净有息负债 -2.17 亿元，相当于其 2026 年上半年营业收入的 -22\%。 流动比率 5.42、速动比率 4.16，短期指标尚可。 2026 年上半年经营活动现金流净额 0.08 亿元，经营现金流与利润同向为正，内生资金可以支撑运营。

\pfl{重大战略转型} 公告口径上，近三年（2023 年 9 月---2026 年 9 月）公司披露重大重组与并购类公告 1 条、对外投资与转型类公告 2 条、回购与股权激励类公告 0 条，合计公告 283 条。 综合判断，公司报告期内未发生实质性的战略转型，资本开支与资本运作均以维持现有业务为主。

\pfl{历史融投资情况} 从现金流量表看，2025 年公司取得借款收到的现金 0.48 亿元、偿还债务支付的现金 0.48 亿元、吸收权益性投资 — 亿元，筹资活动现金流净额 -0.59 亿元（2023 年为取得借款 0.94 亿元、偿还债务 0.95 亿元）。 2026 年上半年筹资活动现金流净额为 -0.30 亿元（取得借款 0.84 亿元、偿还债务 0.44 亿元，未年化），已转为净流出，处于净还债状态。 2025 年分配股利、利润或偿付利息支付的现金合计 0.57 亿元。 股本方面，近三年披露股本变动 4 次（A股上市 1 次、限售股份上市 1 次）；总股本由 2021 年末的 1.91 亿股变为 1.91 亿股（+0.00\%）。 分红方面，近三年现金分红 4 次、累计每股派现 1.17 元（最近一次 2026-06-25）——在利润承压的阶段仍维持了分红。

\pfl{未来潜在融资需求分析} 2026 年 6 月末货币资金 2.84 亿元，而短期债务 0.65 亿元（现金短债比 4.39 倍），2026 年上半年经营现金流净额 0.08 亿元。账面货币资金可以覆盖短期债务，缺口不体现在静态偿付层面。 公司 2026 年 6 月末资产负债率 13.6\%、2026 年上半年 ROE 3.4\%（未年化），资产负债表稳健、盈利水平平淡，暂不存在刚性融资需求；潜在的融资需求以技改扩产、并购或被并购为主。


\subsubsection{佳沃食品（300268）}
\label{co:300268}

\pfl{公司基本情况} 主营业务为“海产品进出口贸易、加工及销售”，按申万行业分类（2021 版）归属三级行业“其他农产品加工”。 公司于 2011-09-27 在深交所创业板首发上市，注册资本 1.74 亿元，总股本 1.74 亿股。 截至 2026-09-24收盘价 12.00 元，流通市值 16.0 亿元，近一年+17.9\%，流通市值在三级行业“其他农产品加工”的 11 家公司中排第 \zhnumber{11} 位，近一年涨跌幅高于全样本中位数（-10.1\%）28.0 个百分点。 按 2025 年年报披露的行业构成，收入占比前三位为动物蛋白 99.9\%；其他 0.1\%。

\pfl{历史股价} 自 2021 年 1 月至 2026 年 9 月，股价由 8.98 元变为 12.00 元，累计 +33.6\%，区间最高 30.33 元（2022 年 5 月）、最低 6.27 元（2024 年 6 月）；当前价相当于区间最高价的 39.6\%，为区间最低价的 1.91 倍。 月度收益的年化波动率 59.4\%，高于全样本中位数 37.4\%，在三级行业“其他农产品加工”的 11 家公司中波动率排第 \zhnumber{1} 位。

\begin{center}
\begin{tikzpicture}
\begin{axis}[
  width=12.6cm, height=3.9cm, scale only axis,
  xmin=2020.922, xmax=2026.888, ymin=5.83, ymax=32.45,
  xtick={2021,2022,2023,2024,2025,2026},
  yearticks,
  yticklabel style={font=\scriptsize},
  ylabel={元/股}, ylabel style={font=\scriptsize},
  grid=major, grid style={LightGray, line width=0.3pt},
  axis line style={InkGray, line width=0.5pt},
  every axis plot/.append style={line width=0.9pt},
]
\addplot[AgriGreen] table[x=x,y=y]{data/clean/profiles/px_300268.dat};
\addplot[SoftRed, dashed, line width=0.7pt] table[x=x,y=y]{data/clean/profiles/pxzz_300268.dat};
\addplot[only marks, mark=*, mark size=1.1pt, SoftRed] table[x=x,y=y]{data/clean/profiles/pxzz_300268.dat};
\end{axis}
\end{tikzpicture}

\captionof{figure}[佳沃食品（300268）股价与周期骨架]{佳沃食品（300268）月度收盘价（元/股）与 ZigZag 周期骨架（阈值 25\%，圆点为周期转折点）}
\end{center}

\pfl{过去周期分析} 以 25\% 的 ZigZag 阈值划分，样本区间内股价形成 4 个主升段与 4 个主跌段。 最大主跌段为 2022 年 12 月 至 2024 年 6 月，18 个月-73.3\%；最大主升段出现在 2021 年 1 月 至 2021 年 5 月，4 个月+157.5\%。 最近一次转折出现在 2026 年 6 月（下跌段自 2026 年 2 月起，4 个月-29.2\%），当前处于该段的延续之中。 公司股价较申万农林牧渔指数提前约 23 个月见底，是板块中较早定价周期反转的公司之一。 该子行业横跨糖、番茄制品、添加剂等多个品种，各自的供给周期与政策（进口配额、关税）是主要变量，公司 2026 年 6 月末资产负债率 35.5\%、半年 ROE -1.4\%（未年化），在本轮下行中属于随行业同步调整的一类。

\pfl{盈利情况} 2026 年上半年营业收入 4.21 亿元（同比 -66.1\%），归母净利润 -0.04 亿元，亏损同比收窄 4.15 亿元。 以年报为趋势对照，2023---2025 年营业收入由 45.05 亿元变为 16.57 亿元（两年年均复合增速 -39.4\%），归母净利润由 -11.52 亿元变为 -4.27 亿元。 按半年报口径，2026 年上半年的 ROE 为 -1.4\%（未年化）、销售净利率 -0.8\%、毛利率 8.4\%，ROE 在“其他农产品加工”11 家公司中排第 \zhnumber{8} 位。 行业同口径半年 ROE 的中位数为 0.75\%，该公司低于中位数 2.14 个百分点。

\begin{center}\small\setlength{\tabcolsep}{4pt}
\captionof{table}[佳沃食品（300268）关键财务指标]{佳沃食品（300268）关键财务指标（合并报表口径；两个半年报期均未年化；现金短债比 $=$ 货币资金 $\div$ 短期债务）}
\begin{tabular}{@{}lrrrrr@{}}
\toprule
指标 & 2023 年 & 2024 年 & 2025 年 & 2025 年半年报 & 2026 年半年报 \\
\midrule
营业收入（亿元） & 45.05 & 34.18 & 16.57 & 12.45 & 4.21 \\
归母净利润（亿元） & -11.52 & -9.24 & -4.27 & -4.19 & -0.04 \\
毛利率（\%） & -5.4 & -10.5 & -2.7 & -6.1 & 8.4 \\
ROE（\%） & --- & --- & --- & --- & -1.4 \\
资产负债率（\%） & 95.6 & 104.9 & 10.9 & 23.8 & 35.5 \\
经营活动现金流净额（亿元） & -0.54 & -1.44 & -5.15 & -6.17 & -2.18 \\
货币资金（亿元） & 1.44 & 2.16 & 1.00 & 0.48 & 0.48 \\
有息负债合计（亿元） & 32.03 & 13.68 & 0.02 & 0.42 & 1.72 \\
现金短债比（倍） & 0.05 & 1.39 & 71.83 & 1.16 & 0.28 \\
\bottomrule
\end{tabular}
\end{center}

\begin{center}
\begin{tikzpicture}
\begin{groupplot}[
  group style={group size=2 by 1, horizontal sep=0.95cm},
  width=7.2cm, height=4.6cm,
  tick label style={font=\tiny},
  title style={font=\scriptsize\heiti},
  yticklabel style={font=\tiny},
  grid=major, grid style={LightGray, line width=0.3pt},
  axis line style={InkGray, line width=0.5pt},
  enlarge x limits={abs=0.8},
  legend style={font=\scriptsize, at={(0.5,-0.22)}, anchor=north,
                legend columns=2, draw=none, fill=none, column sep=6pt},
]
\nextgroupplot[
  ybar, bar width=4pt,
  ymin=-18.19, ymax=63.29,
  xtick={0,1,2,3,4,5.7},
  xticklabels={2021,2022,2023,2024,2025,2026H1},
  ylabel={亿元},
]
\addplot[fill=AgriGreen!75, draw=AgriGreen, bar shift=-2.6pt] table[x=i, y=rev]{data/clean/profiles/fin_300268.dat};
\addplot[fill=SoftRed!60, draw=SoftRed, bar shift=2.6pt] table[x=i, y=ni]{data/clean/profiles/fin_300268.dat};
\addplot[fill=AgriGreen!30, draw=AgriGreen, pattern=north east lines, bar shift=-2.6pt] table[x=x, y=rev]{data/clean/profiles/finh_300268.dat};
\addplot[fill=SoftRed!20, draw=SoftRed, pattern=north east lines, bar shift=2.6pt] table[x=x, y=ni]{data/clean/profiles/finh_300268.dat};
\legend{营业收入（年/半年）, 归母净利润（年/半年）}
\nextgroupplot[
  ymin=-48.4, ymax=119.1,
  xtick={0,1,2,3,4,5.7},
  xticklabels={2021,2022,2023,2024,2025,2026H1},
  ylabel={\%},
]
\addplot[AgriGold, mark=*, mark size=1.3pt] table[x=i, y=roe]{data/clean/profiles/fin_300268.dat};
\addplot[OilBlue, mark=square*, mark size=1.3pt] table[x=i, y=debt]{data/clean/profiles/fin_300268.dat};
\addplot[only marks, AgriGold, mark=*, mark size=2.0pt] table[x=x, y=roe]{data/clean/profiles/finh_300268.dat};
\addplot[only marks, OilBlue, mark=square*, mark size=2.0pt] table[x=x, y=debt]{data/clean/profiles/finh_300268.dat};
\legend{ROE, 资产负债率}
\end{groupplot}
\end{tikzpicture}
\begin{tikzpicture}
\begin{axis}[
  width=15.2cm, height=3.1cm, scale only axis,
  ymin=-10.50, ymax=54.77,
  xtick={0,1,2,3,4,5.7},
  xticklabels={2021,2022,2023,2024,2025,2026H1},
  tick label style={font=\tiny},
  yticklabel style={font=\tiny},
  ylabel={亿元}, ylabel style={font=\scriptsize},
  ybar, bar width=4pt,
  grid=major, grid style={LightGray, line width=0.3pt},
  axis line style={InkGray, line width=0.5pt},
  enlarge x limits={abs=0.8},
  legend style={font=\scriptsize, at={(0.5,-0.30)}, anchor=north,
                legend columns=3, draw=none, fill=none, column sep=10pt},
]
\addplot[fill=AgriGreen!55, draw=AgriGreen, bar shift=-3.0pt] table[x=i, y=cash]{data/clean/profiles/fin_300268.dat};
\addplot[fill=SoftRed!45, draw=SoftRed, bar shift=3.0pt] table[x=i, y=ibd]{data/clean/profiles/fin_300268.dat};
\addplot[fill=AgriGreen!25, draw=AgriGreen, pattern=north east lines, bar shift=-3.0pt] table[x=x, y=cash]{data/clean/profiles/finh_300268.dat};
\addplot[fill=SoftRed!20, draw=SoftRed, pattern=north east lines, bar shift=3.0pt] table[x=x, y=ibd]{data/clean/profiles/finh_300268.dat};
\addplot[OilBlue, mark=*, mark size=1.1pt, line width=0.8pt] table[x=i, y=ocf]{data/clean/profiles/fin_300268.dat};
\addplot[only marks, OilBlue, mark=*, mark size=1.8pt] table[x=x, y=ocf]{data/clean/profiles/finh_300268.dat};
\legend{货币资金, 有息负债, 经营活动现金流净额}
\end{axis}
\end{tikzpicture}

\captionof{figure}[佳沃食品（300268）近年主要财务指标]{佳沃食品（300268）近年主要财务指标（左上：营业收入与归母净利润；右上：ROE 与资产负债率；下：货币资金、有息负债与经营活动现金流净额。
2021---2025 年为年报口径，2026H1 为 2026 年半年报/半年末，与其前一年报之间留出间隔、以斜纹或空心点区分；半年报数据未年化）}
\end{center}

\pfl{财务分析} 2026 年 6 月末资产负债率 35.5\%，较 2025 年末 上升 24.6 个百分点（样本中位数 52.2\%，所处三级行业内排第 \zhnumber{5} 位）。 2026 年 6 月末货币资金 0.48 亿元、短期债务（短期借款与一年内到期非流动负债）1.70 亿元、长期债务 0.01 亿元，有息负债合计 1.72 亿元，现金短债比 0.28 倍——缺口明显。 净有息负债 1.23 亿元，相当于其 2026 年上半年营业收入的 29\%。 流动比率 2.39、速动比率 0.72，短期指标尚可。 2026 年上半年经营活动现金流净额 -2.18 亿元，经营现金流与利润同时为负，现金消耗较快。

\pfl{重大战略转型} 公告口径上，近三年（2023 年 9 月---2026 年 9 月）公司披露重大重组与并购类公告 20 条、对外投资与转型类公告 1 条、回购与股权激励类公告 2 条，合计公告 495 条。 第一大行业仍是“动物蛋白”，收入占比 99.9\%，与 2021 年年报相比没有方向性变化。 综合判断，公司在报告期内有明确的外延扩张或业务结构调整动作，转型方向与融资安排之间存在直接关联。

\pfl{历史融投资情况} 从现金流量表看，2025 年公司取得借款收到的现金 0.92 亿元、偿还债务支付的现金 1.31 亿元、吸收权益性投资 — 亿元，筹资活动现金流净额 5.43 亿元（2023 年为取得借款 3.53 亿元、偿还债务 15.33 亿元）。 2026 年上半年筹资活动现金流净额为 1.66 亿元（取得借款 1.94 亿元、偿还债务 0.26 亿元，未年化），仍处于净融入状态。 2025 年分配股利、利润或偿付利息支付的现金合计 0.23 亿元。 股本方面，近三年披露股本变动 3 次；总股本由 2021 年末的 1.74 亿股变为 1.74 亿股（+0.00\%）。 分红方面，近三年未实施现金分红，在全部样本中属于分红能力偏弱的一端。

\pfl{未来潜在融资需求分析} 2026 年 6 月末货币资金 0.48 亿元，而短期债务 1.70 亿元（现金短债比 0.28 倍），2026 年上半年经营现金流净额 -2.18 亿元。按“短期债务减去账面货币资金”的滚动偿付口径，静态缺口约 1.22 亿元；若再计入上半年亏损的年化额（0.08 亿元），维持一年正常经营所需的外部资金约 1.22---1.30 亿元。 按第 \ref{sec:financing-need} 节的三档划分，公司属于第二档“保壳型”：近三年重大重组与并购类公告 20 条，2026 年上半年归母净利润 -0.04 亿元，融资需求更多体现为“资产置换 + 维持上市地位”，而不是扩产。



